\chapter{Modular curves over \texorpdfstring{$\Q$}{Q}, Jacobians and the modular degree}\label{ch:moddeg}
\chaptermark{Modular curves and the modular degree}

The paper is about one integer, the modular degree $\degphi$ of the $X_0(N)$-optimal curve $E$ of an isogeny class. This
chapter defines it and proves the facts about it that the paper uses. It puts the modular curve $X_0(N)$ of \Cref{ch:modforms}
over $\Q$, builds its Jacobian $J_0(N)$ and the action of the Hecke algebra of \Cref{ch:hecke} on $J_0(N)$, on $H_1$ and on $H^1$ by
correspondences, and proves the identity $T_n^t=w_NT_nw_N$ that makes the pairing $\hat\delta_!$ of Section 4.1 of the paper
perfect, alternating and Hecke self-adjoint. It then constructs the optimal curve as a quotient of $J_0(N)$, characterises it,
and proves the four descriptions of $\degphi$ the paper uses: the degree of a map of curves, the integer $n$ with
$E\to J_0(N)\to E$ equal to $[n]$, the Petersson norm divided by the area of the period lattice (the formula of Zagier and
Watkins, with the Manin constant in the numerator), and the ideal of values of $\hat\delta_!$ on the lattice $M_{f,\lambda}$ of
Diamond, Flach and Guo (Proposition prop:pullback of the paper, proved here with every step expanded, and also over $\Z$). From the
last description the chapter rederives the divisibility $\degphi\mid r_E$ of Agashe, Ribet and Stein at every prime, and explains why the
equality can fail at primes $p$ with $p^2\mid N$ (the curve \curve{99a1}). The deep inputs (the $\Q$-structure of $X_0(N)$, the
Eichler-Shimura relation, modularity, multiplicity one, the theorems of Mazur and of Calegari and Emerton) are stated precisely with
references; the comparison of Ribet and Takahashi with Shimura curves is stated exactly, with its theorem numbers in the sources, and its
$\ell$-adic form is derived from it.

\section*{Core results at a glance}

Throughout, $E/\Q$ has conductor $N$, $f=f_E$ is its newform, $\phi_f\colon\T\to\Z$ the eigencharacter, $I_f=\ker\phi_f$, $\T$ the full Hecke algebra of level $N$ ($U_p$ included), $\ell$ a prime, and
$\varepsilon_N(f)=\pm1$ the eigenvalue of $w_N$ on $f$.

\begin{enumerate}[label=\textbf{\thechapter.\arabic*}]
\item \textbf{Conventions.} $W_N=\sm0{-1}N0$; slash $(f|_k\gamma)=(\det\gamma)^{k/2}(cz+d)^{-k}f(\gamma z)$; $\T$ acts by pushforward $t_*$ on $J_0(N)$, $H_1$, $T_\ell J_0(N)$ and by pullback $t^*$ on $S_2$, $H^1$;
unnormalised Petersson product; $\varphi^*\omega_E=c_\varphi\cdot2\pi if\,dz$ with $c_E>0$; ``optimal'' means $X_0(N)$-optimal; Eisenstein means $\bar\rho_{\mf m}$ reducible. \emph{(convention)}
\Cref{conv:moddeg:setup}.
\item \textbf{$X_0(N)$ over $\Q$.} $X_0(N)$ is a smooth projective geometrically connected curve over $\Q$, a coarse moduli space for $(E,C)$ with $C$ cyclic of order $N$; $j$ and the $w_d$ are defined over $\Q$,
$w_N(E,C)=(E/C,E[N]/C)$; good reduction at $p\nmid N$. \emph{(stated; reference)} \Cref{thm:moddeg:x0-over-q}.
\item \textbf{Cusps over $\Q$.} Cusps $a/c$, $c\mid N$, $a\bmod\gcd(c,N/c)$; those of denominator $c$ are defined over $\Q(\zeta_{\gcd(c,N/c)})$ and permuted simply transitively; $\infty$, $0$ are rational; all cusps are
rational iff $\gcd(c,N/c)\le2$ for all $c\mid N$. \emph{(proof sketched here; reference)} \Cref{prop:moddeg:cusps-over-q}.
\item \textbf{Rational points.} A $G_K$-stable cyclic $C$ gives a $K$-point of $Y_0(N)$, and conversely up to twist when $j\ne0,1728$; $Y_0(N)(\Q)\ne\emptyset$ iff $N\le19$ or $N\in\{21,25,27,37,43,67,163\}$.
\emph{(proved here; the list stated, Mazur, Kenku)} \Cref{prop:moddeg:rational-points}.
\item \textbf{Hecke correspondences.} $\mathcal C_\alpha=(X(\Gamma_\alpha),\pi_1,\pi_\alpha)$, $\Gamma_\alpha=\Gamma\cap\alpha^{-1}\Gamma\alpha$, $\pi_1(\Gamma_\alpha z)=\Gamma z$, $\pi_\alpha(\Gamma_\alpha z)=\Gamma\alpha z$; $T_p$, $U_p=\mathcal C_{\mathrm{diag}(1,p)}$.
\emph{(definition)} \Cref{def:moddeg:hecke-correspondence}.
\item \textbf{Properties.} $\mathcal C_\alpha$ depends only on $\Gamma\alpha\Gamma=\bigsqcup\Gamma\alpha\gamma_i$; $X(\Gamma_{\mathrm{diag}(1,p)})\simeq X_0(Np)$ with $\pi_1$, $\pi_\alpha$ becoming $z\mapsto pz$, $z\mapsto z$; on moduli
$T_p(E,C)=\sum_D(E/D,(C+D)/D)$; defined over $\Q$. \emph{(proved here; moduli sketched)} \Cref{prop:moddeg:hecke-correspondence}.
\item \textbf{Compatibility.} The form $\mathcal C_\alpha^*(f\,dz)=(\sum f|_2\alpha_i)\,dz$ is \Cref{ch:hecke}'s $T_p$ or $U_p$; $t\mapsto t^*\in\End H^1(X_0(N),\Z)$ and $t\mapsto t_*\in\End H_1$ are injective ring homomorphisms;
$\T$ is finite free over $\Z$. \emph{(proved here)} \Cref{prop:moddeg:hecke-compatibility}.
\item \textbf{The transpose.} $\mathcal C_\alpha^t\simeq\mathcal C_{(\det\alpha)\alpha^{-1}}$; $\mathcal C_{\beta\alpha\beta^{-1}}\simeq w_\beta\mathcal C_\alpha w_\beta^{-1}$; $W_N\mathrm{diag}(1,p)W_N^{-1}=\mathrm{diag}(p,1)$; hence $T_p^t\simeq w_NT_pw_N$,
$U_p^t\simeq w_NU_pw_N$, and $(t^*)^\dagger=w_N^*t^*w_N^*$ for the cup product, for every $t\in\T$; $T_p^t\simeq T_p$ for $p\nmid N$. \emph{(proved here)} \Cref{thm:moddeg:transpose}.
\item \textbf{The pairing $\hat\delta_!$.} $\hat\delta_!(x,y)=\langle x,w_N^*y\rangle$ is perfect and alternating on $H^1(X_0(N),R)$ ($R=\Z,\Zl,\Ql$), every $t^*$ is self-adjoint, and $H^1(X_0(N),\Zl)=\bigoplus_{\mf m}H^1_{\mf m}$
orthogonally with $\hat\delta_!$ perfect on each summand. \emph{(proved here)} \Cref{cor:moddeg:delta-hat}.
\item \textbf{$J_0(N)$.} $J_0(N)=\Jac(X_0(N)_\Q)$, $\iota(P)=[(P)-(\infty)]$. \emph{(definition)} \Cref{def:moddeg:j0}.
\item \textbf{Structure.} Abelian variety over $\Q$ of dimension $g=\dim S_2$, good reduction at $p\nmid N$; $J_0(N)(\C)$ is $S_2^\vee$ modulo $H_1(X_0(N),\Z)$; universal property of $\iota$; dimensions $1$, $1$, $1$, $2$, $7$, $9$, $6$, $19$, $49$, $109$, $145$ for
$N=11$, $14$, $20$, $37$, $77$, $99$, $121$, $140$, $539$, $726$, $980$. \emph{(proved here, given \Cref{ch:curves}; good reduction stated)} \Cref{thm:moddeg:j0}.
\item \textbf{Hecke on $J_0(N)$.} $t\mapsto t_*$ embeds $\T$ in $\End_\Q J_0(N)$, commuting with $G_\Q$ on $T_\ell J_0(N)$; $t^t=w_Ntw_N$ (in general not in $\T$); $I_fJ_0(N)$, $J_0(N)/I_fJ_0(N)$, $H^1[I_f]$ are the same
for the Albanese and the Picard actions. \emph{(proved here)} \Cref{prop:moddeg:hecke-on-j0}.
\item \textbf{$T_\ell J_0(N)$ versus $H^1$.} $H^1(X_0(N),\Zl)=(T_\ell J_0(N))^\vee\simeq T_\ell J_0(N)(-1)$; $\Psi=w_{N*}\circ\operatorname{PD}$ is a $\T_\ell$-isomorphism $H^1\to T_\ell J_0(N)$, Galois-equivariant after
the twist; freeness over $\T_{\mf m}$ transfers, and $H^1(X_0(N),\Fl)_{\mf m}[\mf m]\simeq J_0(N)[\mf m]$. \emph{(proved here, given the \'etale statements of \Cref{ch:curves})}
\Cref{prop:moddeg:tate-module-vs-h1}.
\item \textbf{The $f$-lattice.} $V_f=H^1(X_0(N),\Q)[I_f]$, $M_{f,\lambda}=H^1(X_0(N),\Zl)[I_f]$. \emph{(definition)} \Cref{def:moddeg:f-lattice}.
\item \textbf{Properties of $M_{f,\lambda}$.} $V_f\otimes\C=\C\omega_f\oplus\C\overline{\omega_f}$, also the eigenspace of the anemic algebra; $M_{f,\lambda}=H^1(X_0(N),\Z)[I_f]\otimes\Zl$ is saturated of rank $2$; $w_N^*$ acts on it by
$\varepsilon_N(f)$; it is $G_\Q$-stable. \emph{(proved here, given \Cref{ch:hecke})} \Cref{prop:moddeg:f-lattice}.
\item \textbf{Shimura's quotient.} $J_f=J_0(N)/I_fJ_0(N)$ is an abelian variety over $\Q$ of dimension $[K_f:\Q]$ with connected kernel $I_fJ_0(N)$; for rational $f$, $E_f=J_f$ is an elliptic curve, $\T$ acts by $\phi_f$,
$\varphi_{E_f}^*H^1(E_f,\Q)=V_f$; $L(E_f,s)=L(f,s)$. \emph{(proved here for rational $f$; sketched in general; the $L$-function stated)} \Cref{thm:moddeg:shimura-quotient}.
\item \textbf{Eichler-Shimura.} $(T_p)_*=\pi_p+\hat\pi_p$ on $J_0(N)_{\Fp}$ for $p\nmid N$; hence $\Frob_p^2-T_p\Frob_p+p=0$ on $T_\ell J_0(N)$, and $a_p(E_f)=a_p(f)$. \emph{(congruence relation stated; consequences
proved here)} \Cref{thm:moddeg:eichler-shimura}.
\item \textbf{Modularity.} Every $E/\Q$ of conductor $N$ has a newform $f_E$ of level $N$ with $a_p(f_E)=a_p(E)$ for all $p$; equivalently $E$ is isogenous to $E_{f_E}$. \emph{(stated; Wiles, Taylor-Wiles, BCDT,
Carayol)} \Cref{thm:moddeg:modularity}.
\item \textbf{Classes and newforms.} $E\mapsto f_E$ is a bijection from isogeny classes of conductor $N$ to rational newforms of level $N$. \emph{(proof sketched here)} \Cref{cor:moddeg:classes-and-newforms}.
\item \textbf{Modular parametrisation.} A nonconstant $\varphi\colon X_0(N)\to E$ over $\Q$ with $\varphi(\infty)=O$; equivalently a nonzero $h\colon J_0(N)\to E$. \emph{(definition)} \Cref{def:moddeg:parametrisation}.
\item \textbf{Optimal curves.} Every $h\colon J_0(N)\to E$ ($E$ in the class of $f$) kills $I_fJ_0(N)$; every parametrisation is $\psi\circ\varphi_{E_f}$ with $\psi$ an isogeny; $\psi$ an isomorphism $\Leftrightarrow$ $\ker\varphi_*$
connected $\Leftrightarrow$ $\varphi_*$ onto on $H_1$ $\Leftrightarrow$ $\varphi^*H^1$ saturated $\Leftrightarrow$ $\varphi^*$ injective on $E$ $\Leftrightarrow$ $\deg\varphi$ minimal; $E_f$ is the optimal curve. \emph{(proved here,
given Eichler-Shimura and modularity)} \Cref{thm:moddeg:optimal}.
\item \textbf{The modular degree.} $\degphi$ is the degree of $\varphi_E$; $E\to J_0(N)\to E$ is $[\degphi]$; $A\cap B\simeq(\Z/\degphi)^2$; the cup product on $\varphi_E^*H^1(E,\Z)$ has ideal of values $\degphi\Z$. \emph{(proved
here)} \Cref{prop:moddeg:degree-descriptions}.
\item \textbf{Along the class.} If $E[\ell]$ is irreducible, every curve of the class has minimal parametrisation degree of the same $\ell$-adic valuation as $\degphi$; false for \curve{11a} at $\ell=5$.
\emph{(proved here)} \Cref{prop:moddeg:valuation-in-class}.
\item \textbf{Cremona's convention.} Curve $1$ of each class is optimal, except \curve{990h} (optimal curve \curve{990h3}, $\degphi=1728$); optimality determined to conductor $400000$ and for $359{,}009$ classes
beyond. \emph{(convention; data)} \Cref{conv:moddeg:cremona-optimal}.
\item \textbf{Manin constant.} $\varphi^*\omega_E=c_\varphi\cdot2\pi if(z)\,dz$; $c_E=c_{\varphi_E}>0$. \emph{(definition)} \Cref{def:moddeg:manin-constant}.
\item \textbf{$c_\varphi$.} $\varphi^*\omega_E\in\C\omega_f\setminus0$, $c_\varphi\in\Q^\times$, $c_{\psi\circ\varphi_E}=c_Eu_\psi$ with $u_\psi u_{\hat\psi}=\deg\psi$. \emph{(proved here; integrality of $u_\psi$ stated)}
\Cref{prop:moddeg:manin-rational}.
\item \textbf{Known results.} $c_E\in\Z$ (Edixhoven); $c_E=1$ conjecturally; $p\mid c_E\Rightarrow p=2$ or $p^2\mid N$; $c_E=1$ for semistable $E$. \emph{(stated; reference)} \Cref{thm:moddeg:manin-known}.
\item \textbf{Zagier-Watkins.} For every modular parametrisation,
\[
\deg\varphi\cdot\vol(\Lambda_E)=4\pi^2c_\varphi^2\,\langle f,f\rangle\qquad(\text{unnormalised Petersson product}).
\]
\emph{(proved here)} \Cref{thm:moddeg:zagier-watkins}.
\item \textbf{Congruence number.} $r_E$: the largest $r$ with $f\equiv g\pmod r$, $g\in S_2(\Gamma_0(N);\Z)\cap f^\perp$; $\eta_f=\phi_f(\Ann_\T I_f)$. \emph{(definition)} \Cref{def:moddeg:congruence-number}.
\item \textbf{Congruence ideals.} $\T\otimes\Q=\Q\times\T'$, $I_f\cap\Ann I_f=0$, $\eta_f\ne0$; $\eta_f\Zl$ is computed in $\T_{\mf m}$; $r_E<\infty$. \emph{(proved here)} \Cref{prop:moddeg:congruence-ideal}.
\item \textbf{Lemma lem:crit.} $\T/(I_f+\Ann_\T f^\perp)\simeq\Z/r_E$, so $r_E\Z=\eta_f$, also over $\Zl$; a congruent eigenform in $f^\perp$ gives $\ell\mid r_E$. \emph{(stated; proved in \Cref{ch:modgalrep})}
\Cref{lem:moddeg:crit}.
\item \textbf{Agashe-Ribet-Stein.} For optimal $E$: $\degphi\mid r_E$, and $v_p(\degphi)=v_p(r_E)$ when $p^2\nmid N$. \emph{(stated; the divisibility proved here in \Cref{cor:moddeg:ars-divisibility})}
\Cref{thm:moddeg:ars}.
\item \textbf{$M_{f,\lambda}$ is the pullback.} For optimal $E$ and every prime $\ell$: $\varphi_E^*$ is injective with saturated image, $\varphi_E^*H^1(E,\Zl)=M_{f,\lambda}$, and $M_{f,\lambda}\simeq T_\ell E(-1)$ as Galois modules, so
$\ad M_{f,\lambda}\simeq\ad T_\ell E$. \emph{(proved here)} \Cref{prop:moddeg:mf-is-pullback}.
\item \textbf{$w_N$ on $M_{f,\lambda}$.} $\varphi_E\circ w_N=\varepsilon_N(f)\varphi_E+\varphi_E(0)$, so $w_N^*=\varepsilon_N(f)$ on $M_{f,\lambda}$. \emph{(proved here)} \Cref{lem:moddeg:wn-on-mf}.
\item \textbf{The degree as a congruence ideal.} For optimal $E$ and every prime $\ell$,
\[
\hat\delta_!(\varphi_E^*a,\varphi_E^*b)=\varepsilon_N(f)\,\degphi\,\langle a,b\rangle_E,
\]
so the ideal of values of $\hat\delta_!$ on $M_{f,\lambda}$ is $\degphi\Zl$; with the paper's description of \cite{DFG04}, $v_\lambda(\eta^\emptyset_{f,\lambda})=v_\ell(\degphi)$. \emph{(proved here)} \Cref{thm:moddeg:degree-is-eta}.
\item \textbf{Divisibility at every prime.} For every prime $\ell$, $v_\ell(\degphi)\le v_\ell(r_E)$, with equality if $H^1(X_0(N),\Zl)_{\mf m}$ is free over $\T_{\mf m}$. \emph{(proved here, given Lemma lem:crit)}
\Cref{cor:moddeg:ars-divisibility}.
\item \textbf{Lemma lem:hecke.} $v_\ell(\degphi)=v_\ell(\text{values of }\hat\delta_!\text{ on }M_{f,\lambda})\le v_\ell(r_E)$ always; $=v_\ell(r_E)$ if $\ell^2\nmid N$; $=v_\lambda(\eta^\emptyset_{f,\lambda})$ if $\ell\nmid2N$ and $\bar\rho$
irreducible. \emph{(proved here, given \Cref{thm:moddeg:ars} and the identification with \cite{DFG04})} \Cref{thm:moddeg:lem-hecke}.
\item \textbf{The maximal ideal of $(f,\ell)$.} $\mf m_{f,\ell}\ni\ell,I_f$; $\bar\rho_{\mf m}$ with $\tr\bar\rho_{\mf m}(\Frob_p)=T_p$; $\mf m_{f,\ell}$ is Eisenstein iff $E[\ell]$ is reducible. \emph{(definition; $\bar\rho_{\mf m}$ stated)}
\Cref{def:moddeg:maximal-ideal}.
\item \textbf{Multiplicity one.} For $\ell\nmid2N$ and $\mf m$ non-Eisenstein, $\dim_{\T/\mf m}J_0(N)[\mf m]=2$. \emph{(stated; Mazur, Ribet, Ribet-Stein)} \Cref{thm:moddeg:multiplicity-one}.
\item \textbf{Gorenstein.} Then $\T_{\mf m}$ is Gorenstein and $H^1(X_0(N),\Zl)_{\mf m}\simeq\T_{\mf m}^2\simeq T_\ell J_0(N)_{\mf m}$; so $v_\ell(r_E)=v_\ell(\degphi)$ for $\ell\nmid2N$, $E[\ell]$ irreducible. \emph{(proved here, given
multiplicity one)} \Cref{thm:moddeg:gorenstein}.
\item \textbf{Cuspidal subgroup.} $C_N\subseteq J_0(N)$ generated by cuspidal divisors of degree $0$. \emph{(definition)} \Cref{def:moddeg:cuspidal-subgroup}.
\item \textbf{Manin-Drinfeld.} $(T_p)_*(0)=(p+1)(0)$, $(T_p)_*(\infty)=(p+1)(\infty)$ for $p\nmid N$; $T_p-p-1$ is an isogeny; $[(0)-(\infty)]$ is a rational torsion point; $C_N$ is finite. \emph{(proved here; the
finiteness of $C_N$ stated)} \Cref{prop:moddeg:manin-drinfeld}.
\item \textbf{Mazur at prime level.} $J_0(N)(\Q)_{\tors}=C_N\simeq\Z/n$, $n=\operatorname{num}\frac{N-1}{12}$; $\T/I_{\mathrm{Eis}}\simeq\Z/n$; Eisenstein maximal ideals at odd $\ell$ exist iff $\ell\mid n$. \emph{(stated; reference)}
\Cref{thm:moddeg:mazur-prime-level}.
\item \textbf{Shimura subgroup.} $\mathcal S_N=\ker(J_0(N)\to J_1(N))$. \emph{(definition)} \Cref{def:moddeg:shimura-subgroup}.
\item \textbf{Its properties.} $\#\mathcal S_N$ is the saturation index of $\pi^*H^1(X_0(N),\Z)$; $\mathcal S_N$ is finite, multiplicative and Eisenstein; $\mathcal S_N[\mf m]=0$ for $\mf m$ non-Eisenstein. \emph{(proved here; (ii)
stated)} \Cref{prop:moddeg:shimura-subgroup}.
\item \textbf{The image of $0$.} $\varphi_E(0)\in E(\Q)_{\tors}$ and $\varphi_E(0)\equiv c_EL(E,1)\pmod{\Lambda_E}$; its order is the denominator of $c_EL(E,1)/\Omega$. \emph{(proved here, given the Mellin transform of
\Cref{ch:lfunctions})} \Cref{prop:moddeg:image-of-zero}.
\item \textbf{Shimura curves.} $X_0^D(M)$, its Jacobian isogenous to the $D$-new quotient of $J_0(DM)$; for $N=DM$ square-free, $E_D$ is the curve of the isogeny class that is an optimal quotient of $J_0^D(M)$ (in general not the $X_0(N)$-optimal one) and
$\delta_D=\deg(\varphi_{E,D}\colon X_0^D(M)\to E_D)$. \emph{(definition; stated)} \Cref{def:moddeg:shimura-curve}.
\item \textbf{Ribet-Takahashi.} For $N=dpqm$ square-free, $d$ with an even number of prime factors: $\delta^d(pqm)/\delta^{dpq}(m)=n_p(E_2)n_q(E_1)/i_q^2$, $n_r=v_r(\Delta)=\#\Phig_r$ (not the Tamagawa
number), $i_q$ the order of the image of $\Phig_q(J_0^d(pqm))\to\Phig_q(E_1)$; if $E$ is alone in its isogeny class, $\degphi/\delta_D=\prod_{r\mid D}n_r$; if $E[\ell]$ is irreducible (any $\ell$),
$v_\ell(\degphi)=v_\ell(\delta_D)+\sum_{r\mid D}v_\ell(n_r)$. \emph{(stated; \cite[Theorems 1, 2]{RibetTakahashi97}, \cite[Theorem 3.2, Corollary 3.5]{Takahashi01}; the $\ell$-adic form sketched here)}
\Cref{thm:moddeg:ribet-takahashi}; sources, error term and example \curve{57c1} in \Cref{rem:moddeg:ribet-takahashi}.
\item \textbf{Watkins's conjecture.} $2^{\rank E(\Q)}\mid\degphi$; true for all $38{,}042$ optimal curves of conductor below $10{,}000$. \emph{(conjecture; data checked)} \Cref{conj:moddeg:watkins}.
\item \textbf{Calegari-Emerton.} Odd $\degphi$ implies at most two odd bad primes, even analytic rank, and a rational $2$-torsion point, or prime conductor with supersingular reduction at $2$ and $E(\R)$
connected, or CM with $N\in\{27,32,49,243\}$. \emph{(stated; reference; data checked)} \Cref{thm:moddeg:calegari-emerton}.
\end{enumerate}

\section*{Where the paper uses this chapter}

\begin{itemize}
\item \textbf{Section 1 (Introduction) and the abstract.} The $X_0(N)$-optimal curve, the modular parametrisation
$\varphi_E\colon X_0(N)\to E$ and ``minimal degree'' (\Cref{def:moddeg:parametrisation,thm:moddeg:optimal,prop:moddeg:degree-descriptions});
the sentence on the choice of curve when $E[\ell]$ is irreducible (\Cref{prop:moddeg:valuation-in-class}); computability
(\Cref{rem:moddeg:computing-the-degree}); the $abc$ connection (\Cref{rem:moddeg:abc-size}); the Petersson norm and Zagier, Watkins
(\Cref{thm:moddeg:zagier-watkins}); Agashe, Ribet and Stein and the congruence number $r_E$
(\Cref{def:moddeg:congruence-number,thm:moddeg:ars,cor:moddeg:ars-divisibility}).
\item \textbf{Section 1.5 (History).} The parametrisations by Shimura curves of Ribet and Takahashi
(\Cref{def:moddeg:shimura-curve,thm:moddeg:ribet-takahashi,rem:moddeg:ribet-takahashi}).
\item \textbf{Section sec:proof (proof of Theorem thm:qual).} The congruence number and Lemma lem:crit (\Cref{lem:moddeg:crit}, proved in
\Cref{ch:modgalrep}); ``it suffices to prove $\ell\mid r$'' (\Cref{thm:moddeg:ars}); the remark on \curve{99a1} and \curve{726a1}
(\Cref{ex:moddeg:99a1-and-726a1,ex:moddeg:99-at-3}).
\item \textbf{Section 4.1 (sec:dfg).} $H^1(X_0(M^\Sigma),\Zl)$ with its Hecke action, the pairing $\hat\delta_!$ and the sentence
``$w_{M^\Sigma}T_nw_{M^\Sigma}^{-1}$ is the adjoint of $T_n$ for the cup product'' (\Cref{prop:moddeg:hecke-compatibility,thm:moddeg:transpose,cor:moddeg:delta-hat}); the
lattice $M_{f,\lambda}$ (\Cref{def:moddeg:f-lattice}) and the checklist for the identification with Diamond, Flach and Guo's objects, including Proposition prop:integral and its
Remark (\Cref{rem:moddeg:dfg-normalisation}).
\item \textbf{Proposition prop:pullback.} \Cref{prop:moddeg:mf-is-pullback}, \Cref{lem:moddeg:wn-on-mf} and \Cref{thm:moddeg:degree-is-eta}, with a
Commentary on the paper's statement and proof; the paragraph after the proposition is \Cref{prop:moddeg:mf-is-pullback}(d).
\item \textbf{Lemma lem:hecke, Remark rem:free and the paragraph between them.} \Cref{thm:moddeg:lem-hecke}, the divisibility at every prime
(\Cref{cor:moddeg:ars-divisibility}), multiplicity one and the Gorenstein property
(\Cref{def:moddeg:maximal-ideal,thm:moddeg:multiplicity-one,thm:moddeg:gorenstein,rem:moddeg:rem-free-after-39}), and the identification of
$T_\ell J_0(N)$ with $H^1(X_0(N),\Zl)$ as Hecke modules (\Cref{prop:moddeg:tate-module-vs-h1}).
\item \textbf{Section 4.2 and the proof of Theorem thm:main.} $M_{f,\lambda}\simeq H^1(E,\Zl)\simeq T_\ell E(-1)$ and $\ad M_{f,\lambda}\simeq\ad T_\ell E$
(\Cref{prop:moddeg:mf-is-pullback}(d)); the equality in (eq:identity) (\Cref{thm:moddeg:lem-hecke}).
\item \textbf{Remark rem:degree.} The formula of Zagier and Watkins with $c_E^2$ in the numerator (\Cref{thm:moddeg:zagier-watkins,ex:moddeg:zw-numerics}).
\item \textbf{Section sec:data.} Optimality in Cremona's tables and the class \curve{990h} (\Cref{conv:moddeg:cremona-optimal}); \texttt{ellmoddegree}
(\Cref{rem:moddeg:computing-the-degree}).
\item \textbf{Conjecture conj:eis, Section sec:eis, Open question 4.} The cuspidal subgroup, Mazur's Eisenstein ideal, the Shimura subgroup and
the image of the cusp $0$ (\Cref{prop:moddeg:manin-drinfeld,thm:moddeg:mazur-prime-level,def:moddeg:shimura-subgroup,prop:moddeg:image-of-zero,rem:moddeg:eisenstein-and-the-paper}).
\item \textbf{Section sec:two, Open question 5.} Watkins's conjecture and the theorem of Calegari and Emerton (\Cref{conj:moddeg:watkins,thm:moddeg:calegari-emerton}).
\end{itemize}

\section{Overview and standing notation}\label{sec:moddeg:overview}

This section fixes the conventions of the chapter and says which section answers which sentence of the paper. The paper
uses the chapter in six places: the optimal curve and its parametrisation (Sections \ref{sec:moddeg:x0q} to \ref{sec:moddeg:optimal}); the
Petersson norm and the Manin constant (\Cref{sec:moddeg:manin}); the congruence number and the theorem of Agashe, Ribet and Stein
(\Cref{sec:moddeg:congruence}); the cohomology of $X_0(N)$ with its Hecke action and the pairing $\hat\delta_!$ (\Cref{sec:moddeg:hecke,sec:moddeg:jacobian});
the identification of the congruence ideal of Diamond, Flach and Guo with $\degphi$ (\Cref{sec:moddeg:degree-eta,sec:moddeg:mult-one}); and the Eisenstein
case, the Shimura curves and the prime $2$ (\Cref{sec:moddeg:eisenstein,sec:moddeg:shimura-two}).

\begin{convention}[standing conventions of this chapter]\label{conv:moddeg:setup}
\leavevmode
\begin{enumerate}[label=(C\arabic*)]
\item \emph{Curves and levels.} $N\ge1$; $X_0(N)$ is the compact Riemann surface of \Cref{ch:modforms} and, by \Cref{thm:moddeg:x0-over-q}, a smooth
projective curve over $\Q$; $Y_0(N)\subset X_0(N)$ is the complement of the cusps; $\infty$ and $0$ are the cusps $\Gamma_0(N)\infty$ and $\Gamma_0(N)0$;
$g=g(X_0(N))$; $\Gamma=\Gamma_0(N)$. $W_N=\sm0{-1}N0$ and $w_N$ is the involution $z\mapsto-1/(Nz)$ of $X_0(N)$ and the operator $f\mapsto f|_2W_N$
(\Cref{ch:modforms}). The slash operator is $(f|_k\gamma)(z)=(\det\gamma)^{k/2}(cz+d)^{-k}f(\gamma z)$, so that in weight $2$
$(f|_2\gamma)(z)\,dz=\gamma^*(f(z)\,dz)$. The letter $q$ is the variable $e^{2\pi iz}$ of $q$-expansions and also a bad prime of $E$
(Appendix B, B-C1); the two never occur in one display.
\item \emph{Hecke actions.} $\T=\Z[T_n:n\ge1]\subset\End_\C S_2(\Gamma_0(N))$ is the full Hecke algebra of \Cref{ch:hecke}, the operators $U_p=T_p$ for
$p\mid N$ included; $\T_\ell=\T\otimes\Zl$. A correspondence $\mathcal C$ acts on $J_0(N)$, on $H_1(X_0(N),\Z)$ and on $T_\ell J_0(N)$ by pushforward
$\mathcal C_*$ (``Albanese'') and on $H^0(X_0(N),\Omega^1)=S_2(\Gamma_0(N))$, $H^1(X_0(N),\Z)$ and $H^1(X_0(N),\Zl)$ by pullback $\mathcal C^*$, the transpose of
$\mathcal C_*$ (\Cref{ch:curves}). For $t\in\T$ the two resulting actions are written $t_*$ and $t^*$ (\Cref{prop:moddeg:hecke-compatibility}); every
statement ``$\T$ acts on \dots'' names which.
\item \emph{Duals and adjoints} (Appendix B, B-C4). $M^\vee=\Hom(M,\Zl)$ (or $\Hom(M,\Z)$); $A^\dagger$ is the adjoint of an operator for a pairing;
a star after a map is a pullback. The paper's Pontryagin duals occur only inside \texttt{paperlink} environments.
\item \emph{Petersson product, Manin constant} (B-C8, B-C9). $\langle f,g\rangle=\int_{\Gamma_0(N)\backslash\HH}f\bar g\,dx\,dy$ is unnormalised in weight
$2$, and $\langle f,f\rangle=\frac i2\int_{X_0(N)}\omega_f\wedge\overline{\omega_f}$ with $\omega_f=f(z)\,dz$. For an elliptic curve $E/\Q$, $\omega_E$ is the N\'eron
differential of a global minimal model, $\Lambda_E$ its period lattice and $\vol(\Lambda_E)=\frac i2\int_{E(\C)}\omega_E\wedge\bar\omega_E$; $\varphi^*\omega_E=c_\varphi\cdot2\pi if(z)\,dz$
defines $c_\varphi$, and $c_E>0$. PARI's \texttt{mfpetersson} divides by $\mu(N)=[\SL_2(\Z):\Gamma_0(N)]$. The Manin constant is never a bare $c$.
\item \emph{Frobenius} (B-C2(g), B-C3). On $T_\ell J$ the arithmetic Frobenius at $p\nmid N\ell$ acts as the $p$-power Frobenius endomorphism
$\pi_p$ of the reduction; $\hat\pi_p=p\pi_p^{-1}$ is the Verschiebung; on $H^1=(T_\ell J)^\vee$ the geometric Frobenius has the same characteristic
polynomial (\Cref{ch:curves}).
\item \emph{Optimal curves, labels, coefficients.} ``Optimal'' means $X_0(N)$-optimal (\Cref{thm:moddeg:optimal}). Curve labels are Cremona's;
\cite{ARS12} and \cite{Cremona97} print \curve{99A1} for \curve{99a1}. The newform of an elliptic curve has rational coefficients, so $\lambda=\ell$ and
$\mathcal O=\Zl$ throughout, and the paper's $v_\lambda$ is $v_\ell$.
\item \emph{Eisenstein.} A maximal ideal $\mf m$ of $\T$ (or of $\T_\ell$) is \emph{Eisenstein} if the residual representation $\bar\rho_{\mf m}$ of
\Cref{def:moddeg:maximal-ideal} is reducible. Mazur's \emph{Eisenstein ideal} $I_{\mathrm{Eis}}$ at prime level is a specific ideal
(\Cref{thm:moddeg:mazur-prime-level}). The letter $D$ is used for the quaternion discriminant of a Shimura curve (\Cref{sec:moddeg:shimura-two}) and,
in quotations, for the deficit of Conjecture conj:eis; $\mathcal S_N$ is the Shimura subgroup.
\item \emph{Typography.} Double names take plain hyphens: Eichler-Shimura, Agashe-Ribet-Stein, Zagier-Watkins, Manin-Drinfeld,
Ribet-Takahashi, Calegari-Emerton, Deligne-Rapoport, Boston-Lenstra-Ribet, Jacquet-Langlands.
\end{enumerate}
\end{convention}

\noindent\textbf{Standing notation.} Throughout, $E/\Q$ is an elliptic curve of conductor $N$, $f=f_E\in S_2(\Gamma_0(N))$ its newform
(\Cref{thm:moddeg:modularity}), $\phi_f\colon\T\to\Z$, $T_n\mapsto a_n(f)$, its eigencharacter (an upright $\phi$; a curly $\varphi$ is always a map of
curves), $I_f=\ker\phi_f$ (the paper's $I$ in Lemma lem:crit and its $\wp$ in Lemma lem:ddt417), and $f^\perp$ the Petersson orthogonal
complement of $f$ in $S_2(\Gamma_0(N))$, which is $\T$-stable (\Cref{ch:hecke}). $\varepsilon_N(f)=\pm1$ is the eigenvalue of $w_N$ on $f$.

\medskip\noindent\textbf{Three conventions that matter silently.} (1) The factor $(\det\gamma)^{k/2}$ in the slash operator is chosen so that
scalar matrices act trivially in weight $2$: for $\gamma=-N\cdot1$ one gets $(N^2)^{1}(-N)^{-2}f=f$. Since $W_N^2=-N\cdot1$, this is what makes
$f\mapsto f|_2W_N$ an involution of $S_2(\Gamma_0(N))$, so that its eigenvalues $\varepsilon_N(f)$ are $\pm1$ and not $\pm\sqrt{\text{something}}$. It also makes
$(f|_2\gamma)(z)\,dz$ the pullback of $f(z)\,dz$ along $z\mapsto\gamma z$, because $d(\gamma z)/dz=\det\gamma/(cz+d)^2$; every comparison below between Hecke operators on forms and on
cohomology rests on that identity. (2) The sign $\varepsilon_N(f)$ is the eigenvalue of $w_N$, not the root number: the root number of $L(f,s)$ is $-\varepsilon_N(f)$
(\Cref{ch:hecke,ch:lfunctions}). For \curve{11a1}, of rank $0$ and root number $+1$, $\varepsilon_{11}(f)=-1$; for \curve{37a1}, of rank $1$ and root number $-1$, $\varepsilon_{37}(f)=+1$
(PARI's \texttt{ellrootno}). (3) The Hecke algebra acts on two kinds of objects, covariantly (pushforward, on $J_0(N)$ and $H_1$) and contravariantly (pullback, on forms and $H^1$).
The two actions are transposes of each other; for $T_n$ with $(n,N)=1$ the distinction is harmless, but for $U_p$, $p\mid N$, the transpose is a different operator
(\Cref{thm:moddeg:transpose}), and this is the reason the pairing $\hat\delta_!$ of the paper contains $w_N$.

\begin{center}\small
\begin{tabular}{p{7.2cm}p{6.3cm}}
\toprule
sentence or object of the paper & where in this chapter\\
\midrule
``the $X_0(N)$-optimal curve, for which the modular parametrisation $\varphi_E\colon X_0(N)\to E$ has minimal degree'' &\Cref{thm:moddeg:optimal,prop:moddeg:degree-descriptions}\\
``$\degphi$ is $\langle f,f\rangle$ divided by the area of the period lattice'' & \Cref{thm:moddeg:zagier-watkins}\\
``$\degphi$ divides the \emph{congruence number} $r_E$'' &\Cref{thm:moddeg:ars,cor:moddeg:ars-divisibility}\\
$\hat\delta_!$ perfect, alternating, $\T$ self-adjoint & \Cref{thm:moddeg:transpose,cor:moddeg:delta-hat}\\
Proposition prop:pullback & \Cref{prop:moddeg:mf-is-pullback,thm:moddeg:degree-is-eta}\\
Lemma lem:hecke, Remark rem:free & \Cref{thm:moddeg:lem-hecke,thm:moddeg:gorenstein,rem:moddeg:rem-free-after-39}\\
\curve{99A1}: degree $4$, congruence number $12$ & \Cref{ex:moddeg:99a1-and-726a1,ex:moddeg:99-at-3}\\
Conjecture conj:eis, Open question 4 & \Cref{sec:moddeg:eisenstein}\\
Section 1.5, Section sec:two & \Cref{sec:moddeg:shimura-two}\\
\bottomrule
\end{tabular}
\end{center}

\section{The modular curve \texorpdfstring{$X_0(N)$}{X0(N)} over \texorpdfstring{$\Q$}{Q}}\label{sec:moddeg:x0q}

\Cref{ch:modforms} builds $X_0(N)$ as a compact Riemann surface. The paper needs more: the modular parametrisation $\varphi_E\colon X_0(N)\to E$ is
a morphism of curves over $\Q$, its cohomology carries an action of $G_\Q$ (Section 4.1 of the paper works with $H^1$ of a curve over $\Q$), and the cusps
$\infty$ and $0$ are rational points (the base point of $\varphi_E$, and the point whose image is studied in Open question 4). This section states the
existence of the $\Q$-model, describes the cusps over $\Q$, records what is known about rational points, and works out the smallest levels.

\begin{theorem}[$X_0(N)$ over $\Q$; stated]\label{thm:moddeg:x0-over-q}
\leavevmode
\begin{enumerate}[label=(\roman*)]
\item There is a smooth projective geometrically connected curve $X_0(N)_\Q$ over $\Q$ whose complex points are the Riemann surface $X_0(N)$, such
that $Y_0(N)_\Q$ is a coarse moduli space for pairs $(E,C)$ of an elliptic curve and a cyclic subgroup of order $N$: for every algebraically closed
field $K$ of characteristic $0$, $Y_0(N)(K)$ is in bijection with the isomorphism classes of such pairs over $K$, compatibly with $\Aut(K)$, and for
$K=\C$ the bijection is the one of \Cref{ch:modforms}, $\Gamma_0(N)\tau\mapsto(\C/(\Z+\Z\tau),\langle\frac1N\rangle)$.
\item The forgetful map $X_0(N)\to X(1)=\PP^1_j$, $(E,C)\mapsto j(E)$, and the involutions $w_d$ for $d\,\|\,N$ are defined over $\Q$; on moduli
$w_N(E,C)=(E/C,E[N]/C)$.
\item The function field of $X_0(N)_\Q$ is $\Q(j,j_N)$, $j_N(z)=j(Nz)$, and $\Phi_N(j,j_N)=0$ for the modular polynomial $\Phi_N\in\Z[X,Y]$.
\item $X_0(N)_\Q$ has good reduction at every prime $p\nmid N$ (Igusa; Deligne and Rapoport). For $p\,\|\,N$ it has a semistable model over $\Zp$
(Deligne and Rapoport) whose special fibre is two copies of $X_0(N/p)_{\Fp}$ meeting in ordinary double points at the supersingular points; at such a point the
complete local ring over $W(\bar\F_p)$ is $W(\bar\F_p)[[x,y]]/(xy-p^k)$ with $k=\frac12\#\Aut$ of the corresponding supersingular curve with its level structure,
so $k\in\{1,2,3\}$ when $p\ge5$ (for $N=p\in\{2,3\}$ it can reach $12$ or $6$), and the model is regular except at the supersingular points whose pair has automorphisms other than $\pm1$ (possible only for $j=0$ or $1728$),
which Edixhoven resolves.
\end{enumerate}
\end{theorem}

We do not prove this; see \cite[Sections 6.7, 7.2]{Shimura71}, \cite[Sections 7.7, 8.2 to 8.6]{DiamondShurman}, \cite[Sections 8 and 9]{DiamondIm},
\cite{DeligneRapoport73} and \cite{KatzMazur85} (section numbers not verified), and \cite{Edixhoven90} for the local rings at the supersingular points and the
minimal regular model. The idea of (i) and (iii): $j$ and $j_N$ are $\Gamma_0(N)$-invariant
functions on $\HH$ with $q$-expansions in $\Z((q))$, they generate the function field of $X_0(N)$ over $\C$ (\Cref{ch:modforms}), and the subfield
$\Q(j,j_N)$ is linearly disjoint from $\C$ over $\Q$ (the $q$-expansion principle: a function in $\C(j,j_N)$ whose $q$-expansion has rational
coefficients lies in $\Q(j,j_N)$); the smooth projective curve over $\Q$ with this function field is $X_0(N)_\Q$. The polynomial $\Phi_N$ has integer
coefficients because the conjugates $j(\gamma z)$, $\gamma$ running over $\SL_2(\Z)\backslash\{\alpha\in M_2(\Z):\det\alpha=N,\ \alpha\text{ primitive}\}$,
are permuted by $\SL_2(\Z)$ and have $q$-expansions with coefficients in $\Z[\zeta_N]$, permuted by $\Gal(\Q(\zeta_N)/\Q)$ (Kronecker). From now on
$X_0(N)$ also denotes the curve over $\Q$.

Two words in (i) deserve comment. A \emph{coarse} moduli space has the right points over algebraically closed fields, but over a non-closed field $K$ a $K$-rational point need
not come from a pair defined over $K$, and two pairs over $K$ with the same point need not be isomorphic over $K$; the obstruction is the automorphism group of the pair, which is
$\{\pm1\}$ except for $j=0$ and $j=1728$. \Cref{prop:moddeg:rational-points} makes this precise. The simplest case is $N=1$: $Y_0(1)=\mathbb A^1_j$ and $X_0(1)=\PP^1_j$, and two
elliptic curves over $\Q$ with the same $j$-invariant give the same rational point although they need not be isomorphic over $\Q$: for $j\ne0,1728$ they are quadratic
twists of each other, and for example $y^2=x^3-x$ and its quadratic twist by $2$, $y^2=x^3-4x$, both with $j=1728$, are not isomorphic over $\Q$ (an isomorphism would have
to be $x\mapsto u^2x$ with $u^4=4$). Part (iv) is needed in two places: good reduction at $p\nmid N$ gives the Eichler-Shimura relation
(\Cref{thm:moddeg:eichler-shimura}) and the unramifiedness of $T_\ell J_0(N)$ at $p$; the semistable model at $p\,\|\,N$ is what makes $J_0(N)$ semistable there, which is used for
$X_0(37)$ in \Cref{ex:moddeg:small-levels}(c).

\begin{proposition}[cusps over $\Q$]\label{prop:moddeg:cusps-over-q}
\leavevmode
\begin{enumerate}[label=(\roman*)]
\item The cusps of $X_0(N)$ are the classes of $a/c$ with $c\mid N$, $c\ge1$, $\gcd(a,c)=1$, $a$ taken modulo $\gcd(c,N/c)$; there are
$\varphi(\gcd(c,N/c))$ of denominator $c$ ($\varphi$ is Euler's function in this proposition only), and $\sum_{c\mid N}\varphi(\gcd(c,N/c))$ in all.
\item The cusps of denominator $c$ are defined over $\Q(\zeta_{\gcd(c,N/c)})$, and $\Gal(\Q(\zeta_{\gcd(c,N/c)})/\Q)$ permutes them simply transitively.
\item In particular $\infty$ ($c=N$) and $0$ ($c=1$) are $\Q$-rational, and every cusp is $\Q$-rational if and only if $\gcd(c,N/c)\le2$ for every $c\mid N$.
\end{enumerate}
\end{proposition}

\begin{proof}[Sketch of proof]
(i) is the description of the cusps of $\Gamma_0(N)$ in \Cref{ch:modforms}. (ii) Near the cusps the moduli problem degenerates to the Tate curve
$\mathbb G_m/q^{\Z}$ over $\Q((q^{1/N}))$; the cusp of $a/c$ corresponds to a cyclic subgroup of order $N$ generated by an element $\zeta_N^{u}q^{v/N}$, and
the field of definition of that subgroup modulo the automorphisms of the Tate curve with its subgroup is $\Q(\zeta_{\gcd(c,N/c)})$, the element
$\sigma_t\colon\zeta\mapsto\zeta^t$ sending the cusp of $a/c$ to that of $a'/c$ with $a'\equiv t^{\pm1}a\pmod{\gcd(c,N/c)}$ (the sign depends on the normalisation of the
moduli interpretation at the cusps, and only the field of definition is used below). See \cite{Ogg73}, \cite[Section 6.3]{Stein} and
\cite[Chapter 3, exercises]{DiamondShurman} (not verified). (iii) $\varphi(d)=1$ iff $d\le2$; for $c=N$ and $c=1$, $\gcd(c,N/c)=1$.
\end{proof}

\begin{example}[cusps at three levels]\label{ex:moddeg:cusps}
By (i), $X_0(20)$ has $6$ cusps ($\gcd(c,20/c)=1,2,1,1,2,1$ for $c=1,2,4,5,10,20$), all rational. $X_0(99)$ has $8$ cusps: one for each of $c=1,9,11,99$, and two
for each of $c=3$ and $c=33$ ($\gcd=3$), each pair conjugate over $\Q(\zeta_3)=\Q(\sqrt{-3})$. $X_0(121)$ has $12$ cusps: $\infty$, $0$ and ten of denominator $11$,
conjugate over $\Q(\zeta_{11})$. These counts agree with \Cref{ch:modforms}.
\end{example}

\begin{proposition}[rational points of $Y_0(N)$]\label{prop:moddeg:rational-points}
\leavevmode
\begin{enumerate}[label=(\roman*)]
\item Let $K$ be a field of characteristic $0$. If $E/K$ has a cyclic subgroup $C$ of order $N$ stable under $G_K$, the point $(E,C)$ of $Y_0(N)$ is
$K$-rational. Conversely a $K$-rational point of $Y_0(N)$ with $j\ne0,1728$ is represented by such a pair over $K$, unique up to quadratic twist.
\item (Mazur; Kenku; stated.) $Y_0(N)(\Q)\ne\emptyset$ if and only if $N\le19$ or $N$ is one of $21$, $25$, $27$, $37$, $43$, $67$, $163$. $Y_0(N)(\Q)$ is infinite exactly when $X_0(N)$
has genus $0$ ($N\le10$ or $N\in\{12,13,16,18,25\}$), and the rational points of $Y_0(27)$, $Y_0(43)$, $Y_0(67)$ and $Y_0(163)$ are CM points.
\end{enumerate}
\end{proposition}

\begin{proof}
(i) If $E$ and $C$ are defined over $K$, then for $\sigma\in\Aut(\bar K/K)$ the pair $(E^\sigma,C^\sigma)$ is $(E,C)$, so its point is fixed by $\sigma$
(\Cref{thm:moddeg:x0-over-q}(i)), and a point of a curve over $K$ fixed by $\Aut(\bar K/K)$ is $K$-rational. Conversely, let $P\in Y_0(N)(K)$ with
$j(P)\ne0,1728$, and let $E_0/K$ be any curve with $j(E_0)=j(P)$ (\Cref{ch:ellcurves}). Over $\bar K$, $P$ is $(E_0,C)$ for a cyclic $C\subset E_0$ of order $N$.
For $\sigma$ fixing $K$, $(E_0,C^\sigma)\simeq(E_0,C)$; as $\Aut(E_0)=\{\pm1\}$ and $-1$ fixes every subgroup, $C^\sigma=C$. Two pairs over $K$ with the same point
are isomorphic over $\bar K$ by an isomorphism unique up to sign, so they differ by a quadratic twist (\Cref{ch:ellcurves}). (ii) We do not prove this; see
\cite{Mazur78} for prime $N$, \cite{Kenku79,Kenku80a,Kenku80b,Kenku81} for the remaining levels and \cite{Ligozat75} for $N=27$; the list is
\cite[Theorem 1]{Kenku} (checked against the published text: $Y_0(N)(\Q)$ is empty except for $N\le10$, $N=12$, $13$, $16$, $18$, $25$, of genus $0$, and $N=11$, $14$, $15$, $17$,
$19$, $21$, $27$, $37$, $43$, $67$, $163$). Kenku credits the levels of genus one to Ligozat, Ogg and Joly and does not discuss complex multiplication; the CM statement is not in
his paper.
\end{proof}

A consequence used in \Cref{ch:ellcurves} (the proof of Proposition prop:no9 of the paper): a rational cyclic $27$-isogeny is a rational point of $Y_0(27)$, so
it exists only for CM curves \cite{Ligozat75}.

\begin{example}[the smallest levels]\label{ex:moddeg:small-levels}
\leavevmode
\begin{enumerate}[label=(\alph*)]
\item $X_0(11)$ has genus $1$ and the rational point $\infty$, so it is an elliptic curve over $\Q$ with origin $\infty$, and $J_0(11)=X_0(11)$. Its period
lattice is the lattice of periods of $2\pi if_{11}(z)\,dz$, $f_{11}=\eta(z)^2\eta(11z)^2$ (\Cref{ch:modforms}), and its $j$-invariant, computed from the period
ratio (\Cref{ch:ellcurves}), is $-122023936/161051$: $X_0(11)$ is the curve \curve{11a1} $=[0,-1,1,-10,-20]$ (the other two curves of the class have
$j=-52893159101157376/11$ and $-4096/11$). The modular parametrisation of \curve{11a1} is the identity, of degree $1$.
\item In the same way $X_0(14)$ and $X_0(20)$, of genus $1$, are the curves \curve{14a1} and \curve{20a1}, with parametrisations of degree $1$.
\item $X_0(37)$ has genus $2$, and $w_{37}$ has $2$ fixed points (\Cref{ch:modforms}); by Riemann-Hurwitz for the quotient map of degree $2$, $2g-2=2(2g'-2)+\sum_P(e_P-1)$, and each fixed point of the involution is a ramification point with
$e_P=2$ while all other points are unramified; so $2\cdot2-2=2(2g'-2)+2$, and $X_0(37)/w_{37}$ has genus
$g'=1$. With origin the image of $\infty$, it is an elliptic curve over $\Q$ (the quotient map is defined over $\Q$ by \Cref{thm:moddeg:x0-over-q}(ii)), and the
pullback of its holomorphic differential is $w_{37}$-invariant. $S_2(\Gamma_0(37))$ is spanned by the newforms $f_{37a}$ and $f_{37b}$ with $\varepsilon_{37}=+1$ and $-1$, so the
pullback is a multiple of $f_{37a}\,dz$. The quotient has good reduction outside $37$ (it is a quotient of $J_0(37)$, \Cref{thm:moddeg:j0}), and multiplicative reduction at $37$: $J_0(37)$
has semistable reduction there, by the semistable model of \Cref{thm:moddeg:x0-over-q}(iv), and a quotient of an abelian variety with semistable reduction has
semistable reduction (Grothendieck); so its conductor divides $37$, and it is $37$, no elliptic curve over $\Q$ having conductor $1$; and the pullback of the differential of a curve of conductor $37$ along a modular parametrisation is a
multiple of its newform (\Cref{prop:moddeg:manin-rational}). So the quotient lies in the class \curve{37a} (\Cref{cor:moddeg:classes-and-newforms}). That class has one curve,
\curve{37a1}; so $X_0(37)\to X_0(37)/w_{37}\simeq\curve{37a1}$ is a modular parametrisation of degree $2$, and it is the optimal one (\Cref{ex:moddeg:37}).
\end{enumerate}
The identifications in (a) and (b) are computations from the periods; they are confirmed by Cremona's tables \cite{Cremona97,ecdata}, which give
$\deg\varphi=1$ for \curve{11a1}, \curve{14a1}, \curve{20a1} and $2$ for \curve{37a1}.
\end{example}

\begin{paperlink}\label{pl:moddeg:over-q}
The abstract of the paper: ``$\varphi_E\colon X_0(N)\to E$ its modular parametrisation'' is a morphism of curves over $\Q$ sending the rational cusp $\infty$ to $O$
(\Cref{thm:moddeg:x0-over-q}, \Cref{prop:moddeg:cusps-over-q}(iii)). Section 4.1 uses $H^1(X_0(M^\Sigma),\Zl)$ with its Galois action: this is the
\'etale cohomology of the curve over $\Q$ (\Cref{prop:moddeg:tate-module-vs-h1}). The rational isogenies of Section sec:eis and the cyclic $27$-isogeny of
Proposition prop:no9 are rational points of $Y_0(\ell)$ and $Y_0(27)$ (\Cref{prop:moddeg:rational-points}).
\end{paperlink}

The curve is now defined over $\Q$, with rational cusps $\infty$ and $0$. What it still lacks is the structure that links it to modular forms beyond the
identification $S_2=H^0(X_0(N),\Omega^1)$: the Hecke operators, which so far act only on $q$-expansions. The next section realises them geometrically.

\section{Hecke correspondences and their adjoints}\label{sec:moddeg:hecke}

The Hecke operators of \Cref{ch:hecke} act on $q$-expansions. To act on $J_0(N)$, on $H_1(X_0(N),\Z)$ and on $H^1(X_0(N),\Zl)$, as Section 4.1 of the paper and
Remark rem:free require, they must come from correspondences on the curve. This section defines the Hecke correspondences, shows that they induce
\Cref{ch:hecke}'s operators on differentials, so that the same ring $\T$ acts on forms and on cohomology, and proves the identity the paper uses in
Section 4.1: \emph{the adjoint of $T_n$ for the cup product is $w_NT_nw_N$}, with $U_p$ included. \Cref{ch:curves} supplies the general theory (correspondences,
their actions, the adjoint of a correspondence is its transpose); the input specific to $X_0(N)$ is a two-line matrix computation.

For $\alpha\in M_2(\Z)$ with $\det\alpha=m>0$ put $\Gamma_\alpha=\Gamma\cap\alpha^{-1}\Gamma\alpha$. It has finite index in $\Gamma$: if $\gamma=1+NmX\in\Gamma(Nm)$, then
$\alpha\gamma\alpha^{-1}=1+Nm\,\alpha X\alpha^{-1}=1+N\alpha X\alpha'$ with $\alpha'=m\alpha^{-1}\in M_2(\Z)$ (the adjugate of $\alpha$), an integral matrix of determinant $1$ congruent to $1$
modulo $N$, hence in $\Gamma(N)\subseteq\Gamma$; so $\gamma\in\alpha^{-1}\Gamma\alpha$, and $\Gamma(Nm)\subseteq\Gamma_\alpha$. As $\Gamma(Nm)$ has finite index in $\SL_2(\Z)$, so has $\Gamma_\alpha$.

Why correspondences? On the upper half plane the Hecke operator $T_p$ of \Cref{ch:hecke} is a sum over the $p+1$ matrices $\sm1j0p$, $\sm p001$: it sends a point to a formal sum of
$p+1$ points, and on moduli it sends an elliptic curve to the sum of its $p+1$ quotients by subgroups of order $p$. Such a ``multivalued map'' of a curve to itself is a
correspondence: a third curve $C$ with two maps to $X_0(N)$, one used to pull points back (giving the several values) and the other to push them forward. Here $C=X(\Gamma_\alpha)$:
a point of $C$ is a point $\Gamma z$ of $X_0(N)$ together with a choice of one of its values ($p+1$ of them for $T_p$, $p$ for $U_p$), and $\pi_1$, $\pi_\alpha$ forget the choice and
take the value.

\begin{definition}[Hecke correspondences]\label{def:moddeg:hecke-correspondence}
For $\alpha$ as above, the \emph{Hecke correspondence} of $\alpha$ is the correspondence $\mathcal C_\alpha=(X(\Gamma_\alpha),\pi_1,\pi_\alpha)$ from $X_0(N)$ to itself in the sense
of \Cref{ch:curves}, where $X(\Gamma_\alpha)$ is the compactified quotient $\Gamma_\alpha\backslash\HH^*$ and
\[
\pi_1(\Gamma_\alpha z)=\Gamma z,\qquad\pi_\alpha(\Gamma_\alpha z)=\Gamma\alpha z.
\]
($\pi_\alpha$ is well defined: for $\gamma\in\Gamma_\alpha$, $\alpha\gamma z=(\alpha\gamma\alpha^{-1})\alpha z$ with $\alpha\gamma\alpha^{-1}\in\Gamma$.) For a prime $p$ put
$T_p=\mathcal C_{\mathrm{diag}(1,p)}$ if $p\nmid N$ and $U_p=\mathcal C_{\mathrm{diag}(1,p)}$ if $p\mid N$. For $\beta\in\GL_2^+(\Q)$ normalising $\Gamma$, $w_\beta$ is the automorphism
$\Gamma z\mapsto\Gamma\beta z$ of $X_0(N)$; for $\beta=W_N$ it is the involution $w_N$ ($W_N^2=-N$ acts trivially on $\HH$).
\end{definition}

By \Cref{ch:curves}, $\mathcal C_\alpha$ acts on divisors and on $J_0(N)$ by $(\mathcal C_\alpha)_*=\pi_{\alpha*}\pi_1^*$, on $H_1$ by $\pi_{\alpha*}\pi_1^!$, on $H^1$ by
$\mathcal C_\alpha^*=\pi_{1*}\pi_\alpha^*$, and on holomorphic differentials by $\omega\mapsto\Tr_{\pi_1}(\pi_\alpha^*\omega)$.

\begin{proposition}[Hecke correspondences]\label{prop:moddeg:hecke-correspondence}
Let $\alpha\in M_2(\Z)$, $\det\alpha>0$.
\begin{enumerate}[label=(\roman*)]
\item If $\gamma_1,\dots,\gamma_r$ represent $\Gamma_\alpha\backslash\Gamma$, then $\Gamma\alpha\Gamma=\bigsqcup_i\Gamma\alpha\gamma_i$, and $(\mathcal C_\alpha)_*(\Gamma z)=\sum_i(\Gamma\alpha\gamma_iz)$ for
$z\in\HH$ (with multiplicities at the ramification points of $\pi_1$). If $\beta=\gamma\alpha\gamma'$ with $\gamma,\gamma'\in\Gamma$, then $\mathcal C_\beta\simeq\mathcal C_\alpha$: the map
$\Gamma_\alpha z\mapsto\Gamma_\beta\gamma'^{-1}z$ is an isomorphism $X(\Gamma_\alpha)\to X(\Gamma_\beta)$ compatible with the two pairs of maps. So $\mathcal C_\alpha$ depends only on
$\Gamma\alpha\Gamma$.
\item Let $p$ be prime and $\alpha=\mathrm{diag}(1,p)$. Then $\Gamma_\alpha=\Gamma_0(N)\cap\Gamma^0(p)$, $\alpha\Gamma_\alpha\alpha^{-1}=\Gamma_0(Np)$, and $\Gamma_\alpha z\mapsto\Gamma_0(Np)(z/p)$
is an isomorphism $X(\Gamma_\alpha)\simeq X_0(Np)$ under which $\pi_1$ becomes $z\mapsto pz$ and $\pi_\alpha$ becomes $z\mapsto z$. The degree of $\pi_1$ is $p+1$ if $p\nmid N$ and
$p$ if $p\mid N$, with coset representatives giving
\[
\Gamma\alpha\Gamma=\bigsqcup_{j=0}^{p-1}\Gamma\Mtwo1j0p\ \sqcup\ \Gamma\Mtwo p001\ (p\nmid N),\qquad\Gamma\alpha\Gamma=\bigsqcup_{j=0}^{p-1}\Gamma\Mtwo1j0p\ (p\mid N).
\]
\item (Moduli; sketched.) Under \Cref{thm:moddeg:x0-over-q}(i), $(T_p)_*(E,C)=\sum_D(E/D,(C+D)/D)$ over the $p+1$ subgroups $D\subset E$ of order $p$ ($p\nmid N$), and
$(U_p)_*(E,C)=\sum_{D\cap C=0}(E/D,(C+D)/D)$ over the $p$ subgroups $D$ of order $p$ meeting $C$ trivially ($p\mid N$).
\item $T_p$, $U_p$ and the $w_d$ ($d\,\|\,N$) are correspondences defined over $\Q$.
\end{enumerate}
\end{proposition}

\begin{proof}
(i) $\Gamma\alpha\gamma=\Gamma\alpha\gamma'$ iff $\alpha\gamma'\gamma^{-1}\alpha^{-1}\in\Gamma$ iff $\gamma'\gamma^{-1}\in\Gamma\cap\alpha^{-1}\Gamma\alpha=\Gamma_\alpha$; so the cosets $\Gamma\alpha\gamma_i$ are distinct and exhaust $\Gamma\alpha\Gamma$. The fibre
of $\pi_1$ over $\Gamma z$ is $\{\Gamma_\alpha\gamma_iz\}$ (with multiplicities), and $\pi_\alpha(\Gamma_\alpha\gamma_iz)=\Gamma\alpha\gamma_iz$; this is the formula for $\pi_{\alpha*}\pi_1^*$. For $\beta=\gamma\alpha\gamma'$:
$\beta^{-1}\Gamma\beta=\gamma'^{-1}\alpha^{-1}\Gamma\alpha\gamma'$, so $\Gamma_\beta=\gamma'^{-1}\Gamma_\alpha\gamma'$ and $\Gamma_\alpha z\mapsto\Gamma_\beta\gamma'^{-1}z$ is a well-defined isomorphism. It satisfies
$\pi_1(\Gamma_\beta\gamma'^{-1}z)=\Gamma z$ and $\pi_\beta(\Gamma_\beta\gamma'^{-1}z)=\Gamma\gamma\alpha\gamma'\gamma'^{-1}z=\Gamma\alpha z$.

(ii) For $\gamma=\sm abcd\in\Gamma$, $\alpha\gamma\alpha^{-1}=\sm a{b/p}{pc}d$, which lies in $\Gamma$ iff $p\mid b$; so $\Gamma_\alpha=\Gamma\cap\Gamma^0(p)$, and its conjugate
$\alpha\Gamma_\alpha\alpha^{-1}$ consists of the $\sm a{b'}{c'}d$ of determinant $1$ with $Np\mid c'$, which is $\Gamma_0(Np)$. The map $z\mapsto\alpha z=z/p$ therefore induces
$X(\Gamma_\alpha)\simeq X_0(Np)$, and $\Gamma z=\Gamma p(z/p)$, $\Gamma\alpha z=\Gamma(z/p)$ give the two maps. The index $[\Gamma_0(N):\Gamma_0(N)\cap\Gamma^0(p)]$ is $p+1$ for $p\nmid N$ and
$p$ for $p\mid N$. To see this, reduce modulo $p$: the condition $p\mid b$ only depends on $\gamma\bmod p$, so the index is that of the subgroup $\{b=0\}$ in the image $G$ of $\Gamma_0(N)$ in
$\SL_2(\Fp)$. If $p\nmid N$, then $G=\SL_2(\Fp)$ (already $\Gamma(N)$ surjects onto $\SL_2(\Fp)$, because $\SL_2(\Z)\to\SL_2(\Z/Np)=\SL_2(\Z/N)\times\SL_2(\Z/p)$ is onto), and $\{b=0\}$ is
the lower triangular subgroup, of order $p(p-1)$ in a group of order $p(p^2-1)$: index $p+1$. If $p\mid N$, then $c\equiv0$ and $G$ is the upper triangular subgroup of order $p(p-1)$, in which
$\{b=0\}$ is the diagonal subgroup of order $p-1$: index $p$. The displayed matrices lie in $\Gamma\alpha\Gamma$ ($\sm1j0p=\alpha\sm1j01$) and represent distinct cosets; see \Cref{ch:hecke} and \cite[Section 5.2]{DiamondShurman}. For
$p\nmid N$ the matrix $\sm p001$ lies in $\Gamma\alpha\Gamma$ by the computation in the proof of \Cref{thm:moddeg:transpose}(iv).

(iii) (Sketch.) The point $\Gamma z$ is $(\C/\Lambda_z,\langle\frac1N\rangle)$ with $\Lambda_z=\Z+\Z z$. The lattices $\Lambda_{(z+j)/p}$ ($0\le j<p$) and $\frac1p\Z+\Z z=\frac1p\Lambda_{pz}$ are the $p+1$
lattices containing $\Lambda_z$ with index $p$, that is, the subgroups $D$ of order $p$ of $E=\C/\Lambda_z$, and $\C/\Lambda'=E/D$; the image of $\frac1N$ generates $(C+D)/D$
(for $\Lambda'=\frac1p\Lambda_{pz}$ one rescales by $p$, and $\frac1N$, $\frac1p$ generate $\frac1{Np}\Z/\Z$ when $p\nmid N$). For $p\mid N$ the missing lattice $\frac1p\Z+\Z z$ is the one
whose subgroup $D=\langle\frac1p\rangle$ lies in $C=\langle\frac1N\rangle$. See \cite[Sections 5.2, 7.9]{DiamondShurman}.

(iv) $X_0(Np)$, $z\mapsto pz$ and $z\mapsto z$ are defined over $\Q$ by \Cref{thm:moddeg:x0-over-q} (on moduli, with $C'$ the cyclic subgroup of order $Np$ attached to a point of $X_0(Np)$,
$z\mapsto z$ is $(E,C')\mapsto(E,pC')$ and $z\mapsto pz$ is $(E,C')\mapsto(E/C'[p],C'/C'[p])$), and so are the $w_d$.
\end{proof}

\Cref{ch:curves} writes the Hecke correspondence for $p\nmid N$ as $(X_0(Np),z\mapsto z,z\mapsto pz)$. Since $\Gamma_{\mathrm{diag}(p,1)}=\Gamma_0(Np)$ (the computation in (ii)), that is
$\mathcal C_{\mathrm{diag}(p,1)}$, the transpose of $\mathcal C_{\mathrm{diag}(1,p)}$ (\Cref{thm:moddeg:transpose}(i)); for $p\nmid N$ the two are isomorphic (\Cref{thm:moddeg:transpose}(iv)), but for $p\mid N$ the correspondence $(X_0(Np),z\mapsto z,z\mapsto pz)$ is
$U_p^t$, not $U_p$, and $U_p$ itself is $(X_0(Np),z\mapsto pz,z\mapsto z)$.

\begin{proposition}[compatibility with the Hecke operators on forms]\label{prop:moddeg:hecke-compatibility}
\leavevmode
\begin{enumerate}[label=(\roman*)]
\item For $f\in S_2(\Gamma_0(N))$ and $\Gamma\alpha\Gamma=\bigsqcup_i\Gamma\alpha_i$: $\mathcal C_\alpha^*(f(z)\,dz)=\bigl(\sum_if|_2\alpha_i\bigr)(z)\,dz$. For $\alpha=\mathrm{diag}(1,p)$ this is $(T_pf)(z)\,dz$
($p\nmid N$), respectively $(U_pf)(z)\,dz$ ($p\mid N$), with \Cref{ch:hecke}'s operators.
\item There is a unique ring homomorphism $\T\to\End_\Z H^1(X_0(N),\Z)$, $t\mapsto t^*$, with $T_p\mapsto\mathcal C_{\mathrm{diag}(1,p)}^*$ for every prime $p$. It is injective, and its
$\C$-linear extension to $H^1(X_0(N),\C)=S_2\oplus\overline{S_2}$ is $t$ on $S_2$ and $\bar\omega\mapsto\overline{t\omega}$ on $\overline{S_2}$. The transposes $t_*$ on
$H_1(X_0(N),\Z)=H^1(X_0(N),\Z)^\vee$ define a second injective ring homomorphism, which is the action of the correspondences by pushforward. In particular $\T$ is a finite
free $\Z$-module, and $\T_\ell=\T\otimes\Zl$ acts faithfully on $H^1(X_0(N),\Zl)$ and on $H_1(X_0(N),\Zl)$.
\item Since $\T\otimes\C\simeq\Hom_\C(S_2,\C)$ by $(t,g)\mapsto a_1(tg)$ (\Cref{ch:hecke}), every $\psi\in\Hom_\Z(\T,\Z)$ is $t\mapsto a_1(tg_\psi)$ for a unique $g_\psi\in S_2$ with
$a_n(g_\psi)=\psi(T_n)\in\Z$; so $S_2(\Gamma_0(N);\Z)$ has rank $g$. (This is the route from the integral structure of $H^1$ to Shimura's integral basis; \Cref{ch:modforms,ch:hecke}
prove the statement differently.)
\end{enumerate}
\end{proposition}

\begin{proof}
(i) By definition $\mathcal C_\alpha^*\omega=\Tr_{\pi_1}(\pi_\alpha^*\omega)$. On $\HH$, $\pi_\alpha$ is $z\mapsto\alpha z$, and $d(\alpha z)/dz=\det\alpha/(cz+d)^2$, so $\pi_\alpha^*(f(z)\,dz)=(f|_2\alpha)(z)\,dz$
(\Cref{conv:moddeg:setup}(C1)). The local sections of $\pi_1$ near a point $\Gamma z$ off the branch locus are $s_i\colon\Gamma z\mapsto\Gamma_\alpha\gamma_iz$, so
$\Tr_{\pi_1}(g\,dz)=\sum_is_i^*(g\,dz)=\sum_i(g|_2\gamma_i)\,dz$ for a weight-two form $g$ on $\Gamma_\alpha$. With $g=f|_2\alpha$ and $\alpha_i=\alpha\gamma_i$ (part (i) of
\Cref{prop:moddeg:hecke-correspondence}) the sum is $\sum_if|_2\alpha_i$. For $\alpha=\mathrm{diag}(1,p)$ and the representatives of
\Cref{prop:moddeg:hecke-correspondence}(ii): $f|_2\sm1j0p=p\cdot p^{-2}f(\frac{z+j}p)=\frac1pf(\frac{z+j}p)$ and $f|_2\sm p001=p\cdot1^{-2}f(pz)=pf(pz)$. With $f=\sum_na_nq^n$,
\[
\sum_{j=0}^{p-1}\frac1pf\Bigl(\frac{z+j}p\Bigr)=\sum_na_ne^{2\pi inz/p}\cdot\frac1p\sum_{j=0}^{p-1}e^{2\pi inj/p}=\sum_{p\mid n}a_nq^{n/p}=\sum_ma_{mp}q^m,
\]
because the inner sum is $1$ if $p\mid n$ and $0$ otherwise, and $pf(pz)=\sum_npa_nq^{pn}=\sum_mpa_{m/p}q^m$ (with $a_{m/p}=0$ if $p\nmid m$). Summing gives $\sum_m(a_{mp}+pa_{m/p})q^m$ (the last term for
$p\nmid N$ only), which is \Cref{ch:hecke}'s $T_pf$, respectively $U_pf$, in weight $2$, where the normalising factor $(\det\alpha)^{k/2-1}$ is $1$.

The class of a holomorphic $\omega$ is mapped by $\mathcal C^*=\pi_{1*}\pi_\alpha^*$ on $H^1$ to the class of $\Tr_{\pi_1}\pi_\alpha^*\omega$: \Cref{ch:curves} computes $\pi_{1*}$ on a representative
$\theta$ vanishing near the branch points; if $\theta=\omega+dF$, then $\Tr\theta-\Tr\omega=d(\Tr F)$ on the complement of the branch points, so the two forms have the same
integrals over loops avoiding them, and such loops span $H_1$.

(ii) Let $P$ be the polynomial ring $\Z[x_p:p\text{ prime}]$, let $\rho_1\colon P\to\End_\C S_2$ send $x_p$ to $T_p$ (or $U_p$), and $\rho_2\colon P\to\End_\Z H^1(X_0(N),\Z)$, $x_p\mapsto\mathcal C_{\mathrm{diag}(1,p)}^*$. The image
of $\rho_1$ is $\T$, because every $T_n$ is a polynomial in the $T_p$ and $U_p$ (\Cref{ch:hecke}). Each $\mathcal C^*$ on $H^1(X_0(N),\C)$ preserves $S_2$ (pullbacks and traces of
holomorphic forms are holomorphic, \Cref{ch:curves}), acts there as $\rho_1(x_p)$ by (i), and commutes with complex conjugation, being defined over $\Z$. So for
$F\in P$, $\rho_2(F)$ is $\rho_1(F)$ on $S_2$ and its conjugate on $\overline{S_2}$, and $\rho_2(F)=0$ iff $\rho_1(F)=0$ ($H^1(X_0(N),\Z)$ is a lattice in $H^1(X_0(N),\C)$, which is
$S_2\oplus\overline{S_2}$ by \Cref{ch:curves}). Hence $\ker\rho_1=\ker\rho_2$, and $t\mapsto t^*$ is the isomorphism $\T=P/\ker\rho_1\to\rho_2(P)$. Transposition reverses products,
and $\T$ is commutative, so $t\mapsto t_*$ is again a ring homomorphism; it is the pushforward action because $\mathcal C^*$ is the transpose of $\mathcal C_*$ (\Cref{ch:curves}). $\T$ is
a subring of the endomorphism ring of a lattice of rank $2g$, hence free of finite rank over $\Z$; $\T\otimes\Zl\to\End(H^1(X_0(N),\Z))\otimes\Zl$ is injective because
$\Zl$ is flat over $\Z$.

(iii) $\Hom_\Z(\T,\Z)$ is a lattice in $\Hom_\Z(\T,\C)=\Hom_\C(\T\otimes\C,\C)$, which the perfect pairing identifies with $S_2$; and $a_n(g_\psi)=a_1(T_ng_\psi)=\psi(T_n)$.
\end{proof}

Before the computation, the moduli picture shows what to expect. The transpose of a correspondence runs it backwards: $(T_p)^t$ sends $(E',C')$ to the sum of the pairs $(E,C)$ from
which $T_p$ reaches $(E',C')$, that is, it goes back along the dual isogenies $E'=E/D\to E$. For $p\nmid N$ the kernel $E[p]/D$ of the dual is just another subgroup of order $p$ of $E'$,
so going back is again of the form ``divide by a subgroup of order $p$'', and $T_p^t=T_p$. For $p\mid N$ the kernel of the dual of $E\to E/D$ with $D\cap C=0$ is $E[p]/D$, and since
$C[p]+D=E[p]$ this is the subgroup of order $p$ of the level structure $(C+D)/D$: $U_p$ divides by subgroups meeting $C$ trivially, while its transpose divides by the subgroup
contained in $C$. These are different operations, and $w_N$, which replaces $C$ by $E[N]/C$ on the quotient curve, exchanges ``contained in the level structure'' and ``meeting
it trivially''. The theorem makes this precise with matrices.

\begin{theorem}[the transpose of a Hecke correspondence]\label{thm:moddeg:transpose}
\leavevmode
\begin{enumerate}[label=(\roman*)]
\item For $\alpha\in M_2(\Z)$ with $\det\alpha>0$ put $\alpha'=(\det\alpha)\alpha^{-1}$. Then $\mathcal C_\alpha^t\simeq\mathcal C_{\alpha'}$: the map $u\colon X(\Gamma_\alpha)\to X(\Gamma_{\alpha'})$,
$\Gamma_\alpha z\mapsto\Gamma_{\alpha'}\alpha z$, is an isomorphism with $\pi_1\circ u=\pi_\alpha$ and $\pi_{\alpha'}\circ u=\pi_1$.
\item For $\beta\in\GL_2^+(\Q)$ normalising $\Gamma$, $\mathcal C_{\beta\alpha\beta^{-1}}\simeq w_\beta\,\mathcal C_\alpha\,w_\beta^{-1}$ (whenever $\beta\alpha\beta^{-1}\in M_2(\Z)$), so
$\mathcal C_{\beta\alpha\beta^{-1}}^*=(w_\beta^{-1})^*\mathcal C_\alpha^*w_\beta^*$ on $H^1$.
\item $W_N\,\mathrm{diag}(1,p)\,W_N^{-1}=\mathrm{diag}(p,1)=\mathrm{diag}(1,p)'$. Hence, for every prime $p$,
\[
T_p^t\simeq w_NT_pw_N\ (p\nmid N),\qquad U_p^t\simeq w_NU_pw_N\ (p\mid N),
\]
and for every $t\in\T$, $(t^*)^\dagger=w_N^*t^*w_N^*$ for the cup product on $H^1(X_0(N),\Z)$, and likewise for $t_*$ and the intersection pairing on $H_1(X_0(N),\Z)$.
\item For $p\nmid N$, $\Gamma\,\mathrm{diag}(1,p)\,\Gamma=\Gamma\,\mathrm{diag}(p,1)\,\Gamma$; so $T_p^t\simeq T_p$, $T_p^*$ is self-adjoint for the cup product and commutes with $w_N^*$. For $p\mid N$,
$U_p^*$ is in general not self-adjoint for the cup product (\Cref{ex:moddeg:ex-u2}).
\end{enumerate}
\end{theorem}

\begin{proof}
(i) Scalars act trivially on $\HH$ and by conjugation, so $\alpha'^{-1}\Gamma\alpha'=\alpha\Gamma\alpha^{-1}$ and $\Gamma_{\alpha'}=\Gamma\cap\alpha\Gamma\alpha^{-1}=\alpha\Gamma_\alpha\alpha^{-1}$. Hence
$\Gamma_\alpha z\mapsto\Gamma_{\alpha'}\alpha z$ is well defined and bijective, an isomorphism of the compact Riemann surfaces (it is induced by an automorphism of $\HH^*$). Then
$\pi_1(u(\Gamma_\alpha z))=\Gamma\alpha z=\pi_\alpha(\Gamma_\alpha z)$, and $\pi_{\alpha'}(u(\Gamma_\alpha z))=\Gamma\alpha'\alpha z=\Gamma(\det\alpha)z=\Gamma z=\pi_1(\Gamma_\alpha z)$. So $u$ is an isomorphism from
$\mathcal C_\alpha^t=(X(\Gamma_\alpha),\pi_\alpha,\pi_1)$ to $\mathcal C_{\alpha'}=(X(\Gamma_{\alpha'}),\pi_1,\pi_{\alpha'})$.

(ii) Put $\alpha_\beta=\beta\alpha\beta^{-1}$. As $\beta$ normalises $\Gamma$, $\Gamma_{\alpha_\beta}=\Gamma\cap\beta\alpha^{-1}\beta^{-1}\Gamma\beta\alpha\beta^{-1}=\beta\Gamma_\alpha\beta^{-1}$, and $v\colon\Gamma_\alpha z\mapsto\Gamma_{\alpha_\beta}\beta z$ is an
isomorphism with $\pi_1(v(\Gamma_\alpha z))=\Gamma\beta z=w_\beta(\pi_1(\Gamma_\alpha z))$ and $\pi_{\alpha_\beta}(v(\Gamma_\alpha z))=\Gamma\beta\alpha z=w_\beta(\pi_\alpha(\Gamma_\alpha z))$. In the notation of \Cref{ch:curves},
$(C,v^{-1}\circ\alpha,u\circ\beta)=u\,\mathcal C\,v$; here the triple is $(X(\Gamma_\alpha),w_\beta\pi_1,w_\beta\pi_\alpha)=w_\beta\,\mathcal C_\alpha\,w_\beta^{-1}$, which acts on $H^1$ by
$(w_\beta^{-1})^*\mathcal C_\alpha^*w_\beta^*$ (\Cref{ch:curves}).

(iii) $W_N\,\mathrm{diag}(1,p)=\sm0{-p}N0$ and $W_N^{-1}=\frac1N\sm01{-N}0$, so
\[
W_N\Mtwo100p W_N^{-1}=\frac1N\Mtwo0{-p}N0\Mtwo01{-N}0=\frac1N\Mtwo{pN}00N=\Mtwo p001,
\]
and $\mathrm{diag}(1,p)'=p\,\mathrm{diag}(1,p)^{-1}=\mathrm{diag}(p,1)$. By (i) and (ii), with $\alpha=\mathrm{diag}(1,p)$ and $\beta=W_N$ (which normalises $\Gamma_0(N)$, \Cref{ch:modforms}),
$\mathcal C_\alpha^t\simeq\mathcal C_{\alpha'}=\mathcal C_{W_N\alpha W_N^{-1}}\simeq w_N\mathcal C_\alpha w_N^{-1}$, and $w_N^{-1}=w_N$. Isomorphic correspondences act identically, so by
\Cref{ch:curves} (the adjoint of $\mathcal C^*$ for the cup product is $(\mathcal C^t)^*$) and (ii), $(T_p^*)^\dagger=(T_p^t)^*=w_N^*T_p^*w_N^*$, and the same for $U_p$. Let
$\mathcal A=\{t\in\T:(t^*)^\dagger=w_N^*t^*w_N^*\}$. It contains $1$ and every $T_p$, $U_p$, and it is closed under sums; it is closed under products because, with $w=w_N^*$,
$w^2=1$ and $\T$ commutative, $((st)^*)^\dagger=(s^*t^*)^\dagger=(t^*)^\dagger(s^*)^\dagger=wt^*ws^*w=w(ts)^*w=w(st)^*w$. As $\T$ is generated as a ring by the $T_p$ and $U_p$,
$\mathcal A=\T$. Poincare duality is an isometry carrying $\mathcal C^*$ to $(\mathcal C^t)_*$ (\Cref{ch:curves}), which gives the statement on $H_1$.

(iv) Choose integers $b$, $d$ with $pd-bN=1$ and put $\gamma_1=\sm pbNd$, $\gamma_2=\sm1bN{dp}$, both in $\Gamma_0(N)$. Then
$\gamma_1\,\mathrm{diag}(1,p)=\sm p{bp}N{dp}=\mathrm{diag}(p,1)\,\gamma_2$, so $\mathrm{diag}(p,1)=\gamma_1\mathrm{diag}(1,p)\gamma_2^{-1}$ and the double cosets agree. By
\Cref{prop:moddeg:hecke-correspondence}(i) and (i) above, $T_p^t\simeq\mathcal C_{\mathrm{diag}(p,1)}\simeq T_p$, so $(T_p^*)^\dagger=T_p^*$; with (iii), $T_p^*=w_N^*T_p^*w_N^*$.
\end{proof}

Why a twisted pairing is wanted. The congruence modules of \Cref{ch:modules} are built from a perfect pairing on a lattice for which the whole Hecke algebra is self-adjoint: then
the idempotents of $\T_\ell$ are self-adjoint, the local components are orthogonal, and the orthogonal complement of the $f$-part is the part killed by the annihilator of $I_f$. The cup
product is perfect but, by \Cref{thm:moddeg:transpose}, the $U_p^*$ are not self-adjoint for it. Composing with $w_N^*$ repairs this, at no cost: the new pairing is still perfect,
and it stays alternating because $w_N^*$ is a self-adjoint isometry.

\begin{corollary}[the pairing $\hat\delta_!$]\label{cor:moddeg:delta-hat}
Let $R$ be $\Z$, $\Zl$ or $\Ql$ and $\langle\ ,\ \rangle$ the cup product on $H^1(X_0(N),R)$ (\Cref{ch:curves}).
\begin{enumerate}[label=(\roman*)]
\item $(t^*)^\dagger=w_N^*t^*w_N^*$ for every $t\in\T\otimes R$.
\item The pairing
\[
\hat\delta_!(x,y)=\langle x,w_N^*y\rangle
\]
is perfect and alternating on $H^1(X_0(N),R)$, and every $t^*$, $t\in\T\otimes R$ (the $U_p^*$ included), is self-adjoint for it.
\item Let $R=\Zl$. Then every idempotent of $\T_\ell$ is self-adjoint for $\hat\delta_!$,
\[H^1(X_0(N),\Zl)=\bigoplus_{\mf m}H^1(X_0(N),\Zl)_{\mf m}\]
(over the maximal ideals $\mf m$ of $\T_\ell$) is an orthogonal decomposition, and $\hat\delta_!$ is perfect on each summand.
\end{enumerate}
\end{corollary}

\begin{proof}
(i) is \Cref{thm:moddeg:transpose}(iii), extended $R$-linearly. (ii) $w=w_N^*$ is an isometry of the cup product with $w^2=1$ (\Cref{ch:curves}: an automorphism of a compact Riemann
surface is holomorphic, hence orientation preserving of degree $1$, and $\langle w^*x,w^*y\rangle=\deg(w_N)\langle x,y\rangle$), so $w^\dagger=w^{-1}=w$. The cup
product is perfect, so $\hat\delta_!$, its composite with the automorphism $w$, is perfect: in matrix terms, if $G$ is the Gram matrix of the cup product in a basis and $W$ the matrix of
$w$, the Gram matrix of $\hat\delta_!$ is $GW$, and $\det(GW)=\det G\det W=\pm\det G$ is a unit. For the symmetry, use that the cup product on $H^1$ of a surface is antisymmetric and
that $w$ is a self-adjoint isometry: $\hat\delta_!(y,x)=\langle y,wx\rangle=-\langle wx,y\rangle=-\langle x,w^\dagger y\rangle=-\langle x,wy\rangle=-\hat\delta_!(x,y)$; so
$\hat\delta_!$ is antisymmetric, and alternating because $2\hat\delta_!(x,x)=0$ and $H^1(X_0(N),R)$ is torsion-free. For $t\in\T\otimes R$,
$\hat\delta_!(t^*x,y)=\langle x,(t^*)^\dagger wy\rangle=\langle x,wt^*wwy\rangle=\langle x,wt^*y\rangle=\hat\delta_!(x,t^*y)$. (iii) $\T_\ell$ is a finite free $\Zl$-algebra
(\Cref{prop:moddeg:hecke-compatibility}(ii)), so $\T_\ell=\prod_{\mf m}\T_{\mf m}$ with orthogonal idempotents $e_{\mf m}$ (\Cref{ch:modules}). An element of $\T_\ell$ is
self-adjoint by (ii); for an idempotent $e$, $\hat\delta_!(ex,(1-e)y)=\hat\delta_!(x,e(1-e)y)=0$, so the decomposition is orthogonal, and a perfect pairing restricted to the
summands of an orthogonal decomposition is perfect on each (\Cref{ch:modules}). This is \Cref{ch:curves}'s corollary on twisted pairings, whose hypothesis $t^t=w^*tw^*$ is
(i).
\end{proof}

The smallest case shows the sign. At level $11$, $S_2(\Gamma_0(11))=\C f_{11}$ and $\varepsilon_{11}(f_{11})=-1$, so $w_{11}^*$ is $-1$ on $H^1(X_0(11),\Z)\simeq\Z^2$ (it is $-1$ on
$\omega_f$ and on $\overline{\omega_f}$, which span $H^1(X_0(11),\C)$), and $\hat\delta_!=-\langle\ ,\ \rangle$: in a symplectic basis its Gram matrix is $\sm0{-1}10$, perfect and
alternating, as it must be. The twist only becomes visible when $\T$ contains a $U_p$ that is not self-adjoint for the cup product (\Cref{ex:moddeg:ex-u2}).

\begin{remark}[the Petersson side]\label{rem:moddeg:petersson-analogue}
On $S_2(\Gamma_0(N))$, $T_n$ is self-adjoint for the Petersson product when $(n,N)=1$, and $U_p^\dagger=w_NU_pw_N$ for $p\mid N$ (\Cref{ch:hecke}); a newform is a $w_N$-eigenvector
with eigenvalue $\varepsilon_N(f)=\pm1$. The two statements agree: on $S_2\subset H^1(X_0(N),\C)$, $\langle f,g\rangle=\frac i2\langle\omega_f,\overline{\omega_g}\rangle_{\mathrm{cup}}$, and $\mathcal C^*$
commutes with complex conjugation, so $\frac i2\langle\mathcal C^*\omega_f,\overline{\omega_g}\rangle=\frac i2\langle\omega_f,\overline{(\mathcal C^t)^*\omega_g}\rangle$: the Petersson adjoint of $\mathcal C^*$ on $S_2$ is
$(\mathcal C^t)^*$. This is why $f$ is an eigenvector of the adjoints of the $T_n$ (Lemma lem:crit), with the same eigenvalues.
\end{remark}

\begin{paperlink}\label{pl:moddeg:delta-hat}
In Section 4.1 of the paper the pairing $\hat\delta_!$ ``is perfect and alternating on $H^1(X_0(M^\Sigma),\Z_\ell)$, and the Hecke operators, $U_p$ included, are self-adjoint for it,
because $w_{M^\Sigma}T_nw_{M^\Sigma}^{-1}$ is the adjoint of $T_n$ for the cup product (this is what the composition with $w$ is for; see \cite[\S\S1.6.1, 1.8.2]{DFG04} and \cite[\S4.4]{DDT})''.
This is \Cref{thm:moddeg:transpose}(iii) and \Cref{cor:moddeg:delta-hat}(ii) at level $M^\Sigma$. The paper's $\hat\delta_!$ is that of \cite{DFG04}, which on classes from
$X_0(M^\Sigma)$ the paper describes as ``$\frac u2$ times the cup product on $X_0(M^\Sigma)$ composed with $w_{M^\Sigma}$, for an $\ell$-adic unit $u$, so for odd $\ell$ only a unit separates the two'';
the book's $\hat\delta_!$ (Appendix B, B-C13(f)) omits the factor $\frac u2$, a unit for odd $\ell$, so the two have the same ideal of values on every sublattice (\Cref{rem:moddeg:dfg-normalisation}(3)). The sentence of the paragraph after Lemma lem:hecke, ``the
idempotent $e_{\mathfrak m}$ is self-adjoint for $\hat\delta_!$, so $H^1(X_0(N),\Z_\ell)$ splits orthogonally into its local components and the pairing is perfect on each'', is
\Cref{cor:moddeg:delta-hat}(iii).
\end{paperlink}

So far everything has been said in terms of $H^1$ and $H_1$ of the Riemann surface. The modular parametrisation, however, is a map to an elliptic curve, and the
natural home of maps from a curve to elliptic curves is its Jacobian; the next section brings it in and sorts out which Hecke action lives where.

\section{The Jacobian \texorpdfstring{$J_0(N)$}{J0(N)}, its homology and the \texorpdfstring{$f$}{f}-lattice}\label{sec:moddeg:jacobian}

The modular parametrisation factors through the Jacobian $J_0(N)$, optimality is a statement about the kernel of $J_0(N)\to E$, Remark rem:free speaks of $J_0(N)[\mf m]$
and of the Tate module, and Section 4.1 of the paper defines the lattice $M_{f,\lambda}$ inside $H^1(X_0(N),\Zl)$. This section applies \Cref{ch:curves} to $X_0(N)$, keeps
track of which Hecke action lives on which object, and defines $M_{f,\lambda}$.

\begin{definition}[$J_0(N)$]\label{def:moddeg:j0}
$J_0(N)=\Jac(X_0(N)_\Q)$, the Jacobian of the curve $X_0(N)$ over $\Q$ (\Cref{ch:curves}); $\iota\colon X_0(N)\to J_0(N)$, $P\mapsto[(P)-(\infty)]$, is the Abel-Jacobi map with base
point $\infty$.
\end{definition}

\begin{theorem}[structure of $J_0(N)$]\label{thm:moddeg:j0}
\leavevmode
\begin{enumerate}[label=(\roman*)]
\item $J_0(N)$ is an abelian variety over $\Q$ of dimension $g=\dim S_2(\Gamma_0(N))$, with good reduction at every prime $p\nmid N$.
\item $J_0(N)(\C)=S_2(\Gamma_0(N))^\vee/H_1(X_0(N),\Z)$, where $\gamma\in H_1$ is sent to $f\mapsto\int_\gamma f(z)\,dz$.
\item $\iota$ is defined over $\Q$, and every morphism $X_0(N)\to A$ over $\Q$ to an abelian variety sending $\infty$ to $0$ is $h\circ\iota$ for a unique homomorphism
$h\colon J_0(N)\to A$ over $\Q$.
\item For $N=11$, $14$, $20$, $37$, $77$, $99$, $121$, $140$, $539$, $726$, $980$ the dimension of $J_0(N)$ is $1$, $1$, $1$, $2$, $7$, $9$, $6$, $19$, $49$, $109$, $145$.
\end{enumerate}
\end{theorem}

\begin{proof}
(i) The Jacobian of a smooth projective curve of genus $g$ over $\Q$ is an abelian variety of dimension $g$ over $\Q$ (\Cref{ch:curves}, stated), and $g=\dim S_2$ because
$f\mapsto f\,dz$ identifies $S_2(\Gamma_0(N))$ with the holomorphic differentials (\Cref{ch:modforms}). Good reduction of $J_0(N)$ at $p\nmid N$ follows from that of $X_0(N)$
(\Cref{thm:moddeg:x0-over-q}(iv)); we do not prove the implication (the Jacobian of a smooth proper curve over $\Zp$ is an abelian scheme; see \cite[Section 9.4]{BLR}, not verified).
(ii) is the analytic Jacobian and the Abel-Jacobi theorem of \Cref{ch:curves} with $H^0(X_0(N),\Omega^1)=S_2$. (iii) The base point $\infty$ is rational
(\Cref{prop:moddeg:cusps-over-q}), so the universal property of \Cref{ch:curves} applies over $\Q$. (iv) is \texttt{mfdim([N,2],1)} in PARI, and equals the genus formula of
\Cref{ch:modforms}.
\end{proof}

Two remarks on the statement. The choice of base point matters for (iii): a morphism $X_0(N)\to A$ factors through a homomorphism only after it has been normalised to send the base
point to $0$, and the base point must be rational for this to work over $\Q$; this is why modular parametrisations are required to send $\infty$ to $O$. And in genus one (ii) is
concrete: for $N=11$, $S_2(\Gamma_0(11))^\vee=\C$ by $\lambda\mapsto\lambda(f_{11})$, the lattice is $\{\int_\gamma f_{11}(z)\,dz:\gamma\in H_1\}$, and $\iota$ is an isomorphism
$X_0(11)\to J_0(11)$, the Abel-Jacobi map of a curve of genus one with a rational point.

\begin{proposition}[the Hecke action on $J_0(N)$ and $T_\ell J_0(N)$]\label{prop:moddeg:hecke-on-j0}
\leavevmode
\begin{enumerate}[label=(\roman*)]
\item $t\mapsto t_*$ is an injective ring homomorphism $\T\to\End_\Q(J_0(N))$; on $J_0(N)(\C)=S_2^\vee/H_1$ the endomorphism $t_*$ is induced by the transpose of $t$ on $S_2$, and on the
lattice $H_1$ it is $t_*$ of \Cref{prop:moddeg:hecke-compatibility}(ii).
\item On $T_\ell J_0(N)=H_1(X_0(N),\Zl)$ the action is $t_*$, and it commutes with $G_\Q$.
\item For $t\in\T$ put $t^t=w_Ntw_N\in\End_\C S_2$; on $H^1$ its action is $(t^t)^*=w_N^*t^*w_N^*=(t^*)^\dagger$ (\Cref{thm:moddeg:transpose}(iii)), and $t\mapsto t^t$ is a ring isomorphism from
$\T$ onto $w_N\T w_N$. In general $w_N\T w_N\ne\T$ (\Cref{ex:moddeg:ex-u2}). The \emph{Picard} action of $\T$ on $J_0(N)$ is $t\mapsto(t^t)_*$. A newform $f$ of level $N$ is an eigenvector of
every $t^t$ with eigenvalue $\phi_f(t)$. Hence $I_fJ_0(N)$, $J_0(N)/I_fJ_0(N)$ and $H^1(X_0(N),\Zl)[I_f]$ are the same whether $\T$ acts through $t_*$ (Albanese) or through $(t^t)_*$
(Picard).
\end{enumerate}
\end{proposition}

\begin{proof}
(i) A correspondence acts on $J_0(N)$ by $\mathcal C_*$, compatibly with Abel-Jacobi: by the transpose of $\mathcal C^*$ on differentials and by $\mathcal C_*$ on $H_1$
(\Cref{ch:curves}). On generators these are $\T$'s operators, so the ring homomorphism of \Cref{prop:moddeg:hecke-compatibility}(ii) gives an action on the complex torus,
injective because it is injective on $H_1$; it is defined over $\Q$ because the correspondences are (\Cref{prop:moddeg:hecke-correspondence}(iv)). (ii) $T_\ell J=H_1\otimes\Zl$
naturally for homomorphisms of tori (\Cref{ch:curves}); an endomorphism defined over $\Q$ commutes with $G_\Q$ on torsion points. (iii) The first assertions are \Cref{thm:moddeg:transpose}(iii) and the fact that conjugation by an involution is a ring isomorphism. If $f$ is a newform, $f|_2W_N=\varepsilon_N(f)f$
with $\varepsilon_N(f)=\pm1$ (\Cref{ch:hecke}), so $w_Ntw_Nf=\varepsilon_N(f)w_Ntf=\varepsilon_N(f)\phi_f(t)w_Nf=\phi_f(t)f$. The Picard version of $I_fJ_0(N)$ is
$\sum_{t\in I_f}w_{N*}t_*w_{N*}J_0(N)=w_{N*}(I_fJ_0(N))$. Its Lie algebra, in $S_2^\vee$, is the image under the transpose of $w_N^*$ of the Lie algebra of $I_fJ_0(N)$; the latter is
the annihilator of $S_2[I_f]=\C f$ (the Lie algebra of $t_*J_0(N)$ is the image of the transpose of $t$, whose annihilator is $\ker t$), and $w_N^*$ preserves $\C f$; a connected
closed subgroup of a complex torus is determined by its Lie algebra, so $w_{N*}(I_fJ_0(N))=I_fJ_0(N)$. On $H^1$, $x$ is killed by all $(t^t)^*-\phi_f(t)$, $t\in I_f$, iff $w_N^*x\in M_{f,\lambda}$, iff
$x\in M_{f,\lambda}$, because $w_N^*$ preserves $M_{f,\lambda}$ (\Cref{prop:moddeg:f-lattice}(iii)).
\end{proof}

\begin{proposition}[$T_\ell J_0(N)$ and $H^1(X_0(N),\Zl)$ as Hecke modules]\label{prop:moddeg:tate-module-vs-h1}
\leavevmode
\begin{enumerate}[label=(\roman*)]
\item $H^1(X_0(N),\Zl)=\Hom_{\Zl}(T_\ell J_0(N),\Zl)$, with $t^*$ the transpose of $t_*$. As $G_\Q$-modules (\'etale cohomology of the curve over $\Q$),
$H^1_{\mathrm{et}}(X_0(N)_{\Qbar},\Zl)\simeq(T_\ell J_0(N))^\vee\simeq T_\ell J_0(N)(-1)$, and the cup product is $G_\Q$-equivariant with values in $\Zl(-1)$.
\item The map $\Psi=w_{N*}\circ\operatorname{PD}\colon H^1(X_0(N),\Zl)\to H_1(X_0(N),\Zl)=T_\ell J_0(N)$ is an isomorphism of $\T_\ell$-modules for the actions $t^*$ on the source and
$t_*$ on the target. Poincare duality $\operatorname{PD}$ alone intertwines $t^*$ with $(t^t)_*$. After the Tate twist by $\Zl(-1)$, $\Psi$ is $G_\Q$-equivariant.
\item Hence, for every maximal ideal $\mf m$ of $\T_\ell$, the module $T_\ell J_0(N)_{\mf m}$ is free over $\T_{\mf m}$ if and only if $H^1(X_0(N),\Zl)_{\mf m}$ is, and
$H^1(X_0(N),\Fl)_{\mf m}[\mf m]\simeq J_0(N)[\mf m]$ as $\T/\mf m$-vector spaces.
\end{enumerate}
\end{proposition}

\begin{proof}
(i) is \Cref{ch:curves} (the Tate module and $H^1$; the Weil pairing and \'etale cohomology, stated there), with the comparison isomorphism compatible with correspondences.
(ii) By \Cref{ch:curves}, $\operatorname{PD}\circ\mathcal C^*=(\mathcal C^t)_*\circ\operatorname{PD}$, so $\operatorname{PD}\circ t^*=(t^t)_*\circ\operatorname{PD}=(w_Ntw_N)_*\circ\operatorname{PD}$ by
\Cref{thm:moddeg:transpose}(iii) and \Cref{prop:moddeg:hecke-on-j0}(iii). Composing with $w_{N*}$ and using $w_N^2=1$: $\Psi\circ t^*=w_{N*}(w_Ntw_N)_*\operatorname{PD}=t_*w_{N*}\operatorname{PD}=t_*\circ\Psi$.
$\Psi$ is an isomorphism because both factors are. Over $\Qbar$, Poincare duality is the Weil pairing $T_\ell J\simeq(T_\ell J)^\vee(1)$, which is $G_\Q$-equivariant, and $w_N$ is defined over
$\Q$; so $\Psi$ is $G_\Q$-equivariant as a map $H^1\to T_\ell J(-1)$. (iii) An isomorphism of $\T_\ell$-modules maps $\mf m$-components to $\mf m$-components. For the second
statement reduce $\Psi$ modulo $\ell$: $H^1(X_0(N),\Fl)\simeq H_1(X_0(N),\Fl)=J_0(N)[\ell]$ (\Cref{ch:curves}), $\T$-equivariantly, and take $\mf m$-torsion.
\end{proof}

\begin{warning}[four identifications that are easy to confuse]\label{warn:moddeg:four-identifications}
\leavevmode
\begin{enumerate}[label=(\arabic*)]
\item $T_\ell J_0(N)=H_1(X_0(N),\Zl)$: covariant, $\T$ acts by $t_*$.
\item $H^1(X_0(N),\Zl)=\Hom(T_\ell J_0(N),\Zl)$: contravariant, $\T$ acts by $t^*$; Galois-equivariantly it is $T_\ell J_0(N)(-1)$, not $T_\ell J_0(N)$.
\item $\operatorname{PD}\colon H^1\to H_1$ intertwines $t^*$ with $(t^t)_*=(w_Ntw_N)_*$, not with $t_*$.
\item $w_{N*}\circ\operatorname{PD}$ intertwines $t^*$ with $t_*$, and is Galois-equivariant only after the twist by $\Zl(-1)$.
\end{enumerate}
For an elliptic curve the difference between (1) and (2) is the difference between $T_\ell E$, of determinant $\chi_\ell$, and $H^1(E,\Zl)\simeq T_\ell E(-1)$, of determinant
$\chi_\ell^{-1}$. The adjoint does not see it: $\ad(T^\vee)\simeq\ad T\simeq\ad(T(-1))$ (\Cref{ch:reps}).
\end{warning}

\begin{definition}[the $f$-lattice]\label{def:moddeg:f-lattice}
Let $f$ be a newform of level $N$ with rational coefficients and $\ell$ a prime. Put $V_f=H^1(X_0(N),\Q)[I_f]$ and
\[
M_{f,\lambda}=H^1(X_0(N),\Zl)[I_f]=\{x:\ t^*x=\phi_f(t)x\ \text{for all }t\in\T\}.
\]
\end{definition}

The plane $V_f$ is the part of the cohomology that ``belongs to $f$''; over $\Q$ it is a $2$-dimensional space, and the definition picks out a specific $\Zl$-lattice in it. There are
two natural lattices in $V_f\otimes\Ql$: the intersection $M_{f,\lambda}$ with the integral cohomology, and the image of $H^1(X_0(N),\Zl)$ under the projection onto the $f$-component
$e_fH^1(X_0(N),\Ql)$ (with $e_f$ the idempotent of $\T\otimes\Ql$ of \Cref{prop:moddeg:congruence-ideal}(i)). The first is contained in the second, with index measured by congruences
between $f$ and the other forms; the paper and \cite{DFG04} use the first (\Cref{rem:moddeg:dfg-normalisation}(4)). At level $11$, where $\T=\Z$, the two agree and
$M_{f,\lambda}=H^1(X_0(11),\Zl)$.

\begin{proposition}[the $f$-lattice]\label{prop:moddeg:f-lattice}
In the situation of \Cref{def:moddeg:f-lattice}:
\begin{enumerate}[label=(\roman*)]
\item $V_f\otimes\C=H^1(X_0(N),\C)[I_f]=\C\,\omega_f\oplus\C\,\overline{\omega_f}$, of dimension $2$; $V_f\otimes\R$ is spanned by $\operatorname{Re}\omega_f$ and $\operatorname{Im}\omega_f$. The subspace killed by
the anemic algebra $\T^{(N)}=\Z[T_n:(n,N)=1]$ acting through $\phi_f$ is the same plane.
\item $M_{f,\lambda}=H^1(X_0(N),\Z)[I_f]\otimes\Zl=H^1(X_0(N),\Zl)\cap(V_f\otimes\Ql)$; it is a saturated $\Zl$-sublattice of rank $2$.
\item $w_N^*$ preserves $V_f$ and $M_{f,\lambda}$, and acts on both by $\varepsilon_N(f)$.
\item $M_{f,\lambda}$ is stable under $G_\Q$ (in \'etale cohomology).
\end{enumerate}
\end{proposition}

\begin{proof}
(i) Kernels of linear maps over $\Q$ commute with extension of scalars; so the complexification of $V_f$ is the space $H^1(X_0(N),\C)[I_f]$. By \Cref{prop:moddeg:hecke-compatibility}(ii) this is
$S_2[I_f]\oplus\overline{S_2[I_f]}$ (for $t\in I_f$ the class $t\bar\omega$, which is $\overline{t\omega}$, vanishes iff $t\omega$ does). The eigenspace of $\phi_f$ in $S_2$ is the line $\C f$: by the theory of
newforms (\Cref{ch:hecke}), $S_2$ is the direct sum of the spaces spanned by the $g(dz)$, $g$ a newform of level $M\mid N$ and $d\mid N/M$, the anemic algebra acts on each
by the eigenvalues of $g$, distinct newforms have distinct anemic eigenvalue systems (strong multiplicity one), and for $g=f$ of level $N$ only $d=1$ occurs. This argument
uses only $\T^{(N)}$, which gives the last sentence. (ii) The first equality is flatness of $\Zl$ over $\Z$; the second holds because $x\in H^1(X_0(N),\Zl)$ is killed by $I_f$ iff
its image in $H^1(X_0(N),\Ql)$ is. A kernel of endomorphisms of a torsion-free module is saturated (\Cref{ch:modules}). The rank is $\dim V_f=2$. (iii) $w_N^*$ commutes with every
$t^*$, $t\in\T^{(N)}$ (\Cref{thm:moddeg:transpose}(iv)), so it preserves the plane, which by (i) is the $\phi_f$-eigenspace of $\T^{(N)}$. On it $w_N^*\omega_f=(f|_2W_N)\,dz=\varepsilon_N(f)\omega_f$,
and $w_N^*\overline{\omega_f}=\overline{w_N^*\omega_f}=\varepsilon_N(f)\overline{\omega_f}$ since $\varepsilon_N(f)$ is real. $w_N^*$ preserves $H^1(X_0(N),\Zl)$, hence $M_{f,\lambda}$. (iv) $t^*$ commutes with
$G_\Q$, so the kernel of $I_f$ is $G_\Q$-stable.
\end{proof}

\begin{paperlink}\label{pl:moddeg:f-lattice}
Section 4.1 of the paper defines $M_{f^\Sigma,\lambda}$ as ``the kernel on $M(M^\Sigma,1)_{!,\lambda}$ of the kernel of the eigencharacter of $f^\Sigma$ on the full Hecke algebra, $U_p$ included''; with
$\Sigma=\emptyset$ and $H^1(X_0(N),\Zl)$ in place of the integral structure of Diamond, Flach and Guo this is \Cref{def:moddeg:f-lattice}. In the proof of
Proposition prop:pullback the paper notes that the plane $\C f\,dz\oplus\C\overline{f\,dz}$ ``is defined over $\Q$: $H^1(X_0(N),\Q)[I_f]$ is a two-dimensional $\Q$-form of it, and by flatness
$H^1(X_0(N),\Z_\ell)[I_f]=H^1(X_0(N),\Z)[I_f]\otimes\Z_\ell$ is a saturated lattice of rank two in $H^1(X_0(N),\Q_\ell)[I_f]$'': this is \Cref{prop:moddeg:f-lattice}(ii), with $V_f$ the $\Q$-form.
It also says ``By multiplicity one for newforms the plane is also the eigenspace of the $T_n$, $(n,N)=1$, for the eigenvalues $a_n(f)$, and $w_N$ commutes with these $T_n$ (not with the $U_p$,
$p\mid N$)'': this is \Cref{prop:moddeg:f-lattice}(i), (iii). Remark rem:free: ``The identification of $T_\ell J_0(N)$ with $H^1(X_0(N),\Z_\ell)$ as
Hecke modules uses Poincar\'e duality composed with $w_N$, which is not Galois equivariant; freeness transfers nevertheless'' is \Cref{prop:moddeg:tate-module-vs-h1}(ii), (iii):
$\Psi$ is Galois-equivariant after a Tate twist.
\end{paperlink}

The lattice $M_{f,\lambda}$ has been defined purely in terms of Hecke operators. To connect it with an elliptic curve one needs a curve whose $H^1$ is $V_f$: this is
Shimura's quotient, and the bridge between its Frobenius traces and the coefficients of $f$ is the congruence relation.

\section{Eichler-Shimura and modularity}\label{sec:moddeg:eichler-shimura}

The paper starts from ``the newform attached to'' $E$. This section explains the two directions of that attachment: from a newform $f$ with rational coefficients
Shimura's construction produces an elliptic curve $E_f$ as a quotient of $J_0(N)$, whose Frobenius traces are the $a_p(f)$ by the congruence relation of Eichler and
Shimura; and the modularity theorem says that every elliptic curve over $\Q$ is isogenous to one of these. The quotient $E_f$ is the optimal curve of
\Cref{sec:moddeg:optimal}, and the congruence relation is the first step in the construction of the Galois representation of $f$ (\Cref{ch:modgalrep}).

\begin{theorem}[Shimura's quotient]\label{thm:moddeg:shimura-quotient}
Let $f$ be a newform of level $N$ with coefficient field $K_f$ and $I_f=\ker(\T\to K_f,\ T_n\mapsto a_n(f))$.
\begin{enumerate}[label=(\roman*)]
\item $I_fJ_0(N)=\sum_{t\in I_f}t_*J_0(N)$ is an abelian subvariety of $J_0(N)$ defined over $\Q$, and $J_f=J_0(N)/I_fJ_0(N)$ is an abelian variety over $\Q$ of dimension
$[K_f:\Q]$ on which $\T$ acts through $\T/I_f\hookrightarrow K_f$. The quotient map $\pi_f\colon J_0(N)\to J_f$ has connected kernel. (The book writes $J_f$ for Shimura's
$A_f$, which is the paper's adjoint representation; Appendix B, B-C12(c).)
\item If $f$ has rational coefficients, $J_f=:E_f$ is an elliptic curve over $\Q$, $t\in\T$ acts on it as multiplication by $\phi_f(t)$, and $\varphi_{E_f}=\pi_f\circ\iota$ satisfies
$\varphi_{E_f}^*H^1(E_f,\Q)=V_f$ and $\varphi_{E_f}^*\omega_{E_f}\in\C\,\omega_f$.
\item (Shimura; stated.) $L(J_f,s)=\prod_{\sigma\colon K_f\to\C}L(f^\sigma,s)$; in particular $L(E_f,s)=L(f,s)$.
\end{enumerate}
\end{theorem}

\begin{proof}[Proof of (ii), and sketch of (i)]
(i) Let $t_1,\dots,t_r$ generate $I_f$ as an abelian group. Then $I_fJ_0(N)$, which is $\sum_it_{i*}J_0(N)$, is the image of $J_0(N)^r$ under the homomorphism $(x_i)\mapsto\sum_it_{i*}x_i$,
defined over $\Q$; the image of an abelian variety under a homomorphism is an abelian subvariety (connected, as the image of a connected group), and the quotient of an
abelian variety by an abelian subvariety exists over the field of definition (stated; see the references on abelian varieties in \Cref{ch:curves}). $\T$ is commutative, so each $t_*$ preserves
$I_fJ_0(N)$ and acts on $J_f$; and $t-\phi_f(t)$, which lies in $I_f$ when $f$ is rational, maps $J_0(N)$ into $I_fJ_0(N)$, so $t$ acts on $J_f$ as $\phi_f(t)$ in that case.
The kernel $I_fJ_0(N)$ is connected by construction. Over $\C$, the Lie algebra of $t_*J_0(N)$ is the image of the transpose of $t$ on $S_2^\vee$, whose annihilator in $S_2$ is
$\ker t$; so the Lie algebra of $I_fJ_0(N)$ is the annihilator of $S_2[I_f]$, which is spanned by the conjugates $f^\sigma$ and has dimension $[K_f:\Q]$ (\Cref{ch:hecke}). Hence
$\dim J_f=[K_f:\Q]$.
(ii) For rational $f$, $S_2[I_f]=\C f$ and $\dim J_f=1$: $J_f$ is a connected projective group variety of dimension one over $\Q$, an elliptic curve with origin the image of
$0$. Let $x\in H^1(E_f,\Q)$. As $t$ acts on $E_f$ as $\phi_f(t)$, and multiplication by an integer $m$ acts on $H^1$ of an elliptic curve as $m$ (\Cref{ch:curves}),
$t^*\varphi_{E_f}^*x=\varphi_{E_f}^*(\phi_f(t)x)$ (pullback along $\iota$ identifies $t_*$ on $J_0(N)$ with $t^*$ on $H^1(X_0(N))$, \Cref{prop:moddeg:hecke-on-j0}). So
$\varphi_{E_f}^*H^1(E_f,\Q)\subseteq V_f$; both have dimension $2$ ($\varphi_{E_f}^*$ is injective by \Cref{ch:curves}), so they are equal. Then $\varphi_{E_f}^*\omega_{E_f}$ is a
holomorphic form in $V_f\otimes\C=\C\omega_f\oplus\C\overline{\omega_f}$ (\Cref{prop:moddeg:f-lattice}), hence in $\C\omega_f$.
(iii) We do not prove this; see \cite[Theorem 7.14 and Section 7.5]{Shimura71} (theorem number not verified), \cite{Shimura73}, \cite[Section 6.6]{DiamondShurman}. At the primes
$p\nmid N$ it follows from \Cref{thm:moddeg:eichler-shimura}.
\end{proof}

\begin{theorem}[the congruence relation of Eichler and Shimura]\label{thm:moddeg:eichler-shimura}
Let $p\nmid N$ be prime.
\begin{enumerate}[label=(\roman*)]
\item (Stated.) In $\End\bigl(J_0(N)_{\Fp}\bigr)$, the reduction of $(T_p)_*$ is $\pi_p+\hat\pi_p$, where $\pi_p$ is the $p$-power Frobenius endomorphism and $\hat\pi_p=p\pi_p^{-1}$
the Verschiebung.
\item On $T_\ell J_0(N)$, $\ell\ne p$, the arithmetic Frobenius $\Frob_p$ (which acts as $\pi_p$, \Cref{ch:curves}) satisfies $\Frob_p^2-(T_p)_*\Frob_p+p=0$. On each generalised
eigenspace $W_\psi$ of $\T$ in $T_\ell J_0(N)\otimes\Qbar_\ell$ (for a character $\psi\colon\T\to\Qbar_\ell$) the eigenvalues of $\Frob_p$ are roots of $X^2-\psi(T_p)X+p$.
\item If $f$ has rational coefficients, $E_f$ has good reduction at $p$, $\det(X-\Frob_p\mid T_\ell E_f)=X^2-a_p(f)X+p$, and $a_p(E_f)=p+1-\#\tilde E_f(\Fp)=a_p(f)$; on $H^1(E_f,\Zl)$
the geometric Frobenius has the same characteristic polynomial.
\end{enumerate}
\end{theorem}

\begin{proof}
(i) We do not prove this; see \cite[Section 7.4]{Shimura71}, \cite[Section 8.7]{DiamondShurman}, \cite[Section 10]{DiamondIm} (numbering not verified). The idea is
Kronecker's congruence $\Phi_p(X,Y)\equiv(X^p-Y)(X-Y^p)\pmod p$ read on moduli through \Cref{prop:moddeg:hecke-correspondence}(iii): over $\bar\F_p$, for an ordinary curve $E$
the quotient of $E$ by its connected subgroup of order $p$ is the Frobenius twist $E^{[p]}$, and the quotients by the $p$ \'etale subgroups are the curves whose Frobenius twist is
$E$; so modulo $p$ the correspondence $T_p$ is the graph of the Frobenius morphism plus its transpose, which act on the Jacobian as $\pi_p$ and $\hat\pi_p$.
(ii) As $J_0(N)$ has good reduction at $p$ and $\ell\ne p$, reduction modulo $p$ is injective on $\ell$-power torsion (the kernel of reduction is a pro-$p$ group, as for elliptic
curves in \Cref{ch:reduction}), so it identifies $T_\ell J_0(N)$ with $T_\ell$ of the reduction, compatibly with all endomorphisms defined over $\Q$; under this identification the
arithmetic Frobenius becomes $\pi_p$. Since $\pi_p\hat\pi_p=p$, (i) gives $\pi_p^2-(T_p)_*\pi_p+p=\pi_p^2-(\pi_p+\hat\pi_p)\pi_p+\hat\pi_p\pi_p=0$.
On $W_\psi$, $(T_p)_*-\psi(T_p)$ is nilpotent and commutes with $\Frob_p$; so $\Frob_p^2-\psi(T_p)\Frob_p+p$ is nilpotent there, and every eigenvalue of $\Frob_p$ on $W_\psi$ is a
root of $X^2-\psi(T_p)X+p$.
(iii) $E_f$ is a quotient of $J_0(N)$, which has good reduction at $p$, so $E_f$ has good reduction at $p$ (N\'eron-Ogg-Shafarevich, \Cref{ch:galrep}: $T_\ell E_f$ is a quotient of
the unramified $T_\ell J_0(N)\otimes\Ql$). The homomorphism $\pi_f$ reduces to a surjection $J_0(N)_{\Fp}\to(E_f)_{\Fp}$, which commutes with Frobenius and Verschiebung (both are
natural for homomorphisms), and $T_p$ acts on $E_f$ as $a_p(f)$. Hence $a_p(f)\pi_f=\pi_f\circ(T_p)_*=(\pi_p+\hat\pi_p)\circ\pi_f$ after reduction, and as $\pi_f$ is surjective,
$\pi_p+\hat\pi_p=a_p(f)$ in $\End(\tilde E_f)$. By \Cref{ch:ellcurves}, $\pi_p+\hat\pi_p=[p+1-\#\tilde E_f(\Fp)]$ and $\det(X-\Frob_p\mid T_\ell E_f)=X^2-a_p(E_f)X+p$; $\End(\tilde E_f)$ is
torsion-free, so $a_p(E_f)=a_p(f)$. The statement on $H^1$ is \Cref{ch:curves} (B-C3).
\end{proof}

\begin{theorem}[modularity; stated]\label{thm:moddeg:modularity}
For every elliptic curve $E/\Q$ of conductor $N$ there is a newform $f_E\in S_2(\Gamma_0(N))$ with rational coefficients such that $a_p(f_E)=a_p(E)$ for every prime $p$ (with
$a_p(E)\in\{0,\pm1\}$ the usual value at $p\mid N$). Equivalently, $E$ is isogenous over $\Q$ to $E_{f_E}$; and $L(E,s)=L(f_E,s)$.
\end{theorem}

We do not prove this. It is the theorem of Wiles \cite{Wiles95} and Taylor and Wiles \cite{TaylorWiles} for semistable $E$, completed by Breuil, Conrad, Diamond and Taylor
\cite{BCDT}; that the level of $f_E$ is the conductor is Carayol's theorem \cite{Carayol86}. The conjecture goes back to Taniyama, Shimura and Weil; Weil showed that the
functional equations of all twists characterise modular forms. The proof establishes an isomorphism $R=\T$ between a deformation ring of $\bar\rho_{E,\ell}$ and a localised Hecke
algebra (\Cref{ch:congruence}), and it is in the course of that proof that the congruence ideals of \Cref{sec:moddeg:degree-eta} first appear. The equivalence of the two forms:
$E_{f_E}$ has the same $a_p$ as $E$ for $p\nmid N$ (\Cref{thm:moddeg:eichler-shimura}), so $T_\ell E\otimes\Ql\simeq T_\ell E_{f_E}\otimes\Ql$ (Chebotarev and Brauer-Nesbitt, \Cref{ch:galois,ch:reps}; both are
semisimple by Faltings, \Cref{ch:galrep}), and Faltings's isogeny theorem gives an isogeny; conversely isogenous curves have the same $a_p$ and the same conductor.

\begin{corollary}[isogeny classes and newforms]\label{cor:moddeg:classes-and-newforms}
$E\mapsto f_E$ is a bijection from the $\Q$-isogeny classes of elliptic curves of conductor $N$ to the newforms of level $N$ with rational coefficients.
\end{corollary}

\begin{proof}[Sketch of proof]
Well defined: isogenous curves have the same $a_p$ (\Cref{ch:ellcurves}). Injective: curves with the same $a_p$ for almost all $p$ are isogenous (Faltings, as above). Surjective:
for $f$ rational of level $N$, $E_f$ has conductor $N$ (Carayol \cite{Carayol86}; stated) and $a_p(E_f)=a_p(f)$ for $p\nmid N$ (\Cref{thm:moddeg:eichler-shimura}), so $f_{E_f}$ and $f$ are
newforms of level $N$ with the same $a_p$ for almost all $p$, hence equal by strong multiplicity one (\Cref{ch:hecke}).
\end{proof}

\begin{example}[levels $11$ and $99$]\label{ex:moddeg:class-11a}
(a) The class \curve{11a} corresponds to $f_{11}=\eta(z)^2\eta(11z)^2$. It has three curves, \curve{11a1}, \curve{11a2}, \curve{11a3}, with isogeny degrees $5$ (\curve{11a1} to \curve{11a2}
and to \curve{11a3}) and $25$ (\curve{11a2} to \curve{11a3}) (\texttt{allisog} \cite{ecdata}). Since $X_0(11)=J_0(11)$ is \curve{11a1} (\Cref{ex:moddeg:small-levels}) and $\T=\Z$ at level $11$ (so
$I_f=0$ and $E_f=J_0(11)$), \Cref{thm:moddeg:eichler-shimura}(iii) proves $\#\curve{11a1}(\Fp)=p+1-a_p(f_{11})$ for every $p\ne11$: the identity of \Cref{ch:modforms}, checked
there numerically, is a theorem. The modular parametrisations of \curve{11a2} and \curve{11a3} are the two $5$-isogenies out of \curve{11a1}, of degree $5$ each
(\Cref{thm:moddeg:optimal}(ii)).

(b) $\dim S_2(\Gamma_0(99))=9$. The newforms of level $99$ with rational coefficients are those of the four classes \curve{99a}, \curve{99b}, \curve{99c}, \curve{99d}, and there
are no others (\texttt{mfdim([99,2],0)} $=4$); level $33$ has one newform and level $11$ one. By the decomposition of \Cref{ch:hecke}, a newform $g$ of level $M\mid N$ contributes the
forms $g(dz)$, $d\mid N/M$, one for each divisor: the newforms of level $99$ contribute one each, $f_{33}$ contributes $f_{33}(z)$, $f_{33}(3z)$ ($99/33=3$ has two divisors), and $f_{11}$
contributes $f_{11}(z)$, $f_{11}(3z)$, $f_{11}(9z)$ ($99/11=9$ has three). So $9=4\cdot1+1\cdot2+1\cdot3$. (Levels $1$, $3$ and $9$ have no cusp forms of weight $2$.) On the
$3$-dimensional old space of $f_{11}$ every $T_p$ with $p\nmid99$ acts as the scalar $a_p(f_{11})$; this is why the anemic algebra alone cannot cut out the plane of an old form,
while for a newform of level $99$ it can (\Cref{prop:moddeg:f-lattice}(i)).
\end{example}

\begin{paperlink}\label{pl:moddeg:modularity}
Section 1 of the paper: ``Let $E/\Q$ be an elliptic curve of conductor $N$ and let $f=f_E\in S_2(\Gamma_0(N))$ be the newform attached to it'' is
\Cref{thm:moddeg:modularity}; ``the isogeny class of $E$'' corresponds to $f$ by \Cref{cor:moddeg:classes-and-newforms}. Lemma lem:old uses the $a_p$ of newforms as traces of
Frobenius, which is \Cref{thm:moddeg:eichler-shimura}(iii) for rational newforms and its extension to all newforms in \Cref{ch:modgalrep}.
\end{paperlink}

\section[Optimal curves and the modular degree]{Optimal curves, the modular parametrisation and the modular degree}\label{sec:moddeg:optimal}

This section defines the objects of the paper's first sentence: the modular parametrisation, the $X_0(N)$-optimal curve and the modular degree. The main point is
\Cref{thm:moddeg:optimal}: every modular parametrisation of every curve of an isogeny class factors through Shimura's quotient $E_f$, so that ``minimal degree'',
``connected kernel'' and ``saturated pullback on $H^1$'' all single out the same curve. The section ends with Cremona's labelling convention and the table of modular
degrees used in the rest of the book.

\begin{definition}[modular parametrisation]\label{def:moddeg:parametrisation}
Let $E/\Q$ have conductor $N$. A \emph{modular parametrisation} of $E$ is a nonconstant morphism $\varphi\colon X_0(N)\to E$ over $\Q$ with $\varphi(\infty)=O$. By
\Cref{thm:moddeg:j0}(iii), modular parametrisations correspond bijectively to nonzero homomorphisms $h\colon J_0(N)\to E$ over $\Q$, by $\varphi=h\circ\iota$ and $h=\varphi_*$ (the
pushforward of divisor classes, \Cref{ch:curves}).
\end{definition}

\begin{theorem}[optimal curves]\label{thm:moddeg:optimal}
Let $\mathcal I$ be an isogeny class of elliptic curves over $\Q$ of conductor $N$, $f$ its newform (\Cref{cor:moddeg:classes-and-newforms}), $E_f=J_0(N)/I_fJ_0(N)$ and
$\varphi_{E_f}=\pi_f\circ\iota$ (\Cref{thm:moddeg:shimura-quotient}).
\begin{enumerate}[label=(\roman*)]
\item Every homomorphism $h\colon J_0(N)\to E$ over $\Q$ to a curve $E\in\mathcal I$ vanishes on $I_fJ_0(N)$.
\item (Factorisation.) $E_f\in\mathcal I$, and for every $E\in\mathcal I$ and every modular parametrisation $\varphi$ of $E$ there is a unique isogeny $\psi\colon E_f\to E$ over $\Q$ with
$\varphi=\psi\circ\varphi_{E_f}$. Hence $\deg\varphi=\deg\psi\cdot\deg\varphi_{E_f}$, and every isogeny $\psi\colon E_f\to E$ gives a modular parametrisation $\psi\circ\varphi_{E_f}$.
\item (Characterisations.) For $E\in\mathcal I$ and a modular parametrisation $\varphi$ of $E$ the following are equivalent: (a) $\psi$ in (ii) is an isomorphism; (b)
$\ker(\varphi_*\colon J_0(N)\to E)$ is connected; (c) $\varphi_*\colon H_1(X_0(N),\Z)\to H_1(E,\Z)$ is surjective; (d) $\varphi^*H^1(E,\Z)$ is saturated in $H^1(X_0(N),\Z)$; (e)
$\varphi^*\colon E=E^\vee\to J_0(N)$ is injective; (f) $\deg\varphi$ is the minimum of the degrees of all modular parametrisations of all curves in $\mathcal I$.
\item (Uniqueness.) Up to unique isomorphism, $(E_f,\varphi_{E_f})$ is the only pair $(E,\varphi)$ with $E\in\mathcal I$ satisfying these conditions. $E_f$ is the
\emph{$X_0(N)$-optimal curve} (or \emph{strong Weil curve}) of $\mathcal I$, $\varphi_{E_f}$ its \emph{modular parametrisation} $\varphi_E$, and $\degphi$ its \emph{modular degree}.
\end{enumerate}
\end{theorem}

\begin{proof}
(i) We may assume $h\ne0$; then $h$ is surjective and $h_*\colon V_\ell J_0(N)\to V_\ell E$ is a surjective $G_\Q$-equivariant map. Decompose
$V=V_\ell J_0(N)\otimes\Qbar_\ell=\bigoplus_\psi W_\psi$ into the generalised eigenspaces of $\T$ (\Cref{ch:modules}); each $W_\psi$ is $G_\Q$-stable, because $G_\Q$ commutes with $\T$. Suppose
$h_*(W_\psi)\ne0$. It is a $G_\Q$-stable subspace of $V_\ell E\otimes\Qbar_\ell$ and a quotient of $W_\psi$, so for every prime $p\nmid N\ell$ the arithmetic Frobenius has on it an eigenvalue
$\alpha$ which is a root of $X^2-a_p(E)X+p$ (\Cref{ch:ellcurves}) and of $X^2-\psi(T_p)X+p$ (\Cref{thm:moddeg:eichler-shimura}(ii)). Then $a_p(E)=\alpha+p/\alpha=\psi(T_p)$; and $a_p(E)=a_p(f)$
by \Cref{thm:moddeg:modularity}. So $\psi(T_p)=a_p(f)$ for every prime $p\nmid N\ell$. The characters of $\T$ that occur in $V$ take algebraic values and are, after
an identification of the algebraic numbers in $\Qbar_\ell$ and in $\C$, the eigencharacters of eigenforms in $S_2(\Gamma_0(N))$ (\Cref{prop:moddeg:hecke-compatibility}(ii), with $\Psi$ of
\Cref{prop:moddeg:tate-module-vs-h1}); by strong multiplicity one (\Cref{ch:hecke}), which allows any finite set of exceptional primes, an eigenform with $a_p=a_p(f)$ for all
$p\nmid N\ell$ is a multiple of $f$ (\Cref{prop:moddeg:f-lattice}(i)); so $\psi=\phi_f$. Hence $h_*$ kills
$\bigoplus_{\psi\ne\phi_f}W_\psi$. That sum is $I_fV$: $I_f$ kills $W_{\phi_f}$, and for $\psi\ne\phi_f$ some $t\in I_f$ has $\psi(t)\ne0$ (otherwise $\psi$ factors through $\T/I_f=\Z$ and equals
$\phi_f$), and $t$ is invertible on $W_\psi$. (That $I_f$ kills the generalised eigenspace $W_{\phi_f}$, and not only the eigenspace: the anemic algebra, generated by
the $T_n$ with $(n,N)=1$, acts semisimply on $S_2(\Gamma_0(N))$, its operators being normal for the Petersson product (\Cref{ch:hecke}), hence semisimply on $V$; so
$W_{\phi_f}$ lies in the eigenspace of the anemic algebra for the restriction of $\phi_f$, which corresponds to $\C\omega_f\oplus\C\overline{\omega_f}$
(\Cref{prop:moddeg:f-lattice}(i)) and on which every $t\in\T$ acts as $\phi_f(t)$.) Now $V_\ell(I_fJ_0(N))=I_fV_\ell J_0(N)$, as $I_fJ_0(N)$ is the image of $J_0(N)^r$ under $(t_i)$; so $h(I_fJ_0(N))$ is a connected
subgroup of $E$ with trivial rational Tate module, hence $0$.

(ii) By (i) and the universal property of the quotient, $\varphi_*=\psi\circ\pi_f$ for a unique homomorphism $\psi\colon E_f\to E$ over $\Q$ ($\pi_f$ is surjective); $\psi\ne0$, so $\psi$ is an
isogeny, and $\varphi=\varphi_*\circ\iota=\psi\circ\varphi_{E_f}$. Degrees of nonconstant maps of curves multiply (\Cref{ch:curves}). $E_f\in\mathcal I$: a curve $E\in\mathcal I$ is isogenous to
$E_{f_E}=E_f$ by \Cref{thm:moddeg:modularity}. Conversely $\psi\circ\varphi_{E_f}$ is nonconstant, defined over $\Q$, and sends $\infty$ to $\psi(O)=O$.

(iii) (b)$\Leftrightarrow$(c)$\Leftrightarrow$(d)$\Leftrightarrow$(e) is \Cref{ch:curves} (the proposition on the degree and the cup product, part (iii)), with $J(E)=E$. (a)$\Rightarrow$(b):
$\ker\varphi_*=\ker(\psi\circ\pi_f)=\ker\pi_f=I_fJ_0(N)$, connected. (b)$\Rightarrow$(a): $\ker\varphi_*=\pi_f^{-1}(\ker\psi)$ contains $I_fJ_0(N)$ with finite quotient $\ker\psi$; both have dimension
$g-1$, so if $\ker\varphi_*$ is connected it equals the connected closed subgroup $I_fJ_0(N)$ of the same dimension, and $\ker\psi=0$. (a)$\Leftrightarrow$(f): by (ii) the degrees of all
modular parametrisations of curves in $\mathcal I$ are the numbers $\deg\psi\cdot\deg\varphi_{E_f}$, with minimum $\deg\varphi_{E_f}$, attained exactly when $\deg\psi=1$.

(iv) If $(E,\varphi)$ satisfies (a), then $\psi\colon E_f\to E$ is the unique isomorphism with $\varphi=\psi\circ\varphi_{E_f}$.
\end{proof}

The logic of (i) in one sentence: a homomorphism $J_0(N)\to E$ sees only the part of $J_0(N)$ whose Frobenius eigenvalues are those of $E$, the congruence relation says that on the
$\psi$-part those eigenvalues are governed by $\psi(T_p)$, and strong multiplicity one says that only $\psi=\phi_f$ matches. The class \curve{11a} shows the characterisations at
work. There $J_0(11)=X_0(11)=\curve{11a1}$ and $I_f=0$, so $E_f=\curve{11a1}$ with $\varphi_{E_f}$ the identity. The parametrisation of \curve{11a3} is a $5$-isogeny
$\psi\colon\curve{11a1}\to\curve{11a3}$; the kernel of $\varphi_*=\psi$ is a group of order $5$, five points, not connected, so (b) fails; on $H_1$ the image of $\psi_*$ has index $5$ in
$H_1(\curve{11a3},\Z)$ (the index of $\psi_*H_1$ in $H_1$ is $\#\ker\psi$ for an isogeny of complex tori), so (c) fails; and its degree $5$ is not minimal, so (f) fails.

The paper's definition (Section 1: ``the $X_0(N)$-optimal curve, for which the modular parametrisation $\varphi_E\colon X_0(N)\to E$ has minimal degree'') is (f); the proof
of Proposition prop:pullback uses (b) and (c); and Section sec:data speaks of optimality as a property ``determined in the tables'', which is (a) checked by computation
(\Cref{conv:moddeg:cremona-optimal}).

\begin{proposition}[descriptions of the modular degree]\label{prop:moddeg:degree-descriptions}
Let $E$ be optimal, $d=\degphi$, $A=\varphi_E^*(E)\subseteq J_0(N)$ and $B=\ker\varphi_{E*}=I_fJ_0(N)$.
\begin{enumerate}[label=(\roman*)]
\item $d$ is the degree of the morphism of curves $\varphi_E$: the degree of the extension of function fields $\Q(X_0(N))\supseteq\varphi_E^*\Q(E)$, the number of preimages of a general point, the factor by
which $\varphi_E^*$ multiplies degrees of divisors, and the topological degree (\Cref{ch:curves}).
\item $\varphi_{E*}\circ\varphi_E^*=[d]$ on $E$: the composite $E=E^\vee\to J_0(N)\to E$ is multiplication by $d$.
\item $A\simeq E$, and $A\cap B=\varphi_E^*(E[d])\simeq(\Z/d)^2$; so $d$ is the exponent of the finite group $A\cap B$.
\item $\langle\varphi_E^*x,\varphi_E^*y\rangle_{X_0(N)}=d\,\langle x,y\rangle_E$ for $x,y\in H^1(E,\Z)$; the cup product on $\varphi_E^*H^1(E,\Z)$ has ideal of values $d\Z$ and Gram determinant $d^2$.
\end{enumerate}
\end{proposition}

\begin{proof}
(i) is \Cref{ch:curves}. (ii) is the functoriality of Jacobians (\Cref{ch:curves}): $\varphi_*\varphi^*=[\deg\varphi]$ on $J(E)=E$. (iii) $\varphi_E^*$ is injective by
\Cref{thm:moddeg:optimal}(iii)(e), so $A\simeq E$. For $y\in E$, $\varphi_E^*y\in B$ iff $\varphi_{E*}\varphi_E^*y=dy=0$; so $A\cap B=\varphi_E^*(E[d])\simeq E[d]\simeq(\Z/d)^2$, whose exponent is $d$.
(iv) The projection formula (\Cref{ch:curves}) says $\langle\varphi^*x,\varphi^*y\rangle_{X_0(N)}=\langle x,\varphi_*\varphi^*y\rangle_E$, and $\varphi_*\varphi^*$ is multiplication by
$\deg\varphi=d$ on $H^1(E,\Z)$ (dually, on $H_1$: the preimage of a loop on $E$ avoiding the branch points is a union of closed curves whose images cover the loop $d$ times). So the left side is $d\langle x,y\rangle_E$. The
cup product on $H^1(E,\Z)\simeq\Z^2$ is unimodular: in a symplectic basis $a,b$ its Gram matrix is $\sm01{-1}0$, so its values are all of $\Z$ and its Gram determinant is $1$. Multiplying by $d$
gives the Gram matrix $\sm0d{-d}0$ for the basis $\varphi^*a$, $\varphi^*b$: values $d\Z$ and determinant $d^2$.
\end{proof}

\begin{proposition}[the valuation of the degree along the class]\label{prop:moddeg:valuation-in-class}
Let $E$ be optimal in $\mathcal I$, $\ell$ a prime with $E[\ell]$ irreducible, and $E'\in\mathcal I$. Every isogeny $E\to E'$ of minimal degree is cyclic of degree prime to $\ell$, and the
minimal degree of a modular parametrisation of $E'$ has the same $\ell$-adic valuation as $\degphi$. Without irreducibility this fails: in the class \curve{11a} at $\ell=5$ the minimal
degrees are $1$, $5$, $5$.
\end{proposition}

\begin{proof}
If an isogeny $\psi\colon E\to E'$ has $E[m]\subseteq\ker\psi$ with $m>1$, it factors as $\psi'\circ[m]$ with $\deg\psi'=\deg\psi/m^2$ ($[m]$ is surjective with kernel $E[m]$, so
$\psi$ factors through it, and $\deg[m]=m^2$); $\psi'$ is defined over $\Q$ because $\psi$ and $[m]$ are. So an isogeny of minimal degree has a kernel containing no $E[m]$, $m>1$,
and a finite subgroup of $E\simeq(\R/\Z)^2$ with that property is cyclic (a finite abelian group with at most $2$ generators that is not cyclic contains some $(\Z/q)^2=E[q]$, $q$ prime). The
kernel is $G_\Q$-stable, $\psi$ being defined over $\Q$. If its order were divisible by $\ell$, the unique subgroup of order $\ell$ of this cyclic group would be $G_\Q$-stable too, contradicting the
irreducibility of $E[\ell]$ (\Cref{ch:ellcurves}). By
\Cref{thm:moddeg:optimal}(ii) the modular parametrisations of $E'$ are the $\psi\circ\varphi_E$, so the minimal degree is $\min_\psi\deg\psi\cdot\degphi$, and $v_\ell(\min\deg\psi)=0$. For
\curve{11a}: \curve{11a1} is optimal with $\degphi=1$ (\Cref{ex:moddeg:small-levels}), and the isogenies of minimal degree to \curve{11a2} and \curve{11a3} have degree $5$
(\Cref{ex:moddeg:class-11a}); $E[5]$ is reducible (\curve{11a1} has a rational point of order $5$).
\end{proof}

\begin{paperlink}\label{pl:moddeg:class-choice}
Section 1 of the paper: ``When $E[\ell]$ is irreducible, every curve of the class is linked to $E$ by a cyclic isogeny of degree prime to $\ell$, so $v_\ell$ of the degree of its
minimal parametrisation is $v_\ell(\degphi)$: at such an $\ell$ the choice of the optimal curve does not matter.'' This is \Cref{prop:moddeg:valuation-in-class}. The hypothesis
matters: the appendix of the paper notes that \texttt{rule} must be run on the optimal curve (``990h1 returns \texttt{[4, 3]} while the optimal curve 990h3 of its class returns
\texttt{[3, 3]}''), and the class \curve{990h} has $3$-isogenies.
\end{paperlink}

\begin{convention}[Cremona's labels and optimality]\label{conv:moddeg:cremona-optimal}
In Cremona's tables \cite{Cremona97,ecdata} curve number $1$ of each class is the optimal curve, with the single exception of the class \curve{990h}, whose optimal curve is
\curve{990h3} $=[1,-1,1,-1568,-4669]$ (the \texttt{degphi} file lists the class with this equation and $\degphi=1728$; PARI's \texttt{ellmoddegree} also returns $1728$ for it and
$5184=3\cdot1728$ for \curve{990h1}). Optimality is certified by comparing the period lattices of the curves of the class with the lattice of $\varphi^*$ computed from modular
symbols \cite[Chapter II]{Cremona97}; Section sec:data of the paper says: ``Optimality is determined in the tables for every class of conductor at most $400000$ and for
$359{,}009$ further classes up to $500000$.'' The LMFDB numbers curves differently (\Cref{ch:ellcurves}).
\end{convention}

\begin{example}[the running table]\label{ex:moddeg:running-degrees}
The optimal curves used in the book, with rank and torsion order from \texttt{allcurves}, $\degphi$ from \texttt{degphi} \cite{ecdata} (in every case equal to PARI's
\texttt{ellmoddegree}), and $\dim J_0(N)$:
\begin{center}\small
\begin{tabular}{llllll}
\toprule
curve & $N$ & rank & $\#E(\Q)_{\tors}$ & $\degphi$ & $\dim J_0(N)$\\
\midrule
\curve{11a1} & $11$ & $0$ & $5$ & $1$ & $1$\\
\curve{20a1} & $2^2\cdot5$ & $0$ & $6$ & $1$ & $1$\\
\curve{37a1} & $37$ & $1$ & $1$ & $2$ & $2$\\
\curve{99a1} & $3^2\cdot11$ & $1$ & $2$ & $4$ & $9$\\
\curve{121a1} & $11^2$ & $0$ & $1$ & $6$ & $6$\\
\curve{539c1} & $7^2\cdot11$ & $1$ & $2$ & $288=2^5\cdot3^2$ & $49$\\
\curve{726a1} & $2\cdot3\cdot11^2$ & $1$ & $2$ & $144=2^4\cdot3^2$ & $109$\\
\curve{980i1} & $2^2\cdot5\cdot7^2$ & $0$ & $1$ & $2880=2^6\cdot3^2\cdot5$ & $145$\\
\bottomrule
\end{tabular}
\end{center}
Further optimal curves that occur: \curve{14a1} ($1$), \curve{19a1} ($1$), \curve{27a1} ($1$), \curve{37b1} ($2$), \curve{77c1} ($6$), \curve{99b1} ($12$), \curve{99c1} ($12$), \curve{99d1} ($6$),
\curve{121b1} ($4$), \curve{121c1} ($6$), \curve{121d1} ($24$), \curve{140b1} ($60$), \curve{243b1} ($9$), \curve{539d1} ($72$), \curve{990h3} ($1728$). Checks: the curves of degree $1$ are those
of levels of genus one, as they must be (a map of degree one is an isomorphism, so $X_0(N)\simeq E$); \curve{37a1} has degree $2$ as in \Cref{ex:moddeg:37}; \curve{121a1} is the paper's case (b) of Theorem thm:qual with $e_3=1$ attained, \curve{539c1} its
case (c), \curve{980i1} the $\ell=5$ illustration of Example ex:qual, and \curve{726a1} the curve with image $\Ns$ of the appendix.
\end{example}

\begin{example}[$X_0(37)\to\curve{37a1}$ is optimal]\label{ex:moddeg:37}
Let $\varphi\colon X_0(37)\to X_0(37)/w_{37}\simeq\curve{37a1}$ be the quotient map of \Cref{ex:moddeg:small-levels}(c), of degree $2$. By \Cref{thm:moddeg:optimal}(ii), $\varphi=\psi\circ\varphi_E$
with $E$ the optimal curve of the class \curve{37a}, which has only one curve; so $\psi$ is an endomorphism of \curve{37a1} over $\Q$, that is, $[m]$ for an integer $m\ne0$ (\curve{37a1} has no
complex multiplication, $j=110592/37$), of degree $m^2$. As $2=m^2\degphi$, $m=\pm1$ and $\degphi=2$: the quotient by $w_{37}$ is the optimal parametrisation, without any appeal to
tables.
\end{example}

The modular degree has now been described algebraically (degrees, kernels, lattices). The next section gives the analytic description that the paper quotes and
that PARI uses to compute it, and isolates the one constant on which that description depends.

\section{The Manin constant and the formula of Zagier and Watkins}\label{sec:moddeg:manin}

The paper's introduction says: ``Up to explicit factors $\degphi$ is $\langle f,f\rangle$ divided by the area of the period lattice of $E$ (Zagier \cite{Zagier85}, Watkins
\cite{Watkins02})'', and Remark rem:degree gives the exact formula. This section defines the constant that measures the difference between the N\'eron differential of $E$ and the
differential of the newform, proves the exact formula in three lines, and checks it numerically; the warning after the proof records the normalisations that are easy to get
wrong.

\begin{definition}[the Manin constant]\label{def:moddeg:manin-constant}
Let $\varphi$ be a modular parametrisation of $E$ (conductor $N$, newform $f$). The constant $c_\varphi$ is defined by
\[
\varphi^*\omega_E=c_\varphi\cdot2\pi i\,f(z)\,dz=c_\varphi\,f(q)\,\frac{dq}q\qquad(q=e^{2\pi iz}).
\]
The \emph{Manin constant} of $E$ is $c_E=c_{\varphi_E}$ for the optimal curve, the sign of $\omega_E$ being chosen so that $c_E>0$ (Appendix B, B-C9).
\end{definition}

Why such a constant must appear. On $X_0(N)$ there is a distinguished differential, $2\pi if(z)\,dz=f(q)\,dq/q=(1+a_2q+\cdots)\,dq$, normalised by $a_1(f)=1$; on $E$ there is another,
the N\'eron differential, normalised by the integral structure of $E$ (it is $dx/(2y+a_1x+a_3)$ on a global minimal model, determined up to sign). Pulling the second back gives a
differential on $X_0(N)$ in the line of the first (\Cref{prop:moddeg:manin-rational}(i)), and nothing forces the two normalisations to agree; $c_\varphi$ is the ratio. The sign of
$\omega_E$ is not determined by the minimal model (the change $y\mapsto-y-a_1x-a_3$ negates it), which is why one normalises $c_E>0$. Example: for the optimal curve \curve{11a1}
the constant is $1$, and for the parametrisation of \curve{11a3} through the $5$-isogeny it is $5$ (\Cref{ex:moddeg:zw-numerics}).

\begin{proposition}[the constant $c_\varphi$]\label{prop:moddeg:manin-rational}
In \Cref{def:moddeg:manin-constant}:
\begin{enumerate}[label=(\roman*)]
\item $\varphi^*\omega_E$ is a nonzero multiple of $\omega_f=f(z)\,dz$, so $c_\varphi$ is well defined and nonzero.
\item $c_\varphi\in\Q^\times$.
\item If $\varphi=\psi\circ\varphi_{E_0}$ with $E_0$ optimal and $\psi\colon E_0\to E$ an isogeny (\Cref{thm:moddeg:optimal}(ii)), and $\psi^*\omega_E=u_\psi\,\omega_{E_0}$, then $c_\varphi=c_{E_0}u_\psi$;
$u_\psi$ is a nonzero integer (stated), and $u_\psi u_{\hat\psi}=\deg\psi$.
\end{enumerate}
\end{proposition}

\begin{proof}
(i) By \Cref{thm:moddeg:optimal}(ii), $\varphi=\psi\circ\varphi_{E_0}$. The pullback of an invariant differential by an isogeny is a constant multiple of the invariant differential
(\Cref{ch:ellcurves}), so $\varphi^*\omega_E=u_\psi\varphi_{E_0}^*\omega_{E_0}$, and $\varphi_{E_0}^*\omega_{E_0}\in\C\omega_f$ by \Cref{thm:moddeg:shimura-quotient}(ii). It is nonzero because $\varphi$ is
nonconstant (a nonconstant holomorphic map is a local biholomorphism off a finite set). Also $2\pi i\,dz=dq/q$.
(ii) The function $1/j$ lies in $\Q(X_0(N))$ and has $q$-expansion $q-744q^2+\cdots\in\Q[[q]]$ at the cusp $\infty$, which has width $1$; so $q$ is a local parameter at the rational point
$\infty$ for the $\Q$-structure, and the completed local ring there is $\Q[[q]]$. The differential $\varphi^*\omega_E$ is defined over $\Q$, so its expansion $c_\varphi\sum_na_n(f)q^{n-1}dq$ has
rational coefficients; as $a_1(f)=1$, $c_\varphi\in\Q$.
(iii) $c_\varphi$ is computed from $\varphi^*\omega_E=\varphi_{E_0}^*\psi^*\omega_E=u_\psi\varphi_{E_0}^*\omega_{E_0}$. That $u_\psi\in\Z$ uses N\'eron models: $\psi$ extends to the N\'eron models over $\Z$, and the
N\'eron differential generates the invariant differentials of the N\'eron model (\Cref{ch:reduction}; \cite[Chapter II]{Cremona97}); we do not prove it here. From $\hat\psi\circ\psi=[\deg\psi]$
and $[m]^*\omega=m\omega$ (\Cref{ch:ellcurves}), $u_\psi u_{\hat\psi}=\deg\psi$.
\end{proof}

\begin{theorem}[what is known about the Manin constant; stated]\label{thm:moddeg:manin-known}
Let $E/\Q$ be optimal.
\begin{enumerate}[label=(\roman*)]
\item (Edixhoven \cite{Edixhoven91}.) $c_E\in\Z$; hence $c_\varphi\in\Z$ for every modular parametrisation, by \Cref{prop:moddeg:manin-rational}(iii).
\item (Manin's conjecture.) $c_E=1$. (The literature writes $c_E=\pm1$ when the sign of $\omega_E$ is not normalised.)
\item If a prime $p$ divides $c_E$, then $p=2$ or $p^2\mid N$: Mazur \cite{Mazur78} for $p\nmid2N$, Abbes and Ullmo \cite{AU96} for $p\mid N$, $p^2\nmid N$.
\item $c_E=1$ for every semistable optimal $E$ (\v Cesnavi\v cius \cite{Cesnavicius18}, completing the cases left by \cite{Mazur78,AU96} at $p=2$); further results relate $c_E$ to
$\degphi$ (\v Cesnavi\v cius, Neururer and Saha \cite{CNS}), and computations of Cremona and of Agashe, Ribet and Stein \cite{ARS06} confirm $c_E=1$ for all optimal curves in large
ranges of conductors.
\item For non-optimal parametrisations $c_\varphi\ne1$ is possible (\Cref{ex:moddeg:zw-numerics}).
\end{enumerate}
\end{theorem}

We do not prove this. Items (iii) and (iv) are stated as we recall the sources; the exact hypotheses, in particular at $p=2$ and the ranges of the computations, were not
checked against them, and the book uses none of (i) to (iv) in a proof. The paper uses the Manin constant only in Remark rem:degree, where $c_E^2$ appears explicitly, and in
the heuristic sentence of Section 1.

\begin{theorem}[the formula of Zagier and Watkins]\label{thm:moddeg:zagier-watkins}
For every modular parametrisation $\varphi$ of $E$, with the unnormalised Petersson product,
\[
\deg\varphi\cdot\vol(\Lambda_E)=4\pi^2c_\varphi^2\,\langle f,f\rangle.
\]
In particular, for the optimal curve,
\[
\degphi=\frac{4\pi^2c_E^2\langle f,f\rangle}{\vol(\Lambda_E)},\quad\text{which is }\frac{4\pi^2\langle f,f\rangle}{\vol(\Lambda_E)}\text{ if }c_E=1.
\]
\end{theorem}

\begin{proof}
For a nonconstant map $\varphi$ of compact Riemann surfaces and a smooth $2$-form $\rho$ on the target, $\int_{X_0(N)}\varphi^*\rho=\deg\varphi\int_E\rho$ (\Cref{ch:curves}: off the
finitely many branch points $\varphi$ is a $\deg\varphi$-sheeted covering, each sheet maps diffeomorphically and orientation-preservingly onto its image, and a finite set has measure zero).
With $\rho=\frac i2\omega_E\wedge\bar\omega_E$ the right side is $\deg\varphi\cdot\vol(\Lambda_E)$: under $E(\C)\simeq\C/\Lambda_E$ the form $\omega_E$ becomes $du$, and
$\frac i2du\wedge d\bar u$ is the Lebesgue measure on the fundamental parallelogram. On the left, $\varphi^*\omega_E=2\pi ic_\varphi\,\omega_f$ with $c_\varphi$ real
(\Cref{prop:moddeg:manin-rational}(ii)), so $\overline{\varphi^*\omega_E}=-2\pi ic_\varphi\,\overline{\omega_f}$ and
\[
\varphi^*\rho=\frac i2(2\pi ic_\varphi)(-2\pi ic_\varphi)\,\omega_f\wedge\overline{\omega_f}=4\pi^2c_\varphi^2\cdot\frac i2\,\omega_f\wedge\overline{\omega_f}.
\]
Its integral over $X_0(N)$ is $4\pi^2c_\varphi^2\langle f,f\rangle$ by \Cref{conv:moddeg:setup}(C4): on a fundamental domain, $\frac i2\,f\,dz\wedge\overline{f\,dz}=\abs f^2\cdot\frac i2dz\wedge d\bar z=\abs f^2\,dx\wedge dy$,
since $dz\wedge d\bar z=(dx+i\,dy)\wedge(dx-i\,dy)=-2i\,dx\wedge dy$.
\end{proof}

For example, for \curve{37a1}, with the numbers of \Cref{ex:moddeg:zw-numerics}, $4\pi^2\cdot0.37175\ldots/7.33813\ldots=2.0000\ldots$, the degree of the quotient map by $w_{37}$ (\Cref{ex:moddeg:37}).

In the language of physics: the $L^2$ norm of a pulled-back one-form counts every sheet of the cover, so the area of a $d$-sheeted cover is $d$ times the area of the base.

\begin{warning}[normalisations in the formula of Zagier and Watkins]\label{warn:moddeg:zw-normalisation}
(1) The Manin constant stands in the \emph{numerator}: $\degphi=4\pi^2c_E^2\langle f,f\rangle/\vol(\Lambda_E)$. The check is the class \curve{11a} (\Cref{ex:moddeg:zw-numerics}): the parametrisation of
\curve{11a3} has degree $5$ and $c_\varphi=5$, and $4\pi^2\langle f,f\rangle/\vol(\Lambda_{\curve{11a3}})=1/5=5/5^2$. (2) $\langle f,f\rangle$ is unnormalised; PARI's \texttt{mfpetersson} returns
$\langle f,f\rangle/\mu(N)$, and books that divide by $\vol(\Gamma_0(N)\backslash\HH)=\pi\mu(N)/3$ change the constant $4\pi^2$. (3) $\vol(\Lambda_E)$ is the area of a fundamental
parallelogram of the period lattice of the N\'eron differential, $\abs{\operatorname{Im}(\bar\omega_1\omega_2)}$ for PARI's periods \texttt{E.omega} $=[\omega_1,\omega_2]$; it is not the real
period $\Omega_E$. The normalisations of \cite{Zagier85,Watkins02} were not compared with these; the formula above is proved directly.
\end{warning}

\begin{example}[numerical check]\label{ex:moddeg:zw-numerics}
Unnormalised Petersson norms ($\mu(N)$ times \texttt{mfpetersson}), period areas, and the quotient $4\pi^2\langle f,f\rangle/\vol(\Lambda_E)=\deg\varphi/c_\varphi^2$, computed with PARI to $40$ digits
(the last column agrees with a rational number to all of them):
\begin{center}\small
\begin{tabular}{llllll}
\toprule
curve & $\langle f,f\rangle$ & $\vol(\Lambda_E)$ & $4\pi^2\langle f,f\rangle/\vol$ & $\deg\varphi$ & $c_\varphi$\\
\midrule
\curve{11a1} & $0.04690014787349518782$ & $1.85154362345595931771$ & $1$ & $1$ & $1$\\
\curve{11a2} & same & $0.37030872469119186354$ & $5$ & $5$ & $1$\\
\curve{11a3} & same & $9.25771811727979658854$ & $1/5$ & $5$ & $5$\\
\curve{20a1} & $0.08134945682538172857$ & $3.21154782844006508758$ & $1$ & $1$ & $1$\\
\curve{37a1} & $0.37175414751069605028$ & $7.33813274078957673907$ & $2$ & $2$ & $1$\\
\curve{99a1} & $0.53851956871086065179$ & $5.31497510542145355032$ & $4$ & $4$ & $1$\\
\curve{121a1} & $0.42823936001195717858$ & $2.81770204819580190671$ & $6$ & $6$ & $1$\\
\bottomrule
\end{tabular}
\end{center}
In the class \curve{11a} the parametrisations of \curve{11a2} and \curve{11a3} both have degree $5$ (\Cref{thm:moddeg:optimal}(ii)), so the formula forces $c_\varphi=1$ for \curve{11a2} and $c_\varphi=5$ for
\curve{11a3}: the $5$-isogeny $\psi\colon\curve{11a1}\to\curve{11a3}$ has $u_\psi=5$, and its dual $u_{\hat\psi}=1$. For the optimal curves of the table the computation gives $c_E^2=1$, in
agreement with Manin's conjecture (the degrees being those of Cremona's tables).
\end{example}

\begin{remark}[Rankin's method and the computation of $\degphi$]\label{rem:moddeg:rankin-and-computation}
(a) Rankin and Shimura express $\langle f,f\rangle$ through the special value $L^{\nv}(\Sym f,2)$ of the naive symmetric square $L$-function: $\langle f,f\rangle=N\,L^{\nv}(\Sym f,2)/(8\pi^3)$
(Remark rem:degree of the paper, citing \cite{Rankin39,Shimura76} and \cite[Lemma 2.12]{DFG04}; \Cref{ch:lfunctions}). Combined with \Cref{thm:moddeg:zagier-watkins} this is Watkins's
method of computing $\degphi$ \cite{Watkins02}, implemented in PARI's \texttt{ellmoddegree}, which assumes $c_E=1$. (b) The Rankin series $\sum a_n^2n^{-s}$ is not the naive symmetric
square $L$-function (\Cref{ch:lfunctions}; Appendix B, B-C7). (c) In the book's normalisation the ``explicit factors'' of the paper's Section 1 are $4\pi^2c_E^2$.
\end{remark}

\begin{paperlink}\label{pl:moddeg:zw}
Remark rem:degree: ``By the formula of Zagier and Watkins \cite{Zagier85,Watkins02}, $\degphi\cdot\operatorname{vol}(\Lambda_E)=4\pi^2c_E^2\langle f,f\rangle$, with the unnormalised Petersson product
and the Manin constant $c_E>0$ in the numerator'' is \Cref{thm:moddeg:zagier-watkins} for $\varphi=\varphi_E$. The remark's check ``With $c_E=1$ the first formula reproduces $\degphi$ to $30$
digits for 11a1, 121a1, 539d1, 1232l1 and 726a1'' is of the same kind as \Cref{ex:moddeg:zw-numerics}.
\end{paperlink}

The formula of Zagier and Watkins says how large $\degphi$ is; it says nothing about which primes divide it. For that the paper uses congruences between $f$ and
the other forms of level $N$, the subject of the next section.

\section[The congruence number and Agashe-Ribet-Stein]{The congruence number and the theorem of Agashe, Ribet and Stein}\label{sec:moddeg:congruence}

The paper reaches the modular degree through the congruence number: ``By a theorem of Agashe, Ribet and Stein \cite[Theorem 2.2]{ARS12}, $\degphi$ divides the
\emph{congruence number} $r_E$ \dots with $v_p(\degphi)=v_p(r_E)$ for every prime $p$ with $p^2\nmid N$'' (Section 1), and the proof of Theorem thm:qual reduces to ``$\ell\mid r$''. This
section defines $r_E$ and the Hecke congruence ideal, states Lemma lem:crit (proved in \Cref{ch:modgalrep}) and the theorem of Agashe, Ribet and Stein, and computes $r_E$
directly for the curves of the book, \curve{99a1} included. The divisibility half of the theorem is proved in \Cref{sec:moddeg:degree-eta} at every prime.

\begin{definition}[congruence number and congruence ideal]\label{def:moddeg:congruence-number}
Let $f=f_E$. The \emph{congruence number} $r_E$ is the largest integer $r$ such that $f\equiv g\pmod r$ coefficientwise for some $g\in S_2(\Gamma_0(N);\Z)\cap f^\perp$. The
\emph{Hecke congruence ideal} of $f$ is $\eta_f=\phi_f(\Ann_\T I_f)\subseteq\Z$, and for a prime $\ell$, $\eta_f\Zl=(\phi_f\otimes\Zl)(\Ann_{\T_\ell}(I_f\T_\ell))$. (These are the definitions of
\Cref{ch:hecke} for $f=f_E$; $r_E$ depends only on the isogeny class.)
\end{definition}

The two definitions measure the same phenomenon from the two sides of the duality between forms and Hecke operators. Over $\Q$ the space of cusp forms splits as $\Q f$ plus the
complement $f^\perp$; over $\Z$ it does not, because an integral form can be congruent to $f$ modulo $r$ and to a form in $f^\perp$ at the same time. The congruence number is the
largest such $r$, and $r_E>1$ exactly when some prime $\ell$ makes $f$ congruent modulo $\ell$ to a combination of the other forms (for an eigenform, typically, to another
eigenform, as in Lemma lem:crit below). The Hecke version asks the same question inside $\T$: $\T\otimes\Q=\Q\times\T'$, and $\eta_f$ measures how far $\T$ is from the product of its two
projections. At level $11$, $S_2(\Gamma_0(11))=\C f_{11}$, $f^\perp=0$ and $\T=\Z$, so $r_E=1$ and $\eta_f=\phi_f(\Ann_\Z0)=\Z$. The same holds for every conductor below $26$ (they are $11$, $14$, $15$, $17$, $19$, $20$, $21$, $24$, and $X_0(N)$ has genus one for each,
so $f^\perp=0$). At $N=26$, where $X_0(26)$ has genus $2$ and the classes \curve{26a}, \curve{26b} share it, $r_E=2$ for both, as \Cref{thm:moddeg:ars} gives from $\degphi=2$ ($26$ is
square-free); level $37$ is worked out by hand in \Cref{ex:moddeg:ex-level37-hecke}.

\begin{proposition}[congruence ideals of a newform]\label{prop:moddeg:congruence-ideal}
\leavevmode
\begin{enumerate}[label=(\roman*)]
\item $\T\otimes\Q=\Q\times\T'$, with $\phi_f\otimes\Q$ the first projection; hence $I_f\cap\Ann_\T(I_f)=0$ and $\eta_f\ne0$.
\item Let $\mf m=\mf m_{f,\ell}$ be the maximal ideal of $\T_\ell$ containing $\ell$ and $I_f$. Then $\eta_f\Zl$ is the congruence ideal of $\phi_f$ restricted to the local factor $\T_{\mf m}$: the other
local factors of $\T_\ell$ lie in $\ker(\phi_f\otimes\Zl)$, not in its annihilator.
\item $r_E$ is finite: $\Z f\oplus(f^\perp\cap S_2(\Gamma_0(N);\Z))$ has finite index in $S_2(\Gamma_0(N);\Z)$.
\end{enumerate}
\end{proposition}

\begin{proof}
(i) $\T\otimes\Q$ is a finite-dimensional commutative $\Q$-algebra acting faithfully on $S_2(\Gamma_0(N))$ (\Cref{prop:moddeg:hecke-compatibility}); let $A$ be its local factor through which
$\phi_f$ factors and $e$ its idempotent. Then $eS_2$ is the generalised eigenspace of $\phi_f$, which is $\C f$: the anemic algebra acts semisimply on $S_2(\Gamma_0(N))$, its operators being normal for
the Petersson product (\Cref{ch:hecke}), so the generalised $\phi_f$-eigenspace of $\T$ lies in the eigenspace of the anemic algebra for the restriction of
$\phi_f$, which is $\C f$ (\Cref{prop:moddeg:f-lattice}(i)); and $A\otimes\C$ acts faithfully on the line
$eS_2$, so $A\otimes\C$ embeds in $\C$ and $A=\Q$. The rest is \Cref{ch:modules} (the proposition on congruence ideals, parts (ii), (iii)). (ii) is part (v) of the same proposition,
$\Zl$ being complete. (iii) is \Cref{ch:hecke}: both summands are saturated and their ranks add up, and a congruence modulo $r$ gives the element $(f-g)/r$ of order $r$ in the
quotient.

In more detail, (iii) also identifies $r_E$ with the index. Write $S=S_2(\Gamma_0(N);\Z)$ and $L=f^\perp\cap S$. With $e$ the idempotent of the factor $\Q$ in (i), $eS_2=\C f$ and
$(1-e)S_2$ is the sum of the other generalised eigenspaces; as $f^\perp$ is $\T$-stable, $f^\perp=ef^\perp\oplus(1-e)f^\perp$ with $ef^\perp\subseteq\C f\cap f^\perp=0$, so $f^\perp=(1-e)S_2$
by dimensions. Since $e$ has rational entries in a basis of $S$, $S\otimes\Q=\Q f\oplus(L\otimes\Q)$, and the coefficient of $f$ is a linear map $P\colon S\to\Q$ with $P(f)=1$ and kernel $L$. Its
image is a finitely generated subgroup of $\Q$ containing $1$, so it is $\frac1{r'}\Z$ for a positive integer $r'$, and $P$ induces $S/(\Z f\oplus L)\simeq\frac1{r'}\Z/\Z\simeq\Z/r'$. If
$f\equiv g\pmod r$ with $g\in L$, then $h=(f-g)/r\in S$ and $P(h)=1/r$, so $r\mid r'$; conversely, if $P(h)=1/r'$, then $g=f-r'h$ lies in $L$ and $f\equiv g\pmod{r'}$. Hence $r_E=r'$ is
the index of $\Z f\oplus L$ in $S$, and the quotient is cyclic. This is the ``cyclic-quotient argument'' used in \Cref{ex:moddeg:99a1-and-726a1}.
\end{proof}

\begin{lemma}[the paper's Lemma lem:crit; stated, proved in \Cref{ch:modgalrep}]\label{lem:moddeg:crit}
Let $J=\Ann_\T(f^\perp)$. Then $J=\Ann_\T(I_f)$ and $\T/(I_f+J)\simeq\Z/r_E$; equivalently $r_E\Z=\eta_f$. The same holds over $\Zl$, with $r_E$ replaced by its $\ell$-part:
$r_E\Zl=\eta_f\Zl$. If a normalised eigenform $g\in f^\perp$ with coefficients in the ring of integers of a number field satisfies $a_n(g)\equiv a_n(f)\pmod\lambda$ for every $n$, for a
prime $\lambda\mid\ell$, then $\ell\mid r_E$.
\end{lemma}

This chapter uses only $r_E\Zl=\eta_f\Zl$ (in \Cref{cor:moddeg:ars-divisibility,thm:moddeg:lem-hecke}).

\begin{theorem}[Agashe, Ribet and Stein; stated]\label{thm:moddeg:ars}
Let $E$ be optimal of conductor $N$. Then $\degphi\mid r_E$, and $v_p(\degphi)=v_p(r_E)$ for every prime $p$ with $p^2\nmid N$.
\end{theorem}

We do not prove the theorem as a whole; it is \cite[Theorem 2.2]{ARS12}, as the paper cites it (the number is the paper's; we have not checked it against the published
text). This chapter proves the divisibility at every prime (\Cref{cor:moddeg:ars-divisibility}) and the equality at primes $\ell\nmid2N$ for which $\mf m_{f,\ell}$ is not Eisenstein, given
multiplicity one (\Cref{thm:moddeg:gorenstein}); the equality at $p\,\|\,N$, at $p=2$ and in the Eisenstein case is theirs. The idea of their proof: with $A$ and $B$ as in
\Cref{prop:moddeg:degree-descriptions}, $\degphi$ is the exponent of $A\cap B$, equivalently of $H_1(J_0(N),\Z)/(H_1(A,\Z)+H_1(B,\Z))$, while $r_E$ is the exponent of
$S_2(\Z)/(S_2(\Z)[I_f]+S_2(\Z)[J])$; the perfect pairing $S_2(\Z)\times\T\to\Z$ (\Cref{ch:hecke}) and the action of $\T$ on $H_1$ compare the two, and they agree at $p$ when $H_1(X_0(N),\Zp)$ is
locally free of rank two over $\T\otimes\Zp$ at the relevant maximal ideal, which they establish for $p^2\nmid N$ from results of Mazur, Ribet and others. Zagier \cite{Zagier85} first
compared $\degphi$ with congruence primes, and Ribet \cite{Ribet83} proved the divisibility $\degphi\mid r_E$ in an earlier form.

\begin{remark}[the size of the degree]\label{rem:moddeg:abc-size}
Section 1 of the paper: ``the conjectured bound $\degphi\ll_\varepsilon N^{2+\varepsilon}$ is equivalent to a form of the $abc$ conjecture, see \cite[\S2]{ARS12} and the references there''. The
references there include Frey \cite{Frey} and Mai and Murty \cite{MaiMurty}. Through \Cref{thm:moddeg:zagier-watkins} the bound is a bound for $\langle f,f\rangle/\vol(\Lambda_E)$; through
\Cref{thm:moddeg:ars} it bounds the congruence primes of $f$ (compare \cite{Murty99}). Nothing in the paper depends on it.
\end{remark}

\begin{example}[congruence numbers computed directly]\label{ex:moddeg:99a1-and-726a1}
The following algorithm computes $r_E$. (1) Take a $\Q$-basis of $S_2(\Gamma_0(N))$ with rational $q$-expansions (\texttt{mfbasis(mfinit([N,2],1))}), and compute the lattice
$S_2(\Gamma_0(N);\Z)$ of forms with integral coefficients as the saturation of the $\Q$-span of the truncated expansions in $\Z^M$, $M$ beyond Sturm's bound (\texttt{matrixqz(A,-2)}); a
rational form whose first $M$ coefficients are integers is integral by Sturm's theorem (\Cref{ch:modforms}). (2) Compute the matrices of $T_p$ ($U_p$ for $p\mid N$) on this lattice from
$a_n(T_pg)=a_{np}(g)+pa_{n/p}(g)$ (the second term for $p\nmid N$ only). (3) Find the primitive integral functional $\theta$ on the lattice with $\theta\circ T_p=a_p(f)\theta$ for
$p\le23$ (the common eigenline was already one-dimensional): its kernel is $f^\perp\cap S_2(\Gamma_0(N);\Z)$, because $f^\perp$ is the sum of the other generalised eigenspaces. (4) Then $\theta\colon S_2(\Gamma_0(N);\Z)\to\Z$ is onto with
kernel $L=f^\perp\cap S_2(\Z)$, so $\Z f\oplus L$ has index $\abs{\theta(f)}$ (indeed $\theta=\pm r'P$ for the map $P$ of the proof of \Cref{prop:moddeg:congruence-ideal}, both having kernel
$L$ and $\theta$ being onto $\Z$), and $r_E=\abs{\theta(f)}$ by the cyclic-quotient argument given there. The results, with
$\degphi$ from \texttt{ellmoddegree} and \texttt{degphi}:
\begin{center}\small
\begin{tabular}{lllll}
\toprule
curve & $N$ & $\degphi$ & $r_E$ & primes $p$ with $p^2\mid N$\\
\midrule
\curve{11a1} & $11$ & $1$ & $1$ & none\\
\curve{20a1} & $2^2\cdot5$ & $1$ & $1$ & $2$\\
\curve{37a1} & $37$ & $2$ & $2$ & none\\
\curve{77c1} & $7\cdot11$ & $6$ & $6$ & none\\
\curve{88a1} & $2^3\cdot11$ & $8$ & $16$ & $2$\\
\curve{99a1} & $3^2\cdot11$ & $4$ & $12$ & $3$\\
\curve{99b1} & $3^2\cdot11$ & $12$ & $12$ & $3$\\
\curve{99c1} & $3^2\cdot11$ & $12$ & $12$ & $3$\\
\curve{99d1} & $3^2\cdot11$ & $6$ & $6$ & $3$\\
\curve{121a1} & $11^2$ & $6$ & $6$ & $11$\\
\curve{121b1} & $11^2$ & $4$ & $4$ & $11$\\
\curve{121c1} & $11^2$ & $6$ & $6$ & $11$\\
\curve{121d1} & $11^2$ & $24$ & $24$ & $11$\\
\curve{140b1} & $2^2\cdot5\cdot7$ & $60$ & $60$ & $2$\\
\curve{539c1} & $7^2\cdot11$ & $288$ & $288$ & $7$\\
\curve{539d1} & $7^2\cdot11$ & $72$ & $504$ & $7$\\
\curve{726a1} & $2\cdot3\cdot11^2$ & $144$ & $1584$ & $11$\\
\curve{980i1} & $2^2\cdot5\cdot7^2$ & $2880$ & $2880$ & $2$, $7$\\
\bottomrule
\end{tabular}
\end{center}
So \curve{99a1} has $\degphi=4$ and $r_E=12$, as \cite[Table 1]{ARS12} and the remark after the paper's Example ex:qual say, and the discrepancy is at $3$, with $3^2\mid99$. The congruence behind
$3\mid r_E$: $a_n(f_{99a})\equiv a_n(f_{99b})\equiv a_n(f_{99d})\pmod3$ for all $n\le60$, beyond Sturm's bound $24$ for $\Gamma_0(99)$ (which has index $99\cdot\frac43\cdot\frac{12}{11}=144$
in $\SL_2(\Z)$, and $2\cdot144/12=24$), so for all $n$ (\Cref{ch:modforms}). The first coefficients, from \texttt{ellan},
\begin{center}\small
\begin{tabular}{lrrrrrrrr}
\toprule
$n$ & $2$ & $4$ & $5$ & $7$ & $8$ & $10$ & $11$ & $13$\\
\midrule
$a_n(f_{99a})$ & $-1$ & $-1$ & $-4$ & $-2$ & $3$ & $4$ & $-1$ & $-2$\\
$a_n(f_{99b})$ & $-1$ & $-1$ & $2$ & $4$ & $3$ & $-2$ & $-1$ & $-2$\\
$a_n(f_{99d})$ & $2$ & $2$ & $-1$ & $-2$ & $0$ & $-2$ & $-1$ & $4$\\
\bottomrule
\end{tabular}
\end{center}
show the pattern (the coefficients $a_n$ with $3\mid n$ vanish for all three, $3^2$ dividing the level): the differences are $0$ or $\pm3$, $\pm6$. The forms
$f_{99b}$, $f_{99d}$ are newforms of level $99$ different from $f_{99a}$, hence in $f^\perp$, and Lemma lem:crit gives $3\mid r_E$. ($f_{99c}$ is not congruent to $f_{99a}$ modulo $3$:
$a_2=1$ against $-1$.) The paper's curves \curve{726a1} (ratio $11$, $11^2\mid726$) and \curve{539d1} (ratio $7$, $7^2\mid539$) are further examples, and \curve{88a1} (ratio $2$, $2^3\mid88$) a fourth (\Cref{ex:moddeg:ex-88}). In every other case of the table the two
numbers agree, also at primes $p$ with $p^2\mid N$ (\curve{121a1}, \curve{539c1}, \curve{980i1}, \curve{99b1}, \curve{99d1}): the equality can hold there; only the guarantee is lost.
\end{example}

\begin{paperlink}\label{pl:moddeg:99a1}
The remark after the Example following the proof of Theorem thm:qual: ``The last step needs only $\ell^2\nmid N$, and it cannot be dropped there: Table 1 of \cite{ARS12} lists the curve 99A1, of
conductor $9\cdot11$, with modular degree $4$ and congruence number $12$; likewise 726a1 has $\degphi=144$ and $r_E=1584=11\cdot144$, with $11^2\mid726$, and 539d1 has $\degphi=72$ and
$r_E=504=7\cdot72$, with $7^2\mid539$.'' All three pairs are recomputed in \Cref{ex:moddeg:99a1-and-726a1}. The opening of Section sec:proof (``By \cite[Theorem 2.2]{ARS12}, $\deg\varphi_{E_0}$ divides
$r$ and $v_p(\deg\varphi_{E_0})=v_p(r)$ for every prime $p$ with $p^2\nmid N$; as $\ell\nmid N$ it suffices to prove $\ell\mid r$'') uses \Cref{thm:moddeg:ars} at $p=\ell$, where $\ell^2\nmid N$.
\end{paperlink}

The theorem of Agashe, Ribet and Stein compares $\degphi$, a geometric quantity, with $r_E$, a quantity about $q$-expansions. The next section explains the
comparison from the cohomological side, where $\degphi$ turns out to be exactly the congruence ideal attached to the pairing $\hat\delta_!$.

\section{The modular degree as a congruence ideal}\label{sec:moddeg:degree-eta}

This section proves the paper's Proposition prop:pullback in the following form: for the optimal curve the lattice $M_{f,\lambda}$ is exactly $\varphi_E^*H^1(E,\Zl)$, the pairing
$\hat\delta_!$ has ideal of values $\degphi\cdot\Zl$ on it, and this holds for every prime $\ell$ with no hypothesis on $\bar\rho$ or on multiplicity one. The inputs are optimality (the
kernel of $J_0(N)\to E$ is connected, so the pullback is saturated), Hecke equivariance, the Atkin-Lehner sign and the projection formula. Two consequences follow: the divisibility
$\degphi\mid r_E$ of Agashe, Ribet and Stein at every prime, by the inequality half of Lemma lem:ddt417 (\Cref{cor:moddeg:ars-divisibility}); and Lemma lem:hecke of the paper
(\Cref{thm:moddeg:lem-hecke}). The example \curve{99a1} at $\ell=3$ shows that the Hecke-side inequality can be strict.

\begin{proposition}[$M_{f,\lambda}$ is the pullback of $H^1(E)$]\label{prop:moddeg:mf-is-pullback}
Let $E$ be optimal with newform $f$, $\varphi=\varphi_E$, and $\ell$ any prime.
\begin{enumerate}[label=(\alph*)]
\item (Hecke equivariance.) $\varphi_*\circ t_*=\phi_f(t)\varphi_*$ on $H_1(X_0(N),\Z)$ and $t^*\circ\varphi^*=\phi_f(t)\varphi^*$ on $H^1(E,\Z)$, for every $t\in\T$; so $\varphi^*H^1(E,\Z)\subseteq H^1(X_0(N),\Z)[I_f]$.
\item (Saturation.) $\varphi^*\colon H^1(E,\Z)\to H^1(X_0(N),\Z)$ is injective with saturated image.
\item $\varphi^*H^1(E,\Z)=H^1(X_0(N),\Z)[I_f]$, and for every prime $\ell$
\[
M_{f,\lambda}=H^1(X_0(N),\Zl)[I_f]=\varphi^*H^1(E,\Zl).
\]
\item (Galois structure.) In \'etale cohomology $\varphi^*$ is $G_\Q$-equivariant, so $M_{f,\lambda}\simeq H^1_{\mathrm{et}}(E_{\Qbar},\Zl)\simeq(T_\ell E)^\vee\simeq T_\ell E(-1)$ as $\Zl[G_\Q]$-modules, and
$\ad M_{f,\lambda}\simeq\ad T_\ell E$.
\end{enumerate}
\end{proposition}

\begin{proof}
(a) $E=E_f=J_0(N)/I_fJ_0(N)$ and $\varphi_*=\pi_f$ (\Cref{thm:moddeg:optimal}(iv), \Cref{def:moddeg:parametrisation}); $t$ acts on $E_f$ as $\phi_f(t)$ (\Cref{thm:moddeg:shimura-quotient}(ii)), so
$\pi_f\circ t_*=\phi_f(t)\pi_f$ on $J_0(N)$ and on $H_1$. Transposing, with $t^*$ the transpose of $t_*$ and $\varphi^*$ that of $\varphi_*$ (\Cref{prop:moddeg:hecke-compatibility}(ii)), gives
$t^*\varphi^*=\phi_f(t)\varphi^*$; for $t\in I_f$ this is $0$.
(b) $\ker\varphi_*$ is connected (\Cref{thm:moddeg:optimal}(iii)(b)), so $\varphi_*$ is surjective on $H_1$ and $\varphi^*H^1(E,\Z)$ is saturated (\Cref{ch:curves}, the proposition on the degree
and the cup product). The mechanism is elementary: $H^1=\Hom(H_1,\Z)$ and $\varphi^*=\Hom(\varphi_*,\Z)$; a surjection $\varphi_*$ onto the free module $H_1(E,\Z)$ has a section
$s$, and then $s^*\circ\varphi^*=\mathrm{id}$, so $\varphi^*H^1(E,\Z)$ is a direct summand of $H^1(X_0(N),\Z)$, hence saturated. (For the parametrisation $\psi$ of \curve{11a3}, $\psi_*H_1$
has index $5$ and $\psi^*H^1(\curve{11a3},\Z)$ is not saturated: it contains $5x$ but not $x$ for a suitable class $x$.) $\varphi^*$ is injective because $\varphi_*\varphi^*=\deg\varphi\ne0$ and $H^1(E,\Z)$ is torsion-free.
(c) $H^1(X_0(N),\Z)[I_f]$ has rank $2$ (\Cref{prop:moddeg:f-lattice}), as has $\varphi^*H^1(E,\Z)$, which it contains by (a). The quotient $H^1(X_0(N),\Z)[I_f]/\varphi^*H^1(E,\Z)$ is then finite, and it
embeds in $H^1(X_0(N),\Z)/\varphi^*H^1(E,\Z)$, which is torsion-free by (b); so it is $0$. Tensoring with the flat $\Z$-module $\Zl$ and using \Cref{prop:moddeg:f-lattice}(ii) gives the statement
over $\Zl$.
(d) The comparison isomorphism between Betti and \'etale cohomology is compatible with pullbacks (\Cref{ch:curves}), and $\varphi$ is defined over $\Q$; so $\varphi^*\colon H^1_{\mathrm{et}}(E_{\Qbar},\Zl)\to
H^1_{\mathrm{et}}(X_0(N)_{\Qbar},\Zl)$ is $G_\Q$-equivariant, injective with image $M_{f,\lambda}$ by (c). $H^1_{\mathrm{et}}(E_{\Qbar},\Zl)\simeq(T_\ell E)^\vee\simeq T_\ell E(-1)$ by the Weil pairing
(\Cref{ch:curves}; Appendix B, B-C5(b)), and the trace-zero adjoint is insensitive to duals and Tate twists (\Cref{ch:reps}; B-C12(b)).
\end{proof}

\begin{lemma}[$w_N$ on $M_{f,\lambda}$]\label{lem:moddeg:wn-on-mf}
Let $E$ be optimal with newform $f$. As morphisms $X_0(N)\to E$,
\[
\varphi_E\circ w_N=\varepsilon_N(f)\,\varphi_E+\varphi_E(0),
\]
and $w_N^*$ acts on $M_{f,\lambda}=\varphi_E^*H^1(E,\Zl)$ by $\varepsilon_N(f)$. Hence $\hat\delta_!(x,y)=\varepsilon_N(f)\langle x,y\rangle$ for $x,y\in M_{f,\lambda}$.
\end{lemma}

\begin{proof}
Write $\varepsilon_N=\varepsilon_N(f)$ and put $h(P)=\varphi_E(w_NP)-\varphi_E(0)$. Then $h$ is a morphism over $\Q$ (\Cref{thm:moddeg:x0-over-q}(ii)) with $h(\infty)=\varphi_E(0)-\varphi_E(0)=O$, as
$w_N(\infty)=0$. The differential $\omega_E$ is translation invariant, so
\[
h^*\omega_E=w_N^*\varphi_E^*\omega_E=c_E\cdot2\pi i\,(f|_2W_N)(z)\,dz=\varepsilon_N\,c_E\cdot2\pi if(z)\,dz=\varepsilon_N\,\varphi_E^*\omega_E=(\varepsilon_N\varphi_E)^*\omega_E,
\]
using \Cref{conv:moddeg:setup}(C1) and $[-1]^*\omega_E=-\omega_E$. The homomorphisms $h_*$ and $(\varepsilon_N\varphi_E)_*\colon J_0(N)\to E$ are induced on $J_0(N)(\C)=S_2^\vee/H_1$ by the transposes of the
pullbacks of $\omega_E$ (\Cref{ch:curves}, functoriality of the Jacobian), which agree; so $h_*=(\varepsilon_N\varphi_E)_*$ and $h=h_*\circ\iota=\varepsilon_N\varphi_E$. On $H^1$, translations act trivially
(they are homotopic to the identity) and $[\varepsilon_N]^*=\varepsilon_N$, so $w_N^*\varphi_E^*=(\varphi_E\circ w_N)^*=\varepsilon_N\varphi_E^*$, first on $H^1(E,\Z)$ and then on $H^1(E,\Zl)$; with
\Cref{prop:moddeg:mf-is-pullback}(c) this is the action on $M_{f,\lambda}$. Finally $\hat\delta_!(x,y)=\langle x,w_N^*y\rangle=\varepsilon_N\langle x,y\rangle$.
\end{proof}

This proof does not use the description of $M_{f,\lambda}\otimes\C$ as a plane; \Cref{prop:moddeg:f-lattice}(iii) gives a second proof.

\begin{theorem}[the modular degree is the ideal of values of $\hat\delta_!$]\label{thm:moddeg:degree-is-eta}
Let $E$ be optimal and $\ell$ any prime. For $a,b\in H^1(E,\Zl)$,
\[
\hat\delta_!(\varphi_E^*a,\varphi_E^*b)=\varepsilon_N(f)\,\degphi\,\langle a,b\rangle_E .
\]
Hence the ideal of values of $\hat\delta_!$ on $M_{f,\lambda}\times M_{f,\lambda}$ is $\degphi\cdot\Zl$, and the Gram determinant of $\hat\delta_!$ in a basis of $M_{f,\lambda}$ is $(\degphi)^2$ up to a unit of $\Zl$.
The same holds over $\Z$: the ideal of values of $\hat\delta_!$ on $H^1(X_0(N),\Z)[I_f]$ is $\degphi\cdot\Z$. Consequently, if the ideal $\eta^\emptyset_{f,\lambda}$ of Diamond, Flach and Guo is the
ideal of values of $\hat\delta_!$ (or of $\frac u2\hat\delta_!$, $u$ an $\ell$-adic unit) on $M_{f,\lambda}$, as Section 4.1 of the paper describes it for $\ell\nmid2N$, then
$v_\lambda(\eta^\emptyset_{f,\lambda})=v_\ell(\degphi)$, with no multiplicity-one or Gorenstein hypothesis.
\end{theorem}

\begin{proof}
By \Cref{lem:moddeg:wn-on-mf} and the projection formula (\Cref{prop:moddeg:degree-descriptions}(iv)),
$\hat\delta_!(\varphi_E^*a,\varphi_E^*b)=\varepsilon_N(f)\langle\varphi_E^*a,\varphi_E^*b\rangle_{X_0(N)}=\varepsilon_N(f)\degphi\langle a,b\rangle_E$. The cup product on $H^1(E,\Zl)$ is perfect and alternating of rank two, so its
ideal of values is $\Zl$ and its Gram determinant a unit (\Cref{ch:curves}); multiplying by the unit $\varepsilon_N(f)$ and by $\degphi$ gives ideal of values $\degphi\Zl$ and Gram determinant
$(\degphi)^2$ times a unit. By \Cref{prop:moddeg:mf-is-pullback}(c), $\varphi_E^*H^1(E,\Zl)$ is all of $M_{f,\lambda}$. The argument over $\Z$ is the same. The last sentence: $\frac u2$ is a unit for odd
$\ell$.
\end{proof}

In the language of physics: a loop on $E$ lifts to $\degphi$ loops on $X_0(N)$, and their intersection numbers add up.

Two checks. At level $11$, $\varphi_E$ is the identity, $\degphi=1$ and $\varepsilon_{11}(f)=-1$, and the identity reads $\hat\delta_!=-\langle\ ,\ \rangle$ on $H^1(X_0(11),\Zl)$, as computed after
\Cref{cor:moddeg:delta-hat}: the ideal of values is $\Zl$. At level $37$, $\varphi_E$ is the quotient by $w_{37}$, $\degphi=2$ and $\varepsilon_{37}(f_{37a})=+1$, so $\hat\delta_!$ and the cup product
agree on $M_{f,\lambda}$ and take the values $2\Zl$ there; for $\ell=2$ the cup product of the unimodular lattice $H^1(X_0(37),\Z_2)$ vanishes modulo $2$ on the saturated sublattice $M_{f,\lambda}$, and by
\Cref{cor:moddeg:ars-divisibility} this forces $2\mid r_E$, which is the congruence $f_{37a}\equiv f_{37b}\pmod2$ (\Cref{ex:moddeg:ex-level37-hecke}).

\noindent\emph{Commentary.} The paper states this as Proposition prop:pullback: ``Let $E$ be $X_0(N)$-optimal with newform $f$, $\varphi=\varphi_E$, $I_f$ the kernel of the eigencharacter
$\phi_f\colon\T\to\Z$ of $f$ on the full Hecke algebra, and $\ell$ a prime. Then $\varphi^*\colon H^1(E,\Z_\ell)\to H^1(X_0(N),\Z_\ell)$ is injective, with image the kernel lattice
$H^1(X_0(N),\Z_\ell)[I_f]$, and'' $\langle\varphi^*a,\,w_N^*\varphi^*b\rangle_{X_0(N)}=\varepsilon_N(f)\,\degphi\,\langle a,b\rangle_E$ for $a,b\in H^1(E,\Z_\ell)$ (displayed there), ``where $\langle\ ,\ \rangle$
denotes the cup product and $\varepsilon_N(f)=\pm1$ the eigenvalue of $w_N$ on $f$. In particular, for $\ell$ odd, the ideal of values of
$\frac12\langle x,w_N^*y\rangle$ on $H^1(X_0(N),\Z_\ell)[I_f]$ is $\degphi\,\Z_\ell$, and its Gram determinant is $(\degphi)^2$ up to a unit. If $\ell\nmid2N$, then $M_{f,\lambda}=\varphi^*H^1(E,\Z_\ell)$, and''
$v_\lambda(\eta^\emptyset_{f,\lambda})=v_\ell(\degphi)$ (displayed). The book's
\Cref{prop:moddeg:mf-is-pullback,lem:moddeg:wn-on-mf,thm:moddeg:degree-is-eta} prove this. The paper's ``In particular'' sentence is phrased for $\frac12\langle x,w_N^*y\rangle$ and so restricted to
odd $\ell$; the book also states the ideal of values over $\Z$ (on $H^1(X_0(N),\Z)[I_f]$ it is $\degphi\,\Z$) and at $\ell=2$ for $\langle x,w_N^*y\rangle$, both immediate from the displayed
identity. The two proofs differ in one step: the book derives the action of $w_N$ from the identity of maps $\varphi_E\circ w_N=\varepsilon_N(f)\varphi_E+\varphi_E(0)$, the paper from the fact that
the plane is also the eigenspace of the $T_n$, $(n,N)=1$, which commute with $w_N$ (\Cref{prop:moddeg:f-lattice}(i), (iii)). The Galois structure $M_{f,\lambda}\simeq T_\ell E(-1)$, which the paper
states in the paragraph after the proof, is part (d) of \Cref{prop:moddeg:mf-is-pullback}. The last sentence of the paper's proof: ``if $\ell\nmid2N$, the integral structure of \cite{DFG04} is
$H^1(X_0(N),\Z_\ell)$ (Proposition prop:integral), so $M_{f,\lambda}=H^1(X_0(N),\Z_\ell)[I_f]=\varphi^*H^1(E,\Z_\ell)$, and $\hat\delta_!$ is $\frac u2\langle x,w_N^*y\rangle$ with $u$ an
$\ell$-adic unit''. The identification of $H^1(X_0(N),\Zl)$ with the integral structure of \cite{DFG04} holds for every prime $\ell$ with no localisation (\Cref{rem:moddeg:dfg-normalisation}(1));
the hypothesis $\ell\nmid2N$ enters through the standing assumption of Section 4.1 of the paper and the normalisation of $\hat\delta_!$ (\Cref{rem:moddeg:dfg-normalisation}(3)), and no
hypothesis on $\bar\rho$ is needed, in the paper's statement or in \Cref{thm:moddeg:degree-is-eta}. Irreducibility of $\bar\rho$ enters only later, in the results of \cite{DFG04} that the paper applies
to $\eta^\emptyset_{f,\lambda}$, such as their Proposition 1.4(c).

\begin{corollary}[the divisibility of Agashe, Ribet and Stein at every prime]\label{cor:moddeg:ars-divisibility}
Let $E$ be optimal. For every prime $\ell$, $v_\ell(\degphi)\le v_\ell(r_E)$; that is, $\degphi\mid r_E$. Equality holds at $\ell$ if $H^1(X_0(N),\Zl)_{\mf m}$ is free over $\T_{\mf m}$ for $\mf m=\mf m_{f,\ell}$.
\end{corollary}

\begin{proof}
Let $\mathcal O=\Zl$, $\pi=\phi_f\otimes\Zl\colon\T_\ell\to\Zl$, $\wp=\ker\pi=I_f\T_\ell$ and $I=\Ann_{\T_\ell}(\wp)$. By \Cref{prop:moddeg:congruence-ideal} and Lemma lem:crit (\Cref{lem:moddeg:crit}),
$\eta=\pi(I)=r_E\Zl\ne0$. By \Cref{prop:moddeg:congruence-ideal}(ii) and \Cref{ch:modules}, $\wp=\wp_{\mf m}\times\T''$ with $\T''$ the product of the other local factors, so $\eta$ is the
congruence ideal of $\pi|_{\T_{\mf m}}$, and $L=H^1(X_0(N),\Zl)_{\mf m}$ satisfies $L[\wp_{\mf m}]=H^1(X_0(N),\Zl)[\wp]=M_{f,\lambda}$ (an element killed by $\wp$ is killed by $\T''$, hence lies in the
$\mf m$-component). $L$ is a faithful $\T_{\mf m}$-module, free of finite rank over $\Zl$, with the pairing $\hat\delta_!$, which is perfect on $L$ and for which $\T_{\mf m}$ is self-adjoint
(\Cref{cor:moddeg:delta-hat}(iii)). By \Cref{ch:modules} (the module ingredients of Lemma lem:ddt417, and the Gram determinant of a saturated sublattice):
$Q=L/(L[\wp_{\mf m}]\oplus L[I_{\mf m}])$ is killed by $\eta$, and $Q\simeq\coker\bigl(M_{f,\lambda}\to\Hom(M_{f,\lambda},\Zl)\bigr)$, the map induced by $\hat\delta_!$, so $\len Q$ is the valuation of the Gram
determinant, $2v_\ell(\degphi)$ by \Cref{thm:moddeg:degree-is-eta}. $Q$ is a quotient of $\Hom(M_{f,\lambda},\Zl)\simeq\Zl^2$ killed by $r_E$, hence a quotient of $(\Zl/r_E)^2$, and
$2v_\ell(\degphi)=\len Q\le2v_\ell(r_E)$. If $L\simeq\T_{\mf m}^2$, then $Q\simeq(\T_{\mf m}/(\wp_{\mf m}+I_{\mf m}))^2\simeq(\Zl/\eta)^2$ (\Cref{ch:modules}), and equality holds.
\end{proof}

The two module facts used from \Cref{ch:modules} are short enough to see directly; write $M=M_{f,\lambda}$ and drop the subscript $\mf m$. \emph{$\eta$ kills $Q$:} for $t\in I$ and
$x\in L$, $\pi(t)x=tx-(t-\pi(t))x$; here $tx\in L[\wp]$ because $\wp t=0$, and $(t-\pi(t))x\in\wp L\subseteq L[I]$ because $t-\pi(t)\in\wp$ and $I\wp=0$; so $\pi(t)x\in L[\wp]+L[I]$. \emph{$Q$ is the
cokernel of the Gram map:} the map $L\to\Hom(M,\Zl)$, $x\mapsto\hat\delta_!(\,\cdot\,,x)|_M$, is onto, because $\hat\delta_!$ is perfect on $L$ and $M$ is saturated (so every functional on
$M$ extends to $L$); its kernel is the orthogonal complement of $M$, which is $L[I]$ by the self-adjointness of $\T$ (\Cref{ch:modules}); so $L/(M\oplus L[I])\simeq\Hom(M,\Zl)/\{\hat\delta_!(\,\cdot\,,y)|_M:y\in M\}$.
In a basis of $M$ the last group is $\Zl^2$ modulo the column span of the Gram matrix $\sm0{\pm d}{\mp d}0$, $d=\degphi$ (\Cref{thm:moddeg:degree-is-eta}), that is $(\Zl/d)^2$, of length
$2v_\ell(d)$. So in fact $Q\simeq(\Zl/\degphi)^2$, and the inequality is the statement that $(\Zl/\degphi)^2$ is killed by $r_E$.

This is the assembly of the inequality half of Lemma lem:ddt417 in the case at hand ($d=2$); \Cref{ch:congruence} proves the lemma in general. The paper cites \cite{ARS12} for the
divisibility; the corollary derives it from \Cref{thm:moddeg:degree-is-eta}, so that the theorem of Agashe, Ribet and Stein is needed only for the reverse inequality.

\begin{theorem}[Lemma lem:hecke of the paper]\label{thm:moddeg:lem-hecke}
Let $E$ be optimal and $\ell$ a prime.
\begin{enumerate}[label=(\roman*)]
\item $v_\ell(\degphi)=v_\ell(\text{ideal of values of }\hat\delta_!\text{ on }M_{f,\lambda})\le v_\ell(r_E)$, with no further hypothesis.
\item If $\ell^2\nmid N$, then $v_\ell(\degphi)=v_\ell(r_E)$ (\Cref{thm:moddeg:ars}).
\item If $\ell\nmid2N$ and $\bar\rho_{E,\ell}$ is irreducible, then, with the integral structure of \cite{DFG04} identified with $H^1(X_0(N),\Zl)$ by the paper's Proposition prop:integral
(\Cref{rem:moddeg:dfg-normalisation}),
\[
v_\ell(\degphi)=v_\ell(r_E)=v_\lambda(\eta^\emptyset_{f,\lambda}).
\]
\end{enumerate}
\end{theorem}

\begin{proof}
(i) is \Cref{thm:moddeg:degree-is-eta} and \Cref{cor:moddeg:ars-divisibility}; (ii) is \Cref{thm:moddeg:ars}; (iii) combines (ii) with the last sentence of \Cref{thm:moddeg:degree-is-eta}.
\end{proof}

\noindent\emph{Commentary.} The paper states: ``Let $E$ be $X_0(N)$-optimal and $\ell$ an odd prime with $\ell\nmid 2N$ and $\bar\rho$ irreducible. Then $v_\ell(\degphi)=v_\ell(r_E)=v_\lambda(\eta^\emptyset_{f,\lambda})$.'' This is (iii); its proof uses only $\ell\nmid2N$, so the hypothesis that $\bar\rho$ is irreducible, kept in (iii) to match the paper's statement, is not used. The book adds (i) and (ii), which hold at every prime: (i) says that the
cohomological side of the identity needs no hypothesis on $\ell$ or $\bar\rho$, and it gives a proof of the inequality $v_\ell(\degphi)\le v_\ell(r_E)$ at every prime that does not go
through \cite{ARS12}; the divisibility itself is part of \cite[Theorem 2.2]{ARS12} as the paper cites it, so the new content is the route, not the statement. The paragraph after the lemma in the paper notes that Lemma lem:ddt417 alone gives only $v_\ell(r_E)\ge v_\lambda(\eta^\emptyset_{f,\lambda})$, with
equality under freeness. Consequences for the paper, assembled in \Cref{ch:congruence}: the identity (eq:identity) of
Theorem thm:main holds with equality whenever $\ell\notin S_f$, with no multiplicity-one hypothesis; and in Step 4 of the proof of Theorem thm:quant, (iii) applied to the optimal curve
of the class of the twist gives the comparison at level $N_S$ without freeness, so that the first inequality of Theorem thm:quant is an equality whenever $\ell\notin S_f$. Freeness is still
the tool for forms not attached to an optimal quotient, such as the depleted form of Step 2 (\Cref{rem:moddeg:rem-free-after-39}).

\begin{remark}[the identification with the objects of Diamond, Flach and Guo]\label{rem:moddeg:dfg-normalisation}
\Cref{thm:moddeg:degree-is-eta} is a theorem about the ideal of values of $\langle x,w_N^*y\rangle$ on $H^1(X_0(N),\Zl)[I_f]$. That this ideal is the $\eta^\emptyset_{f,\lambda}$ of \cite{DFG04} rests on five
points, which Section 4.1 of the paper establishes. Point (1) is the paper's Proposition prop:integral, which the paper proves and \Cref{ch:congruence} proves in full
(the proposition comparing the integral structure with $H^1(X_0(M),\Zl)$ there); point (2) is proved below and point (5) in \Cref{ch:modules}; \Cref{ch:congruence} checks the normalisation in (3), the kernel lattice in (4) and the
self-adjointness of the Hecke operators against the published text of \cite{DFG04} (Sections 1.5.2, 1.5.3, 1.6.1, 1.6.2, 1.8.2 there).
\begin{enumerate}[label=(\arabic*)]
\item \emph{Which curve.} The integral structure $M(N,1)_{!,\lambda}$ of \cite{DFG04} is defined through a modular curve of full level. The paper's Proposition prop:integral: ``Pull-back along the
covering of $X_0(M^\Sigma)$ by the modular curve of full level identifies $H^1(X_0(M^\Sigma),\Z)$ with the Betti lattice $M(M^\Sigma,1)_{!,B}$ of \cite[\S1.4.2]{DFG04}, compatibly with the Hecke
operators, $U_p$ included; in particular $M(M^\Sigma,1)_{!,\lambda}=H^1(X_0(M^\Sigma),\Z_\ell)$, with no localisation.'' It holds for every prime $\ell$, with no hypothesis on $\bar\rho$. The
proof has four steps: compact support makes the invariants of the full-level cover equal to the cohomology of $Y_0(M^\Sigma)$ (Hochschild-Serre, and the equivariant Leray spectral sequence of
the quotient map, whose $\mathcal R^1$ vanishes because a finite group has no nonzero homomorphism to $\Z$); the image of $H^1_c(Y_0)$ in $H^1(Y_0)$ is $H^1(X_0)$, which is point (2); Hecke
equivariance is checked over $\C$; and the $\lambda$-adic structure is the Betti one tensored with $\Zl$. The point matters because pulling back along a covering of degree $n$ multiplies the
cup product by $n$ (\Cref{ch:curves}): a lattice strictly between $\pi^*H^1(X_0(N),\Zl)$ and its saturation would change the ideal of values by a non-unit. Compact support is what excludes this.
With ordinary cohomology the invariants can be larger than the pull-back; in the paper's words (there $N$ is the full level and $M$ the level of $Y_0(M)$), ``$H^1(Y(N),\Z_\ell)^G/\pi^*H^1(Y_0(M),\Z_\ell)$ is the group of characters $\Gamma_0(M)\to\Q_\ell/\Z_\ell$
of $\ell$-power order that factor through $G$, namely $\gamma\mapsto\psi(d\bmod p)$ for primes $p\mid M$ with $\ell\mid p-1$; for $M=11$ and $\ell=5$ it is the unramified covering
$X_1(11)\to X_0(11)$''. For $X_1(N)\to X_0(N)$ this excess is measured by the Shimura subgroup (\Cref{def:moddeg:shimura-subgroup}, \Cref{ex:moddeg:ex-shimura}); it plays no part in the
comparison of integral structures.
\item \emph{Interior cohomology.} For constant coefficients, the image of $H^1_c(Y_0(N),\Zl)$ in $H^1(Y_0(N),\Zl)$ is $H^1(X_0(N),\Zl)$, and $H^1(\Gamma_0(N),\Zl)=H^1(Y_0(N),\Zl)$. Proof: with $S$ the set
of cusps, the exact sequence $H^0(S)\to H^1_c(Y)\to H^1(X)\to H^1(S)=0$ makes $H^1_c(Y)\to H^1(X)$ onto, and $H^1(X)\to H^1(Y)$ is injective because $H^1(X,Y)$ is, by excision, a sum of
$H^1(D,D\setminus\{P\})=0$ over small discs $D$ around the cusps; so the image of $H^1_c(Y)\to H^1(Y)$ is $H^1(X)$. And $H^1(\Gamma_0(N),\Zl)=\Hom(\Gamma_0(N),\Zl)$ kills the elements of finite order
($\Zl$ is torsion-free), and $\Gamma_0(N)$ modulo the normal subgroup they generate is $\pi_1(Y_0(N))$ (Armstrong's theorem on quotients by groups with fixed points), so
$\Hom(\Gamma_0(N),\Zl)=\Hom(\pi_1(Y_0(N)),\Zl)=H^1(Y_0(N),\Zl)$.
\item \emph{The operator.} $[UwU]_\omega$ is $w_N^*$ with no scalar in weight $2$ and trivial character (in \cite[Section 1.5.3]{DFG04} it ``coincides with $\psi(-1_\infty)N^{k-2}[Uw^{-1}U]_\omega$''),
and on classes pulled back from $X_0(N)$ the pairing $\hat\delta_!$ of \cite{DFG04} is (with $M^\Sigma=N$) ``$\frac u2$ times the cup product on $X_0(M^\Sigma)$ composed with $w_{M^\Sigma}$, for an $\ell$-adic unit
$u$''. The reason, as the paper gives it: at the sufficiently small level $U'=U_0(N)\cap U_1(4)$ the normalised pairing of \cite[Section 1.5.2]{DFG04} is Poincar\'e duality on $X_{U'}$ up to an
$\ell$-adic unit when $\ell\nmid2N$; it is normalised by dividing by the index, so on $U_0(N)$-invariants it is $[U_0(N):U']^{-1}$ times its value on $U'$-invariants; and the covering
$X_{U'}\to X_0(N)$ has degree $[U_0(N):U']/2$, because $-1$ acts trivially, so by the projection formula the cup product on $X_{U'}$ of two pulled-back classes is $[U_0(N):U']/2$ times their cup
product on $X_0(N)$. The book's $\hat\delta_!$ omits the factor $\frac u2$; this is harmless for odd $\ell$.
\item \emph{Kernel, not projection.} $M_{f,\lambda}$ is the kernel lattice $L[\wp]$, as the paper says and as \cite{DFG04} define it ($\mathcal M_f=\mathcal M(N,\psi)_![I_f]$ in their Section 1.6.2,
where $\mathcal M[I]$ is the kernel of $(x_1,\dots,x_r)\colon\mathcal M\to\mathcal M^r$ for generators $x_i$ of $I$, Section 1.1.2), and not the image of $L$ under the projection to the $f$-component (the dual lattice,
which would change the ideal of values by the congruence ideal).
\item \emph{Values versus Gram.} $\eta$ is the ideal of values, equivalently the square root of the Gram determinant (\Cref{ch:modules}), as the paper says; in \cite[Section 1.7.3]{DFG04}
$\eta^\Sigma_f$ is defined by the image of $\wedge^2\mathcal M_{f^\Sigma}$ under the pairing.
\end{enumerate}
Points (3) and (5) matter only up to units, and a unit changes nothing; points (1) and (4) are where a non-unit discrepancy could enter, and both hold exactly, so the paper's
description of $\eta$ in Section 4.1 agrees with \Cref{thm:moddeg:degree-is-eta}.
\end{remark}

\begin{example}[\curve{99a1} at $\ell=3$]\label{ex:moddeg:99-at-3}
Let $E=\curve{99a1}$ and $\ell=3$, so $3^2\mid N$ and the paper's theorems do not apply. By \Cref{thm:moddeg:degree-is-eta} the ideal of values of $\hat\delta_!$ on $M_{f,3}$ is $4\Z_3=\Z_3$, while
$v_3(r_E)=1$ (\Cref{ex:moddeg:99a1-and-726a1}). So the inequality of \Cref{cor:moddeg:ars-divisibility} is strict, and $L=H^1(X_0(99),\Z_3)_{\mf m}$ is not free over $\T_{\mf m}$ for $\mf m=\mf m_{f,3}$.
Since $\rank_{\Z_3}L=2\rank_{\Z_3}\T_{\mf m}$ (\Cref{thm:moddeg:gorenstein}, step (b) of the proof, which does not use $\ell\nmid N$), the freeness criterion of \Cref{ch:modules} shows
$\dim_{\F_3}L/\mf mL\ne2$ (here $\T/\mf m=\F_3$, since $\phi_f$ takes values in $\Z$); it is not $1$, for then $L$ would be cyclic by Nakayama's lemma, a quotient of $\T_{\mf m}$, of
$\Z_3$-rank at most $\rank\T_{\mf m}<2\rank\T_{\mf m}$; so it is at least $3$, and by \Cref{prop:moddeg:tate-module-vs-h1}(iii) and the duality $(W/\mf mW)^\vee\simeq W[\mf m]$ of step (a) of the proof of \Cref{thm:moddeg:gorenstein} (valid for every $\ell$)
$\dim_{\F_3}J_0(99)[\mf m]\ge3$: multiplicity one fails at $\mf m$. Yet $\bar\rho_{\curve{99a1},3}$ is irreducible: the class \curve{99a} has only a $2$-isogeny (\texttt{allisog}), and Cremona's
\texttt{galrep} file lists no non-surjective mod $3$ image for it. This is why the paper's Remark rem:free needs $\ell\nmid N$.
\end{example}

\begin{paperlink}\label{pl:moddeg:lem-hecke}
The paragraph after Lemma lem:hecke in the paper: ``Without the hypotheses the inequality of Lemma lem:ddt417 can be strict: at 99a1 and $\ell=3$, where $3\mid N$,
$\degphi=4$ and $r_E=12$, so $H^1(X_0(99),\Z_3)_{\mathfrak m}$ is not free over $\T_{\mathfrak m}$.'' This is \Cref{ex:moddeg:99-at-3}. The same paragraph: ``the congruence ideal is computed in
$\T_{\mathfrak m}$: the other local factors of $\T\otimes\Z_\ell$ lie in the kernel $\wp$'' is \Cref{prop:moddeg:congruence-ideal}(ii). Section 4.2 of the paper: ``by Proposition prop:pullback,
$\varphi_E^*$ identifies the lattice $M_{f,\lambda}$ of \cite{DFG04} with $H^1(E,\Z_\ell)\simeq T_\ell E(-1)$, so its adjoint is $\ad T$'' is \Cref{prop:moddeg:mf-is-pullback}(d).
\end{paperlink}

\section[Multiplicity one and the Gorenstein property]{Multiplicity one, the Gorenstein property and Remark rem:free}\label{sec:moddeg:mult-one}

The paper's Remark rem:free explains when $H^1(X_0(N),\Zl)_{\mf m}$ is free over the localised Hecke algebra: multiplicity one makes $\T_{\mf m}$ Gorenstein and the module free of rank two. By
\Cref{thm:moddeg:degree-is-eta} (the paper's Proposition prop:pullback), the paper's main theorems do not need this; it remains the conceptual reason for the equality in the theorem of Agashe, Ribet and Stein at $\ell\nmid2N$, and the
tool for forms that are not attached to an optimal quotient. This section defines the maximal ideal attached to $(f,\ell)$, states multiplicity one, proves the Gorenstein property from it
with \Cref{ch:modules}'s criterion, and says what Remark rem:free is still for.

\begin{definition}[the maximal ideal of $(f,\ell)$]\label{def:moddeg:maximal-ideal}
For $f=f_E$ and a prime $\ell$, $\mf m=\mf m_{f,\ell}$ is the maximal ideal of $\T_\ell$ containing $\ell$ and $I_f$, the kernel of $\T_\ell\to\Zl\to\Fl$, $t\mapsto\phi_f(t)\bmod\ell$; $\T_\ell=\prod_{\mf m'}\T_{\mf m'}$ over
its maximal ideals (\Cref{ch:modules}). For every maximal ideal $\mf m$ of $\T$ (or $\T_\ell$) with residue characteristic $\ell$ there is a semisimple continuous representation
$\bar\rho_{\mf m}\colon G_\Q\to\GL_2(\T/\mf m)$, unramified outside $N\ell$, with $\tr\bar\rho_{\mf m}(\Frob_p)=T_p\bmod\mf m$ and $\det\bar\rho_{\mf m}=\bar\chi_\ell$ (\Cref{ch:modgalrep}; stated). For
$\mf m=\mf m_{f,\ell}$, $\bar\rho_{\mf m}\simeq\bar\rho_{E,\ell}^{\mathrm{ss}}$, by \Cref{thm:moddeg:eichler-shimura}(iii), Chebotarev and Brauer-Nesbitt (\Cref{ch:galois,ch:reps}). $\mf m$ is
\emph{Eisenstein} if $\bar\rho_{\mf m}$ is reducible (\Cref{conv:moddeg:setup}(C7)); so $\mf m_{f,\ell}$ is Eisenstein iff $E[\ell]$ is reducible.
\end{definition}

\begin{theorem}[multiplicity one; stated]\label{thm:moddeg:multiplicity-one}
Let $\ell\nmid2N$ and let $\mf m$ be a non-Eisenstein maximal ideal of $\T$ of residue characteristic $\ell$. Then $J_0(N)[\mf m]$ is a $2$-dimensional $\T/\mf m$-vector space, isomorphic to $\bar\rho_{\mf m}$ as a
$G_\Q$-module.
\end{theorem}

We do not prove this. For prime $N$ it is Mazur's \cite[Chapter II]{Mazur77}; in general the argument is Mazur's, through the $q$-expansion principle modulo $\ell$ and the good
reduction of $J_0(N)$ at $\ell$, as written in \cite[Section 5]{Ribet90} (section number not verified). The isomorphism with $\bar\rho_{\mf m}$ follows from the dimension by the theorem of Boston,
Lenstra and Ribet. The paper's route, in Remark rem:free, is through $J_1(N)$: \cite[Theorem 3.5]{RibetStein} gives multiplicity one for $J_1(N)$, ``the only possible exception there having Serre
weight $\ell$, which does not occur in weight $2$ with $\ell\nmid N$'', and $J_0(N)[\mf m]\to J_1(N)[\mf m']$ is injective because its kernel lies in the Shimura subgroup, which is Eisenstein
(\Cref{def:moddeg:shimura-subgroup}).

\begin{theorem}[Gorenstein and free of rank two]\label{thm:moddeg:gorenstein}
In the situation of \Cref{thm:moddeg:multiplicity-one}, with $\mf m$ a maximal ideal of $\T_\ell$: $\T_{\mf m}$ is Gorenstein, and
$H^1(X_0(N),\Zl)_{\mf m}\simeq\T_{\mf m}^2\simeq T_\ell J_0(N)_{\mf m}$ as $\T_{\mf m}$-modules. Consequently, for $E$ optimal, $\ell\nmid2N$ and $E[\ell]$ irreducible, $v_\ell(r_E)=v_\ell(\degphi)$.
\end{theorem}

\begin{proof}
Let $L=H^1(X_0(N),\Zl)_{\mf m}$, with the pairing $\hat\delta_!$, perfect on $L$ with $\T_{\mf m}$ self-adjoint (\Cref{cor:moddeg:delta-hat}(iii)). We check the hypotheses of the freeness criterion of
\Cref{ch:modules} with $d=2$.
(a) $\dim_{\T/\mf m}L/\mf mL=2$. Put $W=L/\ell L=H^1(X_0(N),\Fl)_{\mf m}$; $\hat\delta_!$ induces a perfect pairing on $W$ for which $\T$ is self-adjoint, so $(W/\mf mW)^\vee\simeq W[\mf m]$ (a functional
vanishing on $\mf mW$ is $\langle-,y\rangle$ with $\langle\mf mW,y\rangle=\langle W,\mf my\rangle=0$, that is $y\in W[\mf m]$). Hence $\dim_{\Fl}L/\mf mL=\dim_{\Fl}W[\mf m]=\dim_{\Fl}J_0(N)[\mf m]$ by
\Cref{prop:moddeg:tate-module-vs-h1}(iii), and dividing by $[\T/\mf m:\Fl]$, $\dim_{\T/\mf m}L/\mf mL=2$ by \Cref{thm:moddeg:multiplicity-one}.
(b) $\rank_{\Zl}L=2\rank_{\Zl}\T_{\mf m}$. By \Cref{prop:moddeg:hecke-compatibility}(ii), (iii), $H^1(X_0(N),\C)=S_2\oplus\overline{S_2}$ with $S_2\simeq\Hom_\Z(\T,\C)$ as $\T$-modules, and
$\overline{S_2}\simeq S_2$ as well (complex conjugation on the second factor of $\Hom_\Z(\T,\Z)\otimes\C$ is $\T$-linear). So for $t\in\T$,
$\tr(t\mid H^1(X_0(N),\Q))=2\tr(t\mid\Hom_\Z(\T,\Q))=2\tr(t\mid\T\otimes\Q)$, the trace of a transpose being the same. Both sides are $\Ql$-linear in $t\in\T\otimes\Ql$ and agree on $\T$, so they agree on
the idempotent $e_{\mf m}$, whose trace on a module is the dimension of its $\mf m$-part: $\rank L=\dim e_{\mf m}H^1(X_0(N),\Ql)=2\dim e_{\mf m}(\T\otimes\Ql)=2\rank\T_{\mf m}$.
The criterion gives $L\simeq\T_{\mf m}^2$ and $\T_{\mf m}$ Gorenstein; $T_\ell J_0(N)_{\mf m}\simeq L$ by \Cref{prop:moddeg:tate-module-vs-h1}(ii). For the last statement, $\mf m_{f,\ell}$ is non-Eisenstein when
$E[\ell]$ is irreducible (\Cref{def:moddeg:maximal-ideal}), and \Cref{cor:moddeg:ars-divisibility} gives equality under freeness.
\end{proof}

This is the theorem the paper cites as \cite{Tilouine97} in Remark rem:free (see also \cite[Chapter 2]{Wiles95} and \cite[Section 4]{DDT}). Step (b) holds for every $\ell$ and every $\mf m$; only (a) uses
multiplicity one, and it is (a) that fails for \curve{99a1} at $\ell=3$ (\Cref{ex:moddeg:99-at-3}).

\begin{remark}[what Remark rem:free is still for]\label{rem:moddeg:rem-free-after-39}
After \Cref{thm:moddeg:degree-is-eta,thm:moddeg:lem-hecke}: (a) Remark rem:free is not needed for the equality in Lemma lem:hecke or in (eq:identity) of Theorem thm:main. (b) It is not needed
for the equality in the first inequality of Theorem thm:quant when $\ell\notin S_f$ either: \Cref{thm:moddeg:lem-hecke}(iii) applies to the optimal curve of each of the two classes compared. It
remains the natural tool for comparing the Hecke congruence ideal with the cohomological one for forms that are not attached to an optimal quotient, such as the depleted form $g^S$ of Step 2
of the paper's proof, where it turns an inequality into an equality. (c) It is the conceptual reason for the equality in the theorem of Agashe, Ribet and Stein at $\ell\nmid2N$ with $E[\ell]$
irreducible (\Cref{thm:moddeg:gorenstein}). It fails for $\ell\mid N$ in general (\Cref{ex:moddeg:99-at-3}).
\end{remark}

\begin{paperlink}\label{pl:moddeg:rem-free}
Remark rem:free, sentence by sentence. ``Freeness of $H^1(X_0(N),\Z_\ell)_{\mathfrak m}$ over $\T_{\mathfrak m}$ holds when $J_0(N)[\mathfrak m]$ is two-dimensional, since the localised Hecke algebra is then
Gorenstein and the Tate module is free of rank two'': \Cref{thm:moddeg:gorenstein}, step (a) being the two-dimensionality and step (b) a rank count. ``$J_0(N)[\mathfrak m]$ is two-dimensional for
$\ell\nmid2N$ and $\bar\rho$ irreducible'': \Cref{thm:moddeg:multiplicity-one} (stated). ``The identification of $T_\ell J_0(N)$ with $H^1(X_0(N),\Z_\ell)$ as Hecke modules uses Poincar\'e duality
composed with $w_N$'': \Cref{prop:moddeg:tate-module-vs-h1}. ``After Proposition prop:pullback this is not needed for Theorem thm:main or for Theorem thm:quant'': \Cref{rem:moddeg:rem-free-after-39}.
\end{paperlink}

\section{Cuspidal subgroup, Eisenstein ideal, Shimura subgroup}\label{sec:moddeg:eisenstein}

When $E[\ell]$ is reducible every theorem of the paper is out of reach: the maximal ideal $\mf m_{f,\ell}$ is Eisenstein, multiplicity one and the results of Diamond, Flach and Guo are
unavailable, and the modular degree can fall short of the local bound (the deficit $D$ of Conjecture conj:eis). Open question 4 asks for an interpretation of that deficit ``presumably
through the cuspidal subgroup of $J_0(N)$ \cite{Mazur77}''. This section defines the cuspidal subgroup, proves that the cusps $0$ and $\infty$ differ by a torsion point and computes its
image on the optimal curve, states Mazur's theorem at prime level, and defines the Shimura subgroup used in Remark rem:free.

\begin{definition}[the cuspidal subgroup]\label{def:moddeg:cuspidal-subgroup}
The \emph{cuspidal subgroup} $C_N\subseteq J_0(N)(\Qbar)$ is the subgroup generated by the classes of the divisors of degree $0$ supported on the cusps. It is $G_\Q$-stable, because $G_\Q$ permutes the
cusps (\Cref{prop:moddeg:cusps-over-q}); the \emph{rational cuspidal subgroup} is generated by the classes of the $G_\Q$-invariant cuspidal divisors of degree $0$, and lies in $J_0(N)(\Q)$.
\end{definition}

\begin{proposition}[Manin-Drinfeld for $0$ and $\infty$]\label{prop:moddeg:manin-drinfeld}
\leavevmode
\begin{enumerate}[label=(\roman*)]
\item For $p\nmid N$, $(T_p)_*(\infty)=(p+1)(\infty)$ and $(T_p)_*(0)=(p+1)(0)$ as divisors.
\item For $p\nmid N$ every eigenvalue of $T_p$ on $S_2(\Gamma_0(N))$ has absolute value at most $2\sqrt p<p+1$, so $(T_p-p-1)_*$ is an isogeny of $J_0(N)$.
\item The class $x=[(0)-(\infty)]$ lies in $\ker(T_p-p-1)_*$, so it is a torsion point of $J_0(N)(\Q)$.
\item (Manin; Drinfeld; stated.) $C_N$ is finite.
\end{enumerate}
\end{proposition}

\begin{proof}
(i) By \Cref{prop:moddeg:hecke-correspondence}(i), $(T_p)_*(\Gamma c)=\sum_i(\Gamma\alpha_ic)$ for a cusp $c$, with the representatives $\alpha_i$ of \Cref{prop:moddeg:hecke-correspondence}(ii), counted with the
multiplicities of $\pi_1$ (total $p+1$): every point of $X(\Gamma_\alpha)$ over $\Gamma c$ is $\Gamma_\alpha\gamma c$ for some $\gamma\in\Gamma$ and maps to $\Gamma\alpha\gamma c$, which is one of the $\Gamma\alpha_ic$. Now
$\sm1j0p\infty=\infty$ and $\sm p001\infty=\infty$, so all images of $\infty$ are $\infty$. And $\sm1j0p0=j/p$, $\sm p001 0=0$; for $1\le j<p$ the cusp $j/p$ has denominator $p$ with
$\gcd(p,N)=1$, so it is $\Gamma_0(N)$-equivalent to $0$ (\Cref{prop:moddeg:cusps-over-q}(i): the cusps of denominator prime to $N$ form the single class of $0$). So all images of $0$ are $0$.
(ii) By \Cref{thm:moddeg:eichler-shimura}(ii) the eigenvalues $\alpha$ of Frobenius on $W_\psi$ are roots of $X^2-\psi(T_p)X+p$, and $\abs\alpha=\sqrt p$ by Weil's theorem (\Cref{ch:curves}); so
$\psi(T_p)=\alpha+p/\alpha$ has absolute value at most $2\sqrt p$ for every eigencharacter $\psi$ (this is the Ramanujan-Petersson bound in weight $2$, \Cref{ch:modforms}). Hence $T_p-p-1$ is
invertible on $S_2$, its transpose on $S_2^\vee$ is invertible, and the induced endomorphism of $J_0(N)(\C)=S_2^\vee/H_1$ is surjective with finite kernel.
(iii) By (i), $(T_p-p-1)_*x=[(p+1)(0)-(p+1)(\infty)]-(p+1)x=0$; an element of the finite kernel of an isogeny is torsion; $0$ and $\infty$ are rational (\Cref{prop:moddeg:cusps-over-q}).
(iv) We do not prove this; see \cite[Chapter II]{Mazur77}, \cite[Section 6.6]{Stein} and \cite{Ogg73} (not verified). The proof is (i) to (iii) for every cusp, with the cusps not defined over
$\Q$ handled by the Galois action.
\end{proof}

\begin{theorem}[Mazur at prime level; stated]\label{thm:moddeg:mazur-prime-level}
Let $N$ be prime and $n$ the numerator of $(N-1)/12$.
\begin{enumerate}[label=(\roman*)]
\item The class of $(0)-(\infty)$ has order $n$ and generates $C_N$.
\item $J_0(N)(\Q)_{\tors}=C_N\simeq\Z/n$.
\item The \emph{Eisenstein ideal} $I_{\mathrm{Eis}}\subseteq\T$, generated by the $T_p-1-p$ ($p\ne N$) and $1+w_N$, satisfies $\T/I_{\mathrm{Eis}}\simeq\Z/n$; it annihilates $C_N$ and the Shimura subgroup $\mathcal S_N$,
which has order $n$, is of multiplicative type, and is Cartier dual to $C_N$.
\item For $\ell$ odd, a maximal ideal $\mf m$ of $\T$ of residue characteristic $\ell$ is Eisenstein in the sense of \Cref{conv:moddeg:setup}(C7) iff $\mf m\supseteq I_{\mathrm{Eis}}$, and such $\mf m$ exist iff $\ell\mid n$.
\end{enumerate}
\end{theorem}

We do not prove this; see \cite[Chapter II]{Mazur77} (the numbering of Mazur's results, and the cases $\ell=2$ and $\ell=N$ of (iv), were not checked against the text). For $N=11$: $n=5$, and
$J_0(11)(\Q)=\curve{11a1}(\Q)\simeq\Z/5$, generated by the cusp $0$ (\Cref{ex:moddeg:small-levels}); $\ell=5$ is Eisenstein, and indeed $a_p(f_{11})\equiv1+p\pmod5$ for $p\ne11$ ($a_2=-2$, $a_3=-1$, $a_5=1$,
$a_7=-2$), the congruence of $f_{11}$ with the Eisenstein series of \Cref{ch:modforms}. For $N=37$: $n=3$, and the curve \curve{37b1} has $E(\Q)\simeq\Z/3$.

\begin{definition}[the Shimura subgroup]\label{def:moddeg:shimura-subgroup}
Let $\pi\colon X_1(N)\to X_0(N)$ be the forgetful map. The \emph{Shimura subgroup} is $\mathcal S_N=\ker(\pi^*\colon J_0(N)\to J_1(N))$.
\end{definition}

\begin{proposition}[the Shimura subgroup]\label{prop:moddeg:shimura-subgroup}
\leavevmode
\begin{enumerate}[label=(\roman*)]
\item Assume $g(X_0(N))\ge1$. Then $\#\mathcal S_N=[\operatorname{sat}(\pi^*H^1(X_0(N),\Z)):\pi^*H^1(X_0(N),\Z)]$, the saturation taken in $H^1(X_1(N),\Z)$; so $\mathcal S_N=0$ iff
$\pi^*H^1(X_0(N),\Z)$ is saturated.
\item (Stated.) $\mathcal S_N$ is finite and of multiplicative type (Cartier dual to a constant group), and $T_p$ acts on it as $1+p$ for $p\nmid N$: it is Eisenstein.
\item If $\mf m$ is a non-Eisenstein maximal ideal of $\T$, then $\mathcal S_N[\mf m]=0$, so $\pi^*\colon J_0(N)[\mf m]\to J_1(N)$ is injective.
\end{enumerate}
\end{proposition}

\begin{proof}
(i) is \Cref{ch:curves} (the proposition on the degree and the cup product, part (iii), (a)$\Leftrightarrow$(d) with the index formula) for $\pi$. (ii) We do not prove this; see \cite[Chapter II]{Mazur77}
for prime $N$, and \cite{LingOesterle} and \cite{Ribet88} in general (not verified). (iii) If $0\ne x\in\mathcal S_N[\mf m]$, then $T_p-1-p$ and $\mf m$ both kill $x$, so $T_p\equiv1+p\pmod{\mf m}$ for all $p\nmid N\ell$
($\T/\mf m$ acts faithfully on the nonzero $\T/\mf m$-space spanned by the $\T$-translates of $x$). Then $\tr\bar\rho_{\mf m}(\Frob_p)=\tr(1\oplus\bar\chi_\ell)(\Frob_p)$ for those $p$, and by Chebotarev and
Brauer-Nesbitt the semisimple $\bar\rho_{\mf m}$ is $1\oplus\bar\chi_\ell$, reducible: a contradiction.
\end{proof}

For the optimal curve, $\mf m_{f,\ell}$ is Eisenstein exactly when $E$ has a rational $\ell$-isogeny (\Cref{def:moddeg:maximal-ideal}; \Cref{ch:galrep}). Then \Cref{thm:moddeg:multiplicity-one,thm:moddeg:gorenstein},
the theorems of \cite{DFG04} and the hypothesis of Lemma lem:hecke all fail, and the modular degree can lose local divisibility: the deficit $D$ of Conjecture conj:eis (\Cref{ch:galrep}).

\begin{proposition}[the image of the cusp $0$]\label{prop:moddeg:image-of-zero}
Let $E$ be optimal. Then $\varphi_E(0)\in E(\Q)_{\tors}$, and under $E(\C)\simeq\C/\Lambda_E$, $P\mapsto\int_O^P\omega_E$,
\[
\varphi_E(0)\equiv c_E\,L(E,1)\pmod{\Lambda_E}.
\]
Since $\Lambda_E\cap\R=\Z\Omega$ with $\Omega$ the least positive real period, $c_EL(E,1)/\Omega\in\Q$ and the order of $\varphi_E(0)$ is its denominator.
\end{proposition}

\begin{proof}
$\varphi_E=\varphi_{E*}\circ\iota$ and $\iota(0)=x=[(0)-(\infty)]$, so $\varphi_E(0)=\varphi_{E*}(x)$, a rational torsion point by \Cref{prop:moddeg:manin-drinfeld}(iii). For $z\in\HH$, the composite
$\HH\to X_0(N)\to E(\C)\to\C/\Lambda_E$ is $z\mapsto\int_{i\infty}^z\varphi_E^*\omega_E=c_E\int_{i\infty}^z2\pi if(\tau)\,d\tau$ (the integral converges at $i\infty$, $f$ being cuspidal), and it extends continuously to the
cusp $0$. Along the imaginary axis, $\tau=iy$:
\[
\int_{i\infty}^02\pi if(\tau)\,d\tau=\int_\infty^0(-2\pi)f(iy)\,dy=2\pi\int_0^\infty f(iy)\,dy=L(f,1),
\]
the last equality being the Mellin transform $(2\pi)^{-s}\Gamma(s)L(f,s)=\int_0^\infty f(iy)y^{s-1}\,dy$ at $s=1$ (\Cref{ch:lfunctions}; stated there). With $L(f,1)=L(E,1)$ (\Cref{thm:moddeg:modularity})
this gives the congruence. $\Lambda_E$ is stable under complex conjugation (E and $\omega_E$ are defined over $\R$), so $\Lambda_E\cap\R$ is a nonzero discrete subgroup $\Z\Omega$; a real number $r$ defines a
point of order $m$ in $\C/\Lambda_E$ iff $m$ is least with $mr\in\Z\Omega$.
\end{proof}

\begin{example}[$L(E,1)/\Omega$ and torsion]\label{ex:moddeg:image-of-zero}
PARI (\texttt{lfun(E,1)/E.omega[1]}, recognised as a rational number to $40$ digits) gives $L(E,1)/\Omega=1/5$, $1/6$, $1/3$, $1/6$, $1/3$, $2/3$, $1$, $1/3$ for \curve{11a1}, \curve{14a1}, \curve{19a1}, \curve{20a1},
\curve{27a1}, \curve{37b1}, \curve{121a1}, \curve{243b1}, whose torsion orders are $5$, $6$, $3$, $6$, $3$, $3$, $1$, $3$; and $L(E,1)=0$ for the rank-one curves \curve{37a1}, \curve{99a1}, \curve{539c1},
\curve{726a1}, so $\varphi_E(0)=O$ for them. With $c_E=1$ (Manin's conjecture, confirmed by computation in this range, \Cref{thm:moddeg:manin-known}(iv)) the order of $\varphi_E(0)$ is the denominator, and in each of these examples $\varphi_E(0)$ generates $E(\Q)_{\tors}$. For \curve{11a1}, $\varphi_E$ is the
identity of $X_0(11)$ and the statement is that the cusp $0$ has order $5$ (\Cref{thm:moddeg:mazur-prime-level}).
\end{example}

\begin{remark}[the Eisenstein case in the paper]\label{rem:moddeg:eisenstein-and-the-paper}
Ihara's lemma, the homothety of lattices behind part (c) of Proposition 1.4 of \cite{DFG04} (\Cref{ch:congruence}) and multiplicity one (\Cref{thm:moddeg:multiplicity-one}) all need $\bar\rho$ irreducible;
\Cref{prop:moddeg:mf-is-pullback} does not, and neither does the identification of the integral structure of \cite{DFG04} with $H^1(X_0(N),\Zl)$ (\Cref{rem:moddeg:dfg-normalisation}(1)), which holds
for every prime $\ell$. Conjecture conj:eis bounds the deficit by the $\ell$-part of the torsion, and
\Cref{ch:ellcurves} proves that its bound $D\le2$ at $\ell=3$ follows from the first (Proposition prop:no9 of the paper). \Cref{prop:moddeg:image-of-zero} links rational torsion with the
cuspidal subgroup: $\varphi_E(0)$ is a rational torsion point, the image of the cuspidal class $[(0)-(\infty)]$, and it often generates $E(\Q)_{\tors}$. This is a heuristic for Open question 4, not an
argument.
\end{remark}

\begin{paperlink}\label{pl:moddeg:eis}
Open question 4: ``Identify the deficit of Conjecture conj:eis with a natural invariant of the Eisenstein ideal at $\ell$, presumably through the cuspidal subgroup of $J_0(N)$ \cite{Mazur77}.''
\Cref{def:moddeg:cuspidal-subgroup,prop:moddeg:manin-drinfeld,thm:moddeg:mazur-prime-level,prop:moddeg:image-of-zero} are the objects named there. Remark rem:free: ``$J_0(N)[\mathfrak m]$ injects into
$J_1(N)[\mathfrak m']$ under the pull-back map, whose kernel, the Shimura subgroup, is Eisenstein'' is \Cref{prop:moddeg:shimura-subgroup}(iii).
\end{paperlink}

\section[Shimura curves, the prime 2, computing the degree]{Shimura curves, the prime \texorpdfstring{$2$}{2} and the computation of the modular degree}\label{sec:moddeg:shimura-two}

Three topics of the paper lie outside its theorems but next to them: the comparison of $\degphi$ with the degree of a parametrisation by a Shimura curve (Section 1.5, the history of the
multiplicative part of (eq:main)), the prime $\ell=2$ (Section sec:two: Watkins's conjecture and the theorem of Calegari and Emerton), and the computation of $\degphi$ on which all the
paper's data rest. This section states what is used, with data checks.

\begin{definition}[Shimura curves; stated]\label{def:moddeg:shimura-curve}
Let $D>1$ be a product of an even number of distinct primes and $M\ge1$ prime to $D$. The \emph{Shimura curve} $X_0^D(M)$ is the curve attached to an Eichler order of level $M$ in the indefinite
quaternion algebra over $\Q$ of discriminant $D$; it is a smooth projective curve over $\Q$ without cusps. By the Jacquet-Langlands correspondence, in Ribet's form \cite{Ribet90}, its Jacobian
$J_0^D(M)$ is isogenous over $\Q$, compatibly with the Hecke operators $T_p$, $p\nmid DM$, to the $D$-new quotient of $J_0(DM)$. Now let $E$ be semistable of conductor $N=DM$ (so $N$ is
square-free, the setting of \cite{RibetTakahashi97,Takahashi01}). Among the curves isogenous to $E$ there is exactly one, written $E_D$, which is an optimal quotient of $J_0^D(M)$, that is, the
target of a surjection $\xi_D\colon J_0^D(M)\to E_D$ with connected kernel, unique up to sign. The composite of $\xi_D$ with its dual $E_D\to J_0^D(M)$ is multiplication by a positive
integer $\delta_D$ (written $\delta^D(M)$ in the sources), the degree of the parametrisation $\varphi_{E,D}\colon X_0^D(M)\to J_0^D(M)\to E_D$ (a base point being chosen; the degree does not
depend on it). For $D=1$, $E_1$ is the $X_0(N)$-optimal curve and $\delta_1=\deg\varphi_{E_1}$. In general $E_D$ is not the $X_0(N)$-optimal curve: Takahashi notes ``(This $E$ generally differs
from $E$ in the optimal quotient $J_0(N)\to E$.)'' \cite[p.~75]{Takahashi01}.
\end{definition}

\begin{theorem}[Ribet and Takahashi; Takahashi; stated]\label{thm:moddeg:ribet-takahashi}
Let $E$ be the $X_0(N)$-optimal curve in an isogeny class of semistable elliptic curves over $\Q$ of conductor $N$. For a prime $r\mid N$ and a curve $E'$ isogenous to $E$ write
$n_r(E')=v_r(\Delta_{E'})=\#\Phig_r(E')$, and $n_r=n_r(E)$.
\begin{enumerate}[label=(\roman*)]
\item Let $N=dpqm$ with $p\ne q$ primes and $d$ a product of an even number of primes ($d=1$ allowed). Let $E_1=E_d$ and $\delta^d(pqm)$ be the optimal quotient of $J=J_0^d(pqm)$ and its degree,
and $E_2=E_{dpq}$ and $\delta^{dpq}(m)$ those of $J_0^{dpq}(m)$ (\Cref{def:moddeg:shimura-curve}). Let $i_q$ be the order of the image of the map $\Phig_q(J)\to\Phig_q(E_1)$ induced by the optimal
quotient. Then $i_q$ divides $n_q(E_1)$ and $n_p(E_2)$, and
\[
\frac{\delta^d(pqm)}{\delta^{dpq}(m)}=\frac{n_p(E_2)\,n_q(E_1)}{i_q^{2}}.
\]
\item If $E$ is the only curve in its isogeny class, then $\degphi/\delta_D=\prod_{r\mid D}n_r$ for every divisor $D$ of $N$ with an even number of prime factors.
\item If $\ell$ is a prime with $E[\ell]$ irreducible ($\ell=2$ and $\ell\mid N$ are allowed), then $\ell\nmid i_q$ in (i), and for every divisor $D$ of $N$ with an even number of prime factors
\[
v_\ell(\degphi)=v_\ell(\delta_D)+\sum_{r\mid D}v_\ell(n_r).
\]
\end{enumerate}
\end{theorem}

We do not prove (i) and (ii). Part (i) is \cite[Theorems 1 and 2]{RibetTakahashi97} (pages 11111 and 11112) in the sharpened form of \cite[Theorem 3.2]{Takahashi01} (page 82), whose part (c) reads
``$\delta/\delta'=(c'_p/i_q)\cdot j_q$. In particular, $\delta'$ divides $\delta$.'', together with $j_q=c_q/i_q$ for the order $j_q$ of the cokernel of the same map (proof of
\cite[Theorem 2.3]{Takahashi01}, page 80); divisibility of $n_p(E_2)$ by $i_q$ is \cite[Theorem 3.2(b)]{Takahashi01}. Part (ii) is \cite[Corollary 3.5]{Takahashi01} (page 83). In both sources
$c_r$ is our $n_r$, not the Tamagawa number: Ribet and Takahashi write ``Note that $c_p$, for instance, may be viewed as the order of the group of components of the fiber at $p$ of the
N\'eron model for $A$.'' \cite[page 11111]{RibetTakahashi97}, and Takahashi defines $c_r$ as the number of connected components of the fibre at $r$, that is $\ord_r(\Delta)$
\cite[pages 75 and 76]{Takahashi01}. The Tamagawa number $c_r=\#\Phig_r(\F_r)$ of this book differs from $n_r$ at nonsplit primes (\Cref{ch:reduction}): for \curve{57c2},
$n_{19}=5$ (PARI \texttt{elllocalred}) while its Tamagawa number at $19$ is $1$, and Takahashi's Table I lists $c'_{19}=5$ for this curve.

\begin{proof}[Sketch of proof of (iii) from (i)]
If $E'[\ell]$ is irreducible for one curve of the class, it is irreducible for all \cite[page 11113]{RibetTakahashi97}, and any two curves of the class are then linked by an isogeny of degree
prime to $\ell$, which induces isomorphisms of the $\ell$-primary parts of the component groups; so $v_\ell(n_r(E'))=v_\ell(n_r)$ for every $E'$ in the class and every $r\mid N$
\cite[page 11111]{RibetTakahashi97}. In (i) the prime $q$ divides the level $pqm$ of $J$, and \cite[Theorem 2.7]{Takahashi01} gives $\ord_\ell(j_q)=\ord_\ell(n_q(E_1))$; since $i_qj_q=n_q(E_1)$,
$\ell\nmid i_q$. Hence (i) gives $v_\ell(\delta^d(pqm))-v_\ell(\delta^{dpq}(m))=v_\ell(n_p)+v_\ell(n_q)$. Write $D=p_1q_1\cdots p_kq_k$ and apply this with $d=p_1q_1\cdots p_{j-1}q_{j-1}$,
$(p,q)=(p_j,q_j)$ and $m=N/(p_1q_1\cdots p_jq_j)$ for $j=1,\dots,k$; the sum telescopes from $\delta^1(N)=\degphi$ to $\delta^D(N/D)=\delta_D$.
\end{proof}

\begin{remark}[the sources, the hypotheses and the error term]\label{rem:moddeg:ribet-takahashi}
\emph{History and hypotheses.} For an isogeny class with one element, Bertolini and Darmon showed that the conjecture of Birch and Swinnerton-Dyer implies $\degphi/\delta_D=\prod_{r\mid D}n_r$ in $\Q^\times/(\Q^\times)^2$, and asked
whether it holds in $\Q^\times$ \cite[page 76]{Takahashi01}. Ribet and Takahashi \cite[Theorem 1]{RibetTakahashi97} prove the one-step formula in the shape
$\delta^{pqD}(M)=\delta^D(pqM)\,\mathcal E(D,p,q,M)^2/(c'_pc_q)$, with an ``error term'' $\mathcal E(D,p,q,M)$ that is a positive divisor of $c'_pc_q$; they show that every prime $\ell$ dividing it has
$A[\ell]$ reducible under the further hypothesis that $M$ is square-free but not a prime (the case of prime $M$ is excluded in their footnote), and their Theorem 2 identifies it as the product of
the order of the image of the map on component groups at $q$ and the order of the cokernel of the map at $p$. Takahashi proves that the map on component groups at a prime of the
discriminant is surjective \cite[Theorem 2.4]{Takahashi01}, a consequence of a result of Bertolini and Darmon, so that the error term is $i_q$, and removes the condition on $M$: ``In this paper,
we present a proof of $(*)$ which works for any $D$ and $M$.'' \cite[page 76]{Takahashi01}, $(*)$ being the formula of (ii). The standing hypothesis of both sources is that $N$ is
square-free; in (iii) no condition on $\ell$ other than the irreducibility of $E[\ell]$ enters. Pasten \cite[Section 1]{Pasten} summarises the formula in the form of \cite{RibetTakahashi97}, with
$M$ square-free and not a prime, and refines it in his Section 6 (a secondary pointer, not compared with the sources here).

\emph{The error term.} Put $\gamma_D=\degphi/(\delta_D\prod_{r\mid D}n_r)\in\Q^\times$. By (iii), $v_\ell(\gamma_D)=0$ whenever $E[\ell]$ is irreducible: $\gamma_D$ is supported on the primes
$\ell$ at which $E[\ell]$ is reducible. By (i) it is a product of the factors $i_q^{-2}$ and of ratios $n_r(E')/n_r$ for curves $E'$ isogenous to $E$. Ribet and Takahashi remark that the error
term is not necessarily divisible by all the primes $\ell$ for which $A[\ell]$ is reducible \cite[page 11111]{RibetTakahashi97}. Example (\cite[Section 4, Example 1 and Table I]{Takahashi01},
with $N=57$): for \curve{57c1}, $\degphi=12$, $n_3=10$, $n_{19}=1$, the curve $E_{57}$ is \curve{57c2}, $\delta_{57}=6$ and $i_3=5$; so $\gamma_{57}=12/(6\cdot10)=1/5$, and $E[5]$ is
reducible (\curve{57c1} has a rational $5$-isogeny to \curve{57c2}). For the isolated curve \curve{57a1}, $\degphi=4$, $n_3=2$, $n_{19}=1$, and (ii) gives $\delta_{57}=2$, as computed in
\cite[page 11111]{RibetTakahashi97}. (The values of $\degphi$ are those of \texttt{degphi} \cite{ecdata}, the $n_r$ those of PARI's \texttt{elllocalred}; $\delta_{57}$ and $i_3$ are taken from the sources.)

\emph{Use in the paper.} For $E[\ell]$ irreducible and $N$ square-free, (iii) gives
$v_\ell(\degphi)\ge\sum_{r\mid D}v_\ell(n_r)=\sum_{r\mid D}t_\ell(E,r)$, the last equality being $t_\ell(E,r)=v_\ell(n_r)$ at a multiplicative prime (Lemma lem:tamagawa-exp of the paper). This is
the multiplicative part of (eq:main) at the primes of $D$, an even number of them; the paper's Theorem thm:main contains it, under its own hypotheses on $\ell$, for every multiplicative prime. The
paper's own summary, in Section 1.5, is that they ``compared, for semistable $E$, $\degphi$ with the degree of a parametrisation by a Shimura curve and found the product of the $n_q$ over the primes of the quaternion
discriminant, exactly when $E$ is the only curve in its isogeny class and otherwise up to primes where $E[\ell]$ is reducible''; this matches (ii) and (iii).

\emph{The idea of the proof.} Proposition 1 of \cite{RibetTakahashi97} (Proposition 3.1 of \cite{Takahashi01}) is an exact sequence $0\to\mathcal X_p(J')\to\mathcal X_q(J)\to\mathcal X_q(J'')^2\to0$
of character groups of the tori of the reductions, with $J=J_0^d(pqm)$, $J'=J_0^{dpq}(m)$, $J''=J_0^d(qm)$, compatible with the Hecke operators and the monodromy pairings; for $d=1$ it is
\cite[Theorem 4.1]{Ribet90}, as both sources cite it. Grothendieck's exact sequence $0\to\mathcal X\to\Hom(\mathcal X,\Z)\to\Phi\to0$ expresses each component group as the cokernel of the
monodromy pairing, and comparing the $f$-isotypic lines on the two sides gives (i).
\end{remark}

\begin{conjecture}[Watkins]\label{conj:moddeg:watkins}
For every optimal $E/\Q$, $2^{\rank E(\Q)}\mid\degphi$.
\end{conjecture}

Data: all $38{,}042$ optimal curves of conductor below $10{,}000$ satisfy it (\texttt{degphi} and \texttt{allcurves} \cite{ecdata}); for instance \curve{37a1}, \curve{99a1}, \curve{121b1}, \curve{539c1}, \curve{726a1}
have rank $1$ and degrees $2$, $4$, $4$, $288$, $144$. The conjecture is due to Watkins \cite{Watkins02}; partial results are in \cite{Dummigan06,CalegariEmerton09,Yazdani11,CaroPasten22,HatleyKundu,BKP}, which the
paper cites without statements, and which this book does not state either.

\begin{theorem}[Calegari and Emerton; stated]\label{thm:moddeg:calegari-emerton}
Let $E/\Q$ be optimal of odd modular degree. Then $N$ has at most two odd prime factors, $E$ has even analytic rank, and either (a) $E$ has a rational point of order $2$, or (b) $N$ is prime, $E$
is supersingular at $2$ and $E(\R)$ is connected, or (c) $E$ has complex multiplication and $N\in\{27,32,49,243\}$.
\end{theorem}

We do not prove this. The statement is the one of Section sec:two of the paper: ``a curve of odd modular degree has at most two odd bad primes and even analytic rank, and either it has a
rational $2$-torsion point, or its conductor is prime, it is supersingular at $2$ and $E(\R)$ is connected, or it has complex multiplication and $N\in\{27,32,49,243\}$'', from
\cite{CalegariEmerton09} (as Theorem 1.1 of the arXiv version; the published numbering was not checked). Data (\texttt{degphi}, \texttt{allcurves}, and PARI for $a_2$, the discriminant and $j$): $103$
optimal curves of conductor below $10{,}000$ have odd modular degree, and all satisfy the theorem; $55$ have a rational point of order $2$, $46$ have prime conductor, no rational point of order $2$,
$a_2$ even and negative discriminant, and two, \curve{27a1} ($\degphi=1$) and \curve{243b1} ($\degphi=9$), both with $j=0$ and torsion $\Z/3$, are covered only by (c). The odd-degree optimal curves of
conductor below $100$ are \curve{11a1}, \curve{14a1}, \curve{15a1}, \curve{17a1}, \curve{19a1}, \curve{20a1}, \curve{21a1}, \curve{24a1}, \curve{27a1}, \curve{32a1}, \curve{33a1}, \curve{36a1}, \curve{46a1}, \curve{49a1},
\curve{52a1}, \curve{57b1}, \curve{67a1}, \curve{73a1}, \curve{89b1}.

\begin{remark}[computing the modular degree]\label{rem:moddeg:computing-the-degree}
Three methods. (a) Zagier's \cite{Zagier85}: compute $\varphi_{E*}$ on $H_1(X_0(N),\Z)$ from modular symbols and read $\degphi$ off the index of \Cref{prop:moddeg:degree-descriptions}(iv), the cup product
on $\varphi_E^*H^1(E,\Z)$ having ideal of values $\degphi\Z$ (\cite[Chapter II]{Cremona97}). (b) Watkins's \cite{Watkins02}: compute $\langle f,f\rangle$ from the special value of the symmetric square $L$-function
with the bad Euler factors and apply \Cref{thm:moddeg:zagier-watkins} with $c_E=1$; PARI's \texttt{ellmoddegree} implements this, and, as \Cref{ex:moddeg:zw-numerics} shows for \curve{11a3}, returns
$\deg\varphi/c_\varphi^2$ (here $1/5$) when applied to a non-optimal curve. (c) Cremona's comparison of the period lattices of the curves of a class with the lattice of $\varphi^*$ certifies optimality
(\Cref{conv:moddeg:cremona-optimal}). The paper's appendix: ``The function \texttt{ellmoddegree} needs about $32$ MB of stack with PARI 2.17 for conductors in the thousands''. \Cref{ch:computation}
has the sessions.
\end{remark}

\begin{paperlink}\label{pl:moddeg:two}
Section 1 of the paper: ``The modular degree is computable \cite{Cremona97,Watkins02}'' is \Cref{rem:moddeg:computing-the-degree}; Section sec:two quotes Watkins's conjecture
(\Cref{conj:moddeg:watkins}) and the theorem of Calegari and Emerton (\Cref{thm:moddeg:calegari-emerton}); Section 1.5 summarises Ribet and Takahashi (\Cref{thm:moddeg:ribet-takahashi,rem:moddeg:ribet-takahashi}).
\end{paperlink}

\section{Exercises}\label{sec:moddeg:exercises}

\begin{exercise}\label{ex:moddeg:ex-cusps}
List the cusps of $X_0(50)$ and of $X_0(72)$ by denominator, with their fields of definition, and find all $N\le30$ for which some cusp of $X_0(N)$ is not rational.
\emph{Answer.} $X_0(50)$: $12$ cusps; four each of denominators $5$ and $10$, conjugate over $\Q(\zeta_5)$, the other four rational. $X_0(72)$: $16$ cusps; two each of denominators $3$, $6$, $12$, $24$, conjugate
over $\Q(\sqrt{-3})$, the other eight rational. $N=9,16,18,25,27$ (\Cref{prop:moddeg:cusps-over-q}(iii)).
\end{exercise}

\begin{exercise}\label{ex:moddeg:ex-u2}
Let $f=f_{11}$ and $V=\operatorname{span}\{f(z),f(2z)\}=S_2(\Gamma_0(22))$. (a) Using $a_2(f)=-2$ and $f|_2W_{11}=-f$, show that in the basis $f(z)$, $f(2z)$ the matrices (columns are images) of $U_2$ and
$w_{22}$ are $U=\sm{-2}1{-2}0$ and $W=\sm0{-1/2}{-2}0$. (b) Show that $\T=\Z[U_2]$ at level $22$ (the $T_p$, $p\nmid22$, act as $a_p(f)$ and $U_{11}$ as $a_{11}(f)=1$), that $WUW=\sm0{-1/2}4{-2}$ is not of the form $a+bU$, and deduce that $w_{22}\T w_{22}\ne\T$
(\Cref{prop:moddeg:hecke-on-j0}(iii)). (c) With $A=\langle f,f\rangle_{\Gamma_0(22)}$, the Petersson Gram matrix in this basis is $G=\sm A{-A/3}{-A/3}{A/4}$ (check the off-diagonal entry with PARI's
\texttt{mfpetersson}). Verify $W^{\mathsf T}GW=G$ and $G^{-1}U^{\mathsf T}G=WUW$: the Petersson adjoint of $U_2$ is $w_{22}U_2w_{22}$, and $U_2$ is not self-adjoint.
\emph{Hint.} $U_2f=\sum a_{2n}q^n=a_2f-2f(2z)$ and $U_2(f(2z))=f$; $W_{22}=W_{11}\,\mathrm{diag}(2,1)$ and $\mathrm{diag}(2,1)W_{22}=2W_{11}$. For (b), $a+bU$ has entries $(a-2b,b;-2b,a)$.
\end{exercise}

\begin{exercise}\label{ex:moddeg:ex-counts}
Using \Cref{thm:moddeg:eichler-shimura}(iii) and $X_0(11)=\curve{11a1}$, predict $\#\curve{11a1}(\Fp)$ for $p=2,3,5,7,13$ from $f_{11}=q-2q^2-q^3+2q^4+q^5+2q^6-2q^7-2q^9-2q^{10}+q^{11}-2q^{12}+4q^{13}+\cdots$, and
check by counting points on $y^2+y=x^3-x^2-10x-20$. \emph{Answer.} $5,5,5,10,10$, from $p+1-a_p$ with $a_p=-2,-1,1,-2,4$: $3+2$, $4+1$, $6-1$, $8+2$, $14-4$.
All five are divisible by $5$, as they must be for $p\ne5,11$: the rational $5$-torsion of \curve{11a1} injects into $\tilde E(\Fp)$, torsion of order prime to $p$ injecting under
good reduction (\Cref{ch:reduction}).
\end{exercise}

\begin{exercise}\label{ex:moddeg:ex-37-lvalue}
Let $\varphi\colon X_0(37)\to X_0(37)/w_{37}=\curve{37a1}$ be the optimal parametrisation (\Cref{ex:moddeg:37}). (a) Show $\varphi\circ w_{37}=\varphi$ and deduce from \Cref{lem:moddeg:wn-on-mf} that
$\varepsilon_{37}(f_{37a})=+1$ and $\varphi(0)=O$. (b) Deduce from \Cref{prop:moddeg:image-of-zero} that $c_EL(\curve{37a1},1)\in\Z\Omega$. (c) The functional equation of $L(f,s)$ has sign $-\varepsilon_N(f)$
(\Cref{ch:hecke,ch:lfunctions}); conclude $L(\curve{37a1},1)=0$. \emph{Hint for (a).} $w_{37}(\infty)=0$, and $\varphi\circ w_{37}=\varepsilon_{37}\varphi+\varphi(0)$ with $\varphi\circ w_{37}=\varphi$ forces $\varepsilon_{37}=1$ and
$\varphi(0)=O$ (compare the pullbacks of $\omega_E$).
\end{exercise}

\begin{exercise}\label{ex:moddeg:ex-11a-degrees}
In the class \curve{11a} every isogeny $\curve{11a1}\to\curve{11a2}$ is $\psi_5\circ[m]$ for the cyclic $5$-isogeny $\psi_5$ and an integer $m\ne0$ (the curves have no complex multiplication), and likewise for
\curve{11a3}. (a) Show that the degrees of the modular parametrisations of \curve{11a2} and of \curve{11a3} are the numbers $5m^2$. (b) Show that a degree-$d$ isogeny $\psi\colon E\to E'$ satisfies
$u_\psi^2\vol(\Lambda_E)=d\vol(\Lambda_{E'})$, and use the volumes of \Cref{ex:moddeg:zw-numerics} to find $u_{\psi_5}=\pm1$ for \curve{11a2} and $\pm5$ for \curve{11a3}. (c) Deduce $c_\varphi=\pm m$, respectively $\pm5m$,
for the parametrisation of degree $5m^2$ (the sign is not normalised for a non-optimal curve), in agreement with \Cref{thm:moddeg:zagier-watkins}. \emph{Hint for (b).} Integrate $\frac i2\psi^*(\omega_{E'}\wedge\bar\omega_{E'})$ over $E$.
\end{exercise}

\begin{exercise}\label{ex:moddeg:ex-manin-class}
Let $E$ be optimal, $\ell$ a prime with $E[\ell]$ irreducible, and $\varphi'$ a modular parametrisation of minimal degree of a curve $E'$ of the class. Show that $v_\ell(c_{\varphi'})=v_\ell(c_E)$.
\emph{Hint.} $c_{\varphi'}=c_Eu_\psi$ with $\psi$ of degree prime to $\ell$ (\Cref{prop:moddeg:valuation-in-class}), and $u_\psi u_{\hat\psi}=\deg\psi$ with both factors integers.
\end{exercise}

\begin{exercise}\label{ex:moddeg:ex-77c}
The class \curve{77c} consists of the optimal curve \curve{77c1} $=[1,1,0,4,11]$ and \curve{77c2} $=[1,1,0,-51,110]$, linked by a cyclic $2$-isogeny $\psi\colon\curve{77c1}\to\curve{77c2}$.
(a) Compute $\degphi$ for \curve{77c1} from \Cref{thm:moddeg:zagier-watkins} with $c_E=1$ (\Cref{thm:moddeg:manin-known}, $77$ being square-free): take the newform of \curve{77c} in PARI (the element of
\texttt{mfeigenbasis(mfinit([77,2],0))} whose coefficients are those of \texttt{ellan}), compute \texttt{mfpetersson} of its symbol, multiply by $\mu(77)=96$, and divide $4\pi^2$ times the result by
$\vol(\Lambda_E)=\abs{\operatorname{Im}(\bar\omega_1\omega_2)}$ for \texttt{E.omega}. (b) Compute $\vol(\Lambda_{\curve{77c2}})$ and, with \Cref{ex:moddeg:ex-11a-degrees}(b), find $u_\psi$; deduce the degree of a minimal
parametrisation of \curve{77c2} and its $c_\varphi$. What does \texttt{ellmoddegree} return for \curve{77c2}?
\emph{Answer.} (a) $\langle f,f\rangle=0.3610248954\ldots$, $\vol=2.3754485978\ldots$, and $4\pi^2\langle f,f\rangle/\vol=6.0000\ldots$, so $\degphi=6$. (b) $\vol(\Lambda_{\curve{77c2}})=1.1877242989\ldots$, half of the
former, so $u_\psi^2=2\cdot\frac12=1$ and $u_\psi=\pm1$; the minimal parametrisation $\psi\circ\varphi_E$ has degree $2\cdot6=12$ (\Cref{thm:moddeg:optimal}(ii)) and $c_\varphi=c_Eu_\psi=\pm1$
(\Cref{prop:moddeg:manin-rational}); \texttt{ellmoddegree} returns $\deg\varphi/c_\varphi^2=12$.
\end{exercise}

\begin{exercise}\label{ex:moddeg:ex-88}
The optimal curve \curve{88a1} $=[0,0,0,-4,4]$ has $\degphi=8$, and the algorithm of \Cref{ex:moddeg:99a1-and-726a1} gives $r_E=16$. At which prime do the two differ, and why does this not contradict
\Cref{thm:moddeg:ars}? Show that $H^1(X_0(88),\Z_2)_{\mf m}$ is not free over $\T_{\mf m}$ for $\mf m=\mf m_{f,2}$. \emph{Answer.} At $2$, with $2^2\mid88$; freeness would give equality in
\Cref{cor:moddeg:ars-divisibility}.
\end{exercise}

\begin{exercise}\label{ex:moddeg:ex-md}
Let $x=[(0)-(\infty)]\in J_0(N)$. (a) For $N=11$, $T_2$ acts on $J_0(11)$ as $a_2=-2$; deduce $5x=0$ from $T_2x=3x$. (b) For $N=37$, the newforms have $(a_2,a_3)=(-2,-3)$ and $(0,1)$. Show $(T_2+2)T_2=0$ and
$(T_3+3)(T_3-1)=0$ on $J_0(37)$, and deduce $15x=21x=0$, so $3x=0$. Compare with \Cref{thm:moddeg:mazur-prime-level}. \emph{Hint.} $\T$ acts faithfully on $J_0(N)$ (\Cref{prop:moddeg:hecke-on-j0}).
\end{exercise}

\begin{exercise}\label{ex:moddeg:ex-up-cusp}
Let $N=p$ be prime. Show $(U_p)_*(\infty)=p(\infty)$ and $(U_p)_*(0)=(0)+(p-1)(\infty)$, so that $(U_p-1)_*x=0$ for $x=[(0)-(\infty)]$. Using $U_p=-w_p$ on $S_2(\Gamma_0(p))$ (all forms are new, and $a_p=-\varepsilon_p$ for a
newform of prime level, \Cref{ch:hecke}), relate this to the generator $1+w_p$ of Mazur's Eisenstein ideal. \emph{Hint.} The representatives are $\sm1j0p$, $0\le j<p$, and $j/p\sim\infty$ for $1\le j<p$ when $N=p$.
\end{exercise}

\begin{exercise}\label{ex:moddeg:ex-37-lattice}
Let $\varphi\colon X_0(37)\to E=\curve{37a1}$ be the quotient by $w_{37}$. Show that $\varphi^*H^1(E,\Z)$ is the $(+1)$-eigenlattice of $w_{37}^*$ on $H^1(X_0(37),\Z)$, that it equals $H^1(X_0(37),\Z)[I_f]$ for $f=f_{37a}$,
and that the cup product on it has ideal of values $2\Z$ and Gram determinant $4$. \emph{Hint.} $\varphi\circ w_{37}=\varphi$; the eigenlattice is saturated of rank $2$; use \Cref{ex:moddeg:37} and
\Cref{prop:moddeg:degree-descriptions}(iv).
\end{exercise}

\begin{exercise}\label{ex:moddeg:ex-j037}
With $E=\curve{37a1}$, $A=\varphi_E^*(E)$ and $B=\ker\varphi_{E*}$ in $J_0(37)$: show that $B$ is an elliptic curve on which $\T$ acts through the eigencharacter of $f_{37b}$, that $A\cap B\simeq(\Z/2)^2$, and that
addition $A\times B\to J_0(37)$ is an isogeny with kernel $\simeq(\Z/2)^2$. Deduce that $B$ lies in the class \curve{37b}. \emph{Hint.} \Cref{prop:moddeg:degree-descriptions}(iii); the Lie algebra of $B$ is the
annihilator of $\C f_{37a}$; \Cref{thm:moddeg:eichler-shimura}(ii) on $T_\ell B$.
\end{exercise}

\begin{exercise}\label{ex:moddeg:ex-level37-hecke}
At level $37$ show that $\T=\Z[T_2]\simeq\Z[x]/(x(x+2))$: check $a_n(f_{37a})\equiv a_n(f_{37b})\pmod2$ for $n\le6$ (Sturm's bound) to get $\T\subseteq\{(a,b)\in\Z^2:a\equiv b\bmod2\}$ under
$t\mapsto(\phi_{37a}(t),\phi_{37b}(t))$, and show that $\Z[T_2]$ is that whole ring. Deduce $\eta_{f_{37a}}=2\Z$ and $r_{\curve{37a1}}=2=\degphi$. \emph{Hint.} $T_2\mapsto(-2,0)$; $\Ann(\ker\phi_{37a})=\Z T_2$.
\end{exercise}

\begin{exercise}[harder]\label{ex:moddeg:ex-shimura}
Let $\pi$ be the forgetful map $X_1(N)\to X_0(N)$, of degree $\#(\Z/N)^\times/\{\pm1\}$ for $N\ge3$. (a) Show $\langle\pi^*x,\pi^*y\rangle_{X_1(N)}=\deg\pi\,\langle x,y\rangle_{X_0(N)}$, and that
$\pi^*H^1(X_0(N),\Z)$ is saturated iff $\mathcal S_N=0$. (b) For $N=11$, both curves have genus $1$, $\deg\pi=5$ and $\#\mathcal S_{11}=5$ (\Cref{thm:moddeg:mazur-prime-level}); show that the cup product on the
saturation of $\pi^*H^1(X_0(11),\Z)$ is unimodular, although on $\pi^*H^1(X_0(11),\Z)$ it has values $5\Z$. (c) Relate this to point (1) of \Cref{rem:moddeg:dfg-normalisation}: with ordinary cohomology the $G$-invariants of $H^1$ of a covering can be larger than the pull-back (here by the characters of
order $5$ of the paper's Remark after Proposition prop:integral), the ideal of values on the larger lattice is not the ideal of values on $X_0(N)$, and the difference is governed by the Shimura
subgroup, which is Eisenstein. Explain why the integral structure of \cite{DFG04}, built from compactly supported cohomology of a full-level cover, has no such excess.
\emph{Hint.} \Cref{ch:curves} (the proposition on the degree and the cup product, parts (i), (iii), (iv)) and \Cref{prop:moddeg:shimura-subgroup}(i).
\end{exercise}

\section{Notes and references}\label{sec:moddeg:notes}

\emph{Modular curves over $\Q$ and Hecke correspondences.} The $\Q$-structure of $X_0(N)$, the Hecke correspondences and the congruence relation are in Shimura's book \cite{Shimura71} (Chapters 6
and 7; we cite Sections 6.7, 7.2, 7.4 and 7.5 and Theorem 7.14 from memory, as unverified pointers), in Diamond and Shurman \cite{DiamondShurman} (Sections 3.8, 5.2, 6.1 to 6.6, 7.7, 7.9, 8.2 to 8.7, again
unverified pointers), and in the survey of Diamond and Im \cite{DiamondIm} (Sections 8 to 10). The integral models are Deligne and Rapoport's \cite{DeligneRapoport73} and Katz and Mazur's \cite{KatzMazur85}. The
cusps over $\Q$ are in Ogg \cite{Ogg73} and Stein's book \cite{Stein}. The rational points of $Y_0(N)$ are Mazur's \cite{Mazur78} and Kenku's \cite{Kenku79,Kenku80a,Kenku80b,Kenku81}, with Ligozat
\cite{Ligozat75} for $N=27$. The identity $T_n^t=w_NT_nw_N$ (\Cref{thm:moddeg:transpose}) is classical; the paper cites \cite[\S\S1.6.1, 1.8.2]{DFG04} and \cite[\S4.4]{DDT} for it; the proof here is
the matrix computation $W_N\,\mathrm{diag}(1,p)\,W_N^{-1}=\mathrm{diag}(p,1)$ combined with \Cref{ch:curves}'s theorem that the adjoint of a correspondence is its transpose. \Cref{ex:moddeg:ex-u2} shows that
$w_N\T w_N$ need not be $\T$; the book therefore treats $t\mapsto t^t$ as an isomorphism onto another ring of operators, not as an involution of $\T$.

\emph{Optimal curves and the modular degree.} The factorisation \Cref{thm:moddeg:optimal} and Cremona's conventions are in \cite[Chapter II]{Cremona97}; Stevens \cite{Stevens89} discusses optimal curves
for $X_1(N)$ as well. The proof of \Cref{thm:moddeg:optimal}(i) given here uses only the congruence relation, modularity and strong multiplicity one. Cremona's \texttt{ecdata} files \cite{ecdata} were
used for every curve-level number in the chapter (conductor below $10{,}000$), and every modular degree printed was recomputed with PARI's \texttt{ellmoddegree} \cite{PARI}.

\emph{The Manin constant and Zagier-Watkins.} Zagier \cite{Zagier85} computes $\degphi$ from modular symbols and compares it with congruence primes; Watkins \cite{Watkins02} computes it from the symmetric
square $L$-value; Rankin and Shimura \cite{Rankin39,Shimura76} give $\langle f,f\rangle$ (\Cref{ch:lfunctions}). The proof of \Cref{thm:moddeg:zagier-watkins} given here is direct, so the normalisations of
\cite{Zagier85,Watkins02} were not compared with the book's. The Manin constant: Manin's conjecture; Edixhoven \cite{Edixhoven91}; Mazur \cite{Mazur78}; Abbes and Ullmo \cite{AU96}; Agashe, Ribet and Stein
\cite{ARS06}; \v Cesnavi\v cius \cite{Cesnavicius18}; \v Cesnavi\v cius, Neururer and Saha \cite{CNS}. The statements in \Cref{thm:moddeg:manin-known}(iii), (iv) are recalled, not checked against these
texts.

\emph{Congruence numbers.} Ribet \cite{Ribet83} introduced congruence primes and proved $\degphi\mid r_E$ in an early form; Agashe, Ribet and Stein \cite{ARS12} prove \Cref{thm:moddeg:ars} (Theorem 2.2 and
Table 1 as the paper cites them; numbers not verified) and discuss the $abc$ connection (Section 2, with \cite{Frey,MaiMurty}; see also \cite{Murty99}). The congruence numbers of
\Cref{ex:moddeg:99a1-and-726a1} were computed for this book by the algorithm described there; they agree with the paper's values for \curve{99a1}, \curve{726a1} and \curve{539d1}.

\emph{The congruence ideal of Diamond, Flach and Guo.} \Cref{prop:moddeg:mf-is-pullback,lem:moddeg:wn-on-mf,thm:moddeg:degree-is-eta} are the book's version of the paper's Proposition prop:pullback; the
argument is the projection formula (compare \cite{Zagier85} and \cite{ARS12}, whose arguments we have not compared in detail). The objects of \cite{DFG04} (Sections 1.5.3, 1.6.1, 1.6.2, 1.7.3, equations (13), (14), as the paper cites them; checked against the published text in \Cref{ch:congruence}) are
compared with $H^1(X_0(N),\Zl)$ in \Cref{ch:congruence}, which proves the paper's Proposition prop:integral; \Cref{rem:moddeg:dfg-normalisation} lists what that comparison involves. Lemma lem:ddt417 is \cite[Lemma 4.17]{DDT}, inequality half.

\emph{Multiplicity one, Gorenstein, Eisenstein.} Mazur \cite{Mazur77} (Chapter II) for prime level and the Eisenstein ideal; Ribet \cite{Ribet90} (Section 5, unverified) for multiplicity one in general;
Ribet and Stein \cite{RibetStein} (Theorem 3.5 as the paper cites it) for $J_1(N)$; Tilouine \cite{Tilouine97}, Wiles \cite[Chapter 2]{Wiles95} and \cite[Section 4]{DDT} for the Gorenstein property; Ling and Oesterl\'e
\cite{LingOesterle} and Ribet \cite{Ribet88} for the Shimura subgroup. The proof of \Cref{thm:moddeg:gorenstein} from multiplicity one is \Cref{ch:modules}'s freeness criterion.

\emph{Modularity, Shimura curves, the prime $2$.} Modularity: \cite{Wiles95,TaylorWiles,BCDT}, with Carayol \cite{Carayol86} for the level. Shimura curves and the degree comparison: \cite{Ribet90,RibetTakahashi97,Takahashi01,Pasten}; \Cref{thm:moddeg:ribet-takahashi} was checked against \cite{RibetTakahashi97} (Theorems 1 and 2, Propositions 1 to 3) and \cite{Takahashi01}
(Theorems 2.3, 2.4, 2.7, 3.2, Corollaries 3.4, 3.5, Table I); Pasten's summary (Section 1 of \cite{Pasten}) is a secondary pointer that we have not compared with them. Watkins's conjecture \cite{Watkins02} and the results of \cite{Dummigan06,CalegariEmerton09,Yazdani11,CaroPasten22,HatleyKundu,BKP}; \Cref{thm:moddeg:calegari-emerton} is quoted from the paper and the arXiv version of \cite{CalegariEmerton09}. The data checks of
\Cref{conj:moddeg:watkins,thm:moddeg:calegari-emerton} use all $38{,}042$ optimal curves of conductor below $10{,}000$.
